\chapter{Celestial Mechanics \& Gravitation}

\section*{Theory problems}

\probhead{2.1}{}{2007}{Chiang Mai, Thailand}{TH 1.6}{2}
% source id: 2007_TH_1.6   tag: comet period from Kepler third law
A Sun-orbiting periodic comet is the farthest at 31.5\,A.U. and the closest at
0.5\,A.U. What is the orbital period of this comet?

\probhead{2.2}{}{2007}{Chiang Mai, Thailand}{TH 1.7}{2}
% source id: 2007_TH_1.7   tag: areal velocity, Kepler second law
For the comet in question 1.6, what is the area (in square A.U. per year) swept
by the line joining the comet and the Sun?

\probhead{2.3}{}{2008}{Bandung, Indonesia}{TH S14}{20}
% source id: 2008_TH_S14   tag: orbit mass/energy from u-p data
Consider a Potentially Hazardous Object (PHO) moving in a closed orbit under
the influence of the Earth's gravitational force. Let $u$ be the inverse of the
distance of the object from the Earth and $p$ be the magnitude of its linear
momentum. As the object travels, the graph of $u$ as a function of $p$ passes
through points A and B as shown in the following table. Find the mass and the
total energy of the object, and express $u$ as a function of $p$ and sketch the
shape of the $u$ curve from A to B.

\begin{center}
\begin{tabular}{lcc}
\toprule
 & $p$ ($\times 10^{9}$\,$\mathrm{kg\,m\,s^{-1}}$)
 & $u$ ($\times 10^{-8}$\,$\mathrm{m^{-1}}$) \\
\midrule
A & 0.052 & 5.15 \\
B & 1.94 & 194.17 \\
\bottomrule
\end{tabular}
\end{center}

\probhead{2.4}{}{2008}{Bandung, Indonesia}{TH S8}{20}
% source id: 2008_TH_S8   tag: horizontal tidal acceleration due to Moon
Gravitational forces of the Sun and the Moon lead to the raising and lowering
of sea water surfaces. Let $\varphi$ be the difference in longitude between
points A and B, where both points are at the equator and A is on the sea
surface. Derive the horizontal acceleration of sea water at position A due to
the Moon's gravitational force at the time when the Moon is above point B
according to observers on the Earth (express it in $\varphi$, the radius $R$ of
Earth, and the Earth--Moon distance $r$).

\probhead{2.5}{}{2009}{Tehran, Iran}{TH P12}{10}
% source id: 2009_TH_P12   tag: Earth orbit growth from solar mass loss
Calculate how much the radius of the Earth's orbit increases as a result of the
Sun losing mass due to the thermo-nuclear reactions in its center in 100 years.
Assume the Earth's orbit remains circular during this period.

\probhead{2.6}{}{2009}{Tehran, Iran}{TH P14}{10}
% source id: 2009_TH_P14   tag: satellite distance from zenith angle
A satellite is orbiting around the Earth in a circular orbit in the plane of
the equator. An observer in Tehran at the latitude of $\varphi = 35.6^\circ$
observes that the satellite has a zenith angle of $z = 46.0^\circ$, when it
transits the local meridian. Calculate the distance of the satellite from the
center of the Earth (in the Earth radius unit).

\probhead{2.7}{High Altitude Projectile}{2009}{Tehran, Iran}{TH P16}{45}
% source id: 2009_TH_P16   tag: projectile as elliptical orbit
A projectile which starts from the surface of the Earth at the sea level is
launched with the initial speed of $v_0 = \sqrt{GM/R}$ and with the projecting
angle (with respect to the local horizon) of $\theta = \frac{\pi}{6}$. $M$ and
$R$ are the mass and radius of the Earth respectively. Ignore the air
resistance and rotation of the Earth.

\begin{description}
\item[(a)] Show that the orbit of the projectile is an ellipse with a
  semi-major axis of $a = R$.
\item[(b)] Calculate the highest altitude of the projectile with respect to the
  Earth surface (in the unit of the Earth radius).
\item[(c)] What is the range of the projectile (distance between launching
  point and falling point) in the units of the earth radii?
\item[(d)] What is eccentricity ($e$) of this elliptical orbit?
\item[(e)] Find the time of flight for the projectile.
\end{description}

\probhead{2.8}{}{2009}{Tehran, Iran}{TH P3}{10}
% source id: 2009_TH_P3   tag: planet radius escapable by jumping
Estimate the radius of a planet that a man can escape its gravitation by
jumping vertically. Assume density of the planet and the Earth are the same.

\probhead{2.9}{}{2009}{Tehran, Iran}{TH P6}{10}
% source id: 2009_TH_P6   tag: escape velocity from uniform cloud
Derive a relation for the escape velocity of an object, launched from the
center of a proto-star cloud. The cloud has uniform density with the mass of
$M$ and radius $R$. Ignore collisions between the particles of the cloud and
the launched object. If the object were allowed to fall freely from the
surface, it would reach the center with a velocity equal to $\sqrt{GM/R}$.

\probhead{2.10}{}{2010}{Beijing, China}{TH L16}{30}
% source id: 2010_TH_L16   tag: parabolic orbit crossing Mars
A spacecraft is launched from the Earth and it is quickly accelerated to its
maximum velocity in the direction of the heliocentric orbit of the Earth, such
that its orbit is a parabola with the Sun at its focus point, and grazes the
Earth orbit. Take the orbit of the Earth and Mars as circles on the same plane,
with radius of $r_E = 1$\,AU and $r_M = 1.5$\,AU, respectively. Make the
following approximation: during most of the flight only the gravity from the
Sun needs to be considered. The figure below shows the trajectory of the
spacecraft (not in scale); the inner circle is the orbit of the Earth, the
outer circle is the orbit of Mars.

\ioaafig{2010_TH_L16-f1.pdf}

\begin{description}
\item[(a)] What is the angle $\psi$ between the path of the spacecraft and the
  orbit of the Mars (see the figure) as it crosses the orbit of the Mars,
  without considering the gravity effect of the Mars?
\item[(b)] Suppose the Mars happens to be very close to the crossing point at
  the time of the crossing, from the point of view of an observer on Mars, what
  is the approaching velocity and direction of approach (with respect to the
  Sun) of the spacecraft before it is significantly affected by the gravity of
  the Mars?
\end{description}

\probhead{2.11}{}{2010}{Beijing, China}{TH S6}{10}
% source id: 2010_TH_S6   tag: walking speed vs orbital velocity on asteroid
A spacecraft landed on the surface of a spherical asteroid with negligible
rotation, whose diameter is 2.2\,km, and its average density is
2.2\,$\mathrm{g/cm^3}$. Can the astronaut complete a circle along the equator
of the asteroid on foot within 2.2 hours? Write your answer ``YES'' or ``NO''
on the answer sheet and explain why with formulae and numbers.

\probhead{2.12}{}{2011}{Katowice, Poland}{TH S1}{10}
% source id: 2011_TH_S1   tag: Oort-cloud comet infall travel time
Most single-appearance comets enter the inner Solar System directly from the
Oort Cloud. Estimate how long it takes a comet to make this journey. Assume
that in the Oort Cloud, 35\,000\,AU from the Sun, the comet was at aphelion.

\probhead{2.13}{}{2011}{Katowice, Poland}{TH S3}{10}
% source id: 2011_TH_S3   tag: Voyager orbit type; Sun's magnitude
On 9 March 2011 the Voyager probe was 116.406\,AU from the Sun and moving at
17.062\,$\mathrm{km\,s^{-1}}$. Determine the type of orbit the probe is on:
(a) elliptical, (b) parabolic, or (c) hyperbolic. What is the apparent
magnitude of the Sun as seen from Voyager?

\probhead{2.14}{}{2011}{Katowice, Poland}{TH S4}{10}
% source id: 2011_TH_S4   tag: Phobos time above Martian horizon
Assuming that Phobos moves around Mars on a perfectly circular orbit in the
equatorial plane of the planet, give the length of time Phobos is above the
horizon for a point on the Martian equator. Use the following data: radius of
Mars $R_{\mathrm{Mars}} = 3393$\,km, rotational period of Mars
$T_{\mathrm{Mars}} = 24.623$\,h, mass of Mars
$M_{\mathrm{Mars}} = 6.421 \times 10^{23}$\,kg, orbital radius of Phobos
$R_P = 9380$\,km.

\probhead{2.15}{}{2011}{Katowice, Poland}{TH S6}{10}
% source id: 2011_TH_S6   tag: tidal spin-down, ancient day count
Tidal forces result in a torque on the Earth. Assuming that, during the last
several hundred million years, both this torque and the length of the sidereal
year were constant and had values of $6.0 \times 10^{16}$\,N\,m and
$3.15 \times 10^{7}$\,s respectively, calculate how many days there were in a
year $6.0 \times 10^{8}$ years ago. Moment of inertia of a homogeneous filled
sphere of radius $R$ and mass $m$ is
\[ I = \frac{2}{5} m R^2. \]

\probhead{2.16}{}{2011}{Katowice, Poland}{TH S7}{10}
% source id: 2011_TH_S7   tag: impulse direction vs orbit types
A satellite orbits the Earth on a circular orbit. The initial momentum of the
satellite is given by the vector $\mathbf{p}$. At a certain time, an explosive
charge is set off which gives the satellite an additional impulse
$\Delta\mathbf{p}$, equal in magnitude to $|\mathbf{p}|$. Let $\alpha$ be the
angle between the vectors $\mathbf{p}$ and $\Delta\mathbf{p}$, and $\beta$
between the radius vector of the satellite and the vector $\Delta\mathbf{p}$.
By thinking about the direction of the additional impulse $\Delta\mathbf{p}$,
consider if it is possible to change the orbit to each of the cases given
below. If it is possible mark YES on the answer sheet and give values of
$\alpha$ and $\beta$ for which it is possible. If the orbit is not possible,
mark NO.

\begin{description}
\item[(a)] a hyperbola with perigee at the location of the explosion.
\item[(b)] a parabola with perigee at the location of the explosion.
\item[(c)] an ellipse with perigee at the location of the explosion.
\item[(d)] a circle.
\item[(e)] an ellipse with apogee at the location of the explosion.
\end{description}

Note that for $\alpha = 180^\circ$ and $\beta = 90^\circ$ the new orbit will be
a line along which the satellite will free fall vertically towards the centre
of the Earth.

\probhead{2.17}{}{2012}{Rio de Janeiro, Brazil}{TH S5}{}
% source id: 2012_TH_S5   tag: Earth/Sun density ratio from g, year
Calculate the ratio between the average densities of the Earth and the Sun,
using \textbf{ONLY} the dataset below:
\begin{itemize}
\item the angular diameter of the Sun, as seen from Earth
\item the gravitational acceleration on Earth's surface
\item the length of the year
\item the fact that one degree in latitude at Earth's surface corresponds to
  111\,km
\end{itemize}

\probhead{2.18}{}{2013}{Volos, Greece}{TH L2}{50}
% source id: 2013_TH_L2   tag: Seneca asteroid orbit, Jupiter mass
A spacecraft is orbiting the Near Earth Asteroid (2608) \textit{Seneca}
(staying continuously very close to the asteroid), transmitting pulsed data to
the Earth. Due to the relative motion of the two bodies (the asteroid and the
Earth) around the Sun, the time it takes for a pulse to arrive at the ground
station varies approximately between 2 and 39 minutes. The orbits of the Earth
and Seneca are coplanar. Assuming that the Earth moves around the Sun on a
circular orbit (with radius $a_{\mathrm{Earth}} = 1$\,AU and period
$T_{\mathrm{Earth}} = 1$\,yr) and that the orbit of Seneca does not intersect
the orbit of the Earth, calculate:

\begin{description}
\item[(1)] the semi-major axis, $a_{\mathrm{Sen}}$, and the eccentricity,
  $e_{\mathrm{Sen}}$, of Seneca's orbit around the Sun;
\item[(2)] the period of Seneca's orbit, $T_{\mathrm{Sen}}$, and the average
  period between two consecutive oppositions, $T_{\mathrm{syn}}$, of the
  Earth--Seneca couple;
\item[(3)] an approximate value for the mass of the planet Jupiter,
  $M_{\mathrm{Jup}}$ (assuming this is the only planet of our Solar system with
  non-negligible mass compared to the Sun). Assume that the presence of
  Jupiter does not influence the orbit of Seneca.
\end{description}

\probhead{2.19}{Cosmic Pendulum}{2014}{Suceava, Romania}{TH L2}{}
% source id: 2014_TH_L2   tag: gravity-gradient pendulum on shuttle
A space shuttle (N) orbits the Earth in the equatorial plane on a circular
trajectory with radius $r$. From the spaceship. The shuttle has an arm
designated to place satellites on the orbit. The arm is a metallic rod
(negligible mass) with length $l \ll r$. The arm is connected to the shuttle
with frictionless mobile articulation. A satellite \textbf{S} is attached to
the arm and let out from the shuttle. At a certain moment the angle between the
rod and the shuttle's orbit radius is $\alpha$, see the figure. You know the
mass of the Earth -- $M$, and the gravitational constant $G$.

\ioaafig{2014_TH_L2-f1.pdf}

\begin{description}
\item[(a)] Find out the values of angle $\alpha$ for which the configuration
  of the system shuttle--rod--satellite remains unchanged regardless to Earth
  (the system is in equilibrium), during orbiting the Earth. For each found
  value of angle $\alpha$, specify the type of the system equilibrium i.e.
  stable or unstable.

  You will take in to account the following assumptions: the initial orbit of
  the shuttle is not affected by the presence of the satellite S, all the
  external friction-type interactions are negligible, the satellite--shuttle
  gravitational interaction is negligible too. The following data are known:
  $m_1$ -- the mass of the shuttle, the mass of the satellite $m_2 \ll m_1$.
\item[(b)] In the moment of one stable equilibrium configuration, the rod with
  the satellite attached is slightly rotated with a very small angle in the
  orbital plane and then released. Demonstrate that the small oscillations of
  the satellite S relative to the shuttle are harmonically ones. Express the
  period $T_0$ of this cosmic pendulum as a function of the orbiting period $T$
  of the shuttle around the Earth.

  It is known the linear harmonic oscillator equation:
  \[ \frac{\mathrm{d}^2 \beta}{\mathrm{d}t^2} + \omega_0^2 \beta = 0; \quad
     \omega_0 = \frac{2\pi}{T_0}, \]
  where $\beta$ -- the instantaneous angular deviation, $T_0$ -- the period of
  the linear harmonic oscillator.
\item[(c)] If we consider that the mass of the satellite S, $m_2$, is not
  negligible by comparison with the shuttle's one $m_1$ in the conditions from
  point (a) the evolution on the orbit of the shuttle would be influenced by
  the presence of the satellite S rigid attached to the shuttle by the rod.
  Identify and determine the consequences on the shuttle's movement after one
  complete rotation around the Earth.
\item[(d)] Propose a special technical maneuver which can cancels the
  influence of the non negligible mass satellite S on the shuttles movement.
\end{description}

\probhead{2.20}{From Romania to Antipod! A ballistic messenger}{2014}{Suceava, Romania}{TH L3}{}
% source id: 2014_TH_L3   tag: ballistic trajectory to antipode
The 8th IOAA organizers plan to send to the \textbf{antipode} the official
flag using a ballistic projectile. The projectile will be launched from
Romania, and the rotation of the Earth will be neglected.

\begin{description}
\item[(a)] Calculate the coordinates of the target-point if the launch-point
  coordinates are: $\varphi_{\mathrm{Romania}} = 44^\circ$ North;
  $\lambda_{\mathrm{Romania}} = 30^\circ$ East.
\item[(b)] Determine the elements of the launching-speed vector, regardless to
  the center of the Earth, in order that the projectile should hit the target.
\item[(c)] Calculate the velocity of the projectile when it hits the target.
\item[(d)] Calculate the minimum velocity of the projectile.
\item[(e)] Calculate the flying-time of the projectile, from the launch-moment
  to the impact one. You will use the value of the gravitational acceleration
  at Earth surface $g_0 = 9.81$\,$\mathrm{m\,s^{-2}}$; the Earth radius
  $R = 6370$\,km.
\item[(f)] Evaluate the possibility that the projectile to be seen with the
  naked eye in the moment that it passes at the maximum distance from the
  Earth. You will use the following values: the Moon albedo
  $\alpha_M = 0.12$; the Moon radius $R_M = 1738$\,km; the Earth--Moon
  distance $r_{EL} = 384\,400$\,km; the apparent magnitude of the full moon
  $m_M = -12.7^{\mathrm{m}}$. You assume that the projectile is perfectly
  metallic sphere with radius
  $r_{\mathrm{projectile}} = 400 \times 10^{-3}$\,m and with perfectly
  reflective surface.
\end{description}

\probhead{2.21}{Lagrange Points}{2014}{Suceava, Romania}{TH P1}{}
% source id: 2014_TH_P1   tag: locate L3 Lagrange point Sun-Earth
The \textit{Lagrange} points are the five positions in an orbital
configuration, where a small object is stationary relative to two big bodies,
only gravitationally interacting with them. For example, an artificial
satellite relative to Earth and Moon, or relative to Earth and Sun. In the
figure below are sketched two possible orbits of Earth relative to Sun and of
a small satellite relative to the Sun. Find out which of the two points
$L_3^1$ and $L_3^2$ could be the real Lagrange point relative to the system
Earth--Sun, and calculate its position relative to Sun. You know the following
data: the Earth--Sun distance $d_{ES} = 15 \times 10^{7}$\,km and the
Earth--Sun mass ratio $M_E/M_S = 1/332\,946$.

\ioaafig{2014_TH_P1-f1.pdf}

\probhead{2.22}{Sun gravitational catastrophe!}{2014}{Suceava, Romania}{TH P2}{}
% source id: 2014_TH_P2   tag: new orbit after Sun's mass halves
In a gravitational catastrophe, the mass of the Sun mass decrease instantly to
half of its actual value. If you consider that the actual Earth orbit is
elliptical, its orbital period is $T_0 = 1$\,year and the eccentricity of the
Earth orbit is $e_0 = 0.0167$.

Find the period of the Earth's orbital motion, after the gravitational
catastrophe, if it occurs on: a) 3rd of July, b) 3rd of January.

\probhead{2.23}{The Astronaut saved by \ldots\ ice from a can!}{2014}{Suceava, Romania}{TH P5}{}
% source id: 2014_TH_P5   tag: reaction propulsion via vapor thrust
An astronaut, with mass $M = 100$\,kg, get out of the space ship for a
repairing mission. He has to repair a satellite standing still relatively to
shuttle, at about $d = 90$\,m distance away from the shuttle. After he finished
his job he realizes that the systems designated to assure his come-back to
shuttle were broken. He also observed that he has air only for 3 minutes. He
also noticed that he possessed a hermetically closed cylindrical can (base
section $S = 30\,\mathrm{cm}^2$) firmly attached to its glove, with
$m = 200$\,g of ice inside. The ice did not completely fill the can.

Determine if the astronaut is able to arrive safely to the shuttle, before his
air reserve is empty. Briefly explain your calculations. Note that he cannot
throw away anything of its equipment, or touch the satellite.

\emph{You may use the following data:} $T = 272$\,K --- the temperature of the
ice in the can; $p_s = 550$\,Pa --- the pressure of the saturated water vapors
at the temperature $T = 272$\,K; $R = 8300$\,J/(kmol$\cdot$K) --- the constant
of perfect gas; $\mu = 18$\,kg/kmol --- the molar mass of the water.

\probhead{2.24}{}{2015}{Magelang, Indonesia}{TH L1}{}
% source id: 2015_TH_L1   tag: moon orbit geometry from surface observer
A moon is orbiting a planet such that the plane of its orbit is perpendicular
to the surface of the planet where an observer is standing. After some
necessary scaling, suppose the orbit satisfies the following equation:
\[ 9\left(\frac{x}{2} + \frac{\sqrt{3}y}{2} - 4\right)^2
   + 25\left(-\frac{\sqrt{3}x}{2} + \frac{y}{2}\right)^2 = 225 \]
Consider Cartesian coordinates where $x$ is on the horizontal plane and $y$ is
on the zenith of the observer. Let $r$ be the radius of the moon. Assume that
the period of rotation of the planet is much larger than the orbital period of
the moon. Ignore the atmospheric refraction.

\begin{description}
\item[(a)] Calculate the semimajor and semiminor axis of the ellipse.
\item[(b)] Calculate the zenith angle of perigee.
\item[(c)] Determine $\tan\frac{\theta}{2}$ where $\theta$ is the elevation
  angle (altitude of the upper tangent of the moon) when the moon looks
  largest to the observer.
\end{description}
\ioaafig{2015_TH_L1-f1.pdf}

\probhead{2.25}{}{2015}{Magelang, Indonesia}{TH S1}{}
% source id: 2015_TH_S1   tag: orbital resonances Gliese 876 planets
Several exoplanets have been observed in the Gliese 876 system
($M_{\mathrm{G}} = 0.33 \pm 0.03\,M_\odot$) as given in the following table,

\begin{center}
\begin{tabular}{|l|l|l|}
\hline
\textbf{Gliese System} & \textbf{Mass} & \textbf{Semi Major Axis (AU)} \\
\hline
Gliese 876 b & $2.276\,M_{\mathrm{J}}$ & 0.2083 \\
\hline
Gliese 876 c & $0.714\,M_{\mathrm{J}}$ & 0.1296 \\
\hline
Gliese 876 d & $6.8\,M_\oplus$ & 0.0208 \\
\hline
Gliese 876 e & $15\,M_\oplus$ & 0.334 \\
\hline
\end{tabular}
\end{center}

where $M_\odot$ is mass of Sun, $M_{\mathrm{J}}$ is mass of Jupiter
($M_{\mathrm{J}} = 1.89813 \times 10^{27}$\,kg), and $M_\oplus$ is mass of
Earth. Assume that all these planets revolve around Gliese 876 in the same
direction. Two planets are said to be in resonant orbits if the synodic period
of one planet with respect to the other planet is an integer multiple of the
orbital period of the second planet.

Find if any of the exoplanets of Gliese 876 system may have resonant orbits.

\probhead{2.26}{}{2015}{Magelang, Indonesia}{TH S11}{}
% source id: 2015_TH_S11   tag: mass from density profile, escape velocity
The mass density of an object is inversely proportional to the radial distance
from the center of the object with a factor of proportionality
$\alpha = 5.0 \times 10^{13}\,\mathrm{kg/m}^2$. If the escape velocity at the
surface of the object is $v_0 = 1.5 \times 10^4$\,m/s, calculate the total
mass of the object.

\probhead{2.27}{}{2015}{Magelang, Indonesia}{TH S13}{}
% source id: 2015_TH_S13   tag: Io tidal heating from surface distortion
Volcanic activity on Io, whose rotation period synchronizes with its orbital
period, was proposed to be the result of tidal heating mainly from Jupiter.
The resultant tidal force on a body is the difference in gravitational force
experienced by the near and far sides of that body due to another body.
Measurements of the surface distortion of Io via satellite radar altimeter
mapping indicate that the surface rises and falls by up to 100\,m during
one-half orbit. Only the surface layers will move by this amount. Interior
layers within Io will move by a smaller amount, and thus we assume that on
average the entire mass of Io is moved through 50\,m. Assume that Io is
considered as two hemispheres each treated as a point mass. Calculate the
\emph{average power} of the tidal heating on Io.

\textbf{Hint:} you can use the following approximation
$(1+x)^n \approx 1 + nx$ for small $x$.

The mass of Io is $m_{\mathrm{Io}} = 8.931938 \times 10^{22}$\,kg\\
The perijove distance is $r_{peri} = 420\,000$\,km\\
The apojove distance is $r_{apo} = 423\,400$\,km\\
The orbital period of Io is 152853\,s\\
The radius of Io is $R_{\mathrm{Io}} = 1821.6$\,km.

\probhead{2.28}{}{2015}{Magelang, Indonesia}{TH S2}{}
% source id: 2015_TH_S2   tag: planet density from satellite orbit
One satellite of a planet has an orbital period of 7 days, 3 hours,
43 minutes, and the semi major axis is 15.3 times the mean radius of the
planet. The Moon has an orbital period of 27 days, 7 hours, 43 minutes and the
semi major axis is 60.3 times the Earth's mean radius. Assume that the mass of
the moon and the satellite is negligible compared to the mass of the planet.
Calculate the ratio of the planet's mean density to that of the Earth.

\probhead{2.29}{}{2015}{Magelang, Indonesia}{TH S5}{}
% source id: 2015_TH_S5   tag: eccentricity from apparent angular rates
An observer is trying to determine an approximate value of the orbital
eccentricity of a man-made satellite. When the satellite was at apogee, it was
observed to have moved by $\Delta\theta_1 = 2'44''$ in a short time. When the
radius vector connecting Earth and the satellite is perpendicular to the major
axis (true anomaly is equal to $90^\circ$), within the same duration of time,
it was observed to have moved by $\Delta\theta_2 = 21'17''$. Assume that the
observer is located at the center of the Earth. Find an approximate value of
the eccentricity of the satellite's orbit.

\probhead{2.30}{True or False}{2016}{Bhubaneswar, India}{TH T1}{10}
% source id: 2016_TH_T1   tag: true/false grab bag: orbits, analemma, expansion
Determine if each of the following statements is True or False. In the Summary
Answersheet, tick the correct answer (TRUE / FALSE) for each statement. No
justifications are necessary for this question.

\begin{description}
\item[(T1.1)] In a photograph of the clear sky on a Full Moon night with a
  sufficiently long exposure, the colour of the sky would appear blue as in
  daytime. \hfill[2]
\item[(T1.2)] An astronomer at Bhubaneswar marks the position of the Sun on
  the sky at 05:00 UT every day of the year. If the Earth's axis were
  perpendicular to its orbital plane, these positions would trace an arc of a
  great circle. \hfill[2]
\item[(T1.3)] If the orbital period of a certain minor body around the Sun in
  the ecliptic plane is less than the orbital period of Uranus, then its orbit
  must necessarily be fully inside the orbit of Uranus. \hfill[2]
\item[(T1.4)] The centre of mass of the solar system is inside the Sun at all
  times. \hfill[2]
\item[(T1.5)] A photon is moving in free space. As the Universe expands, its
  momentum decreases. \hfill[2]
\end{description}

\probhead{2.31}{Mars Orbiter Mission}{2016}{Bhubaneswar, India}{TH T9}{20}
% source id: 2016_TH_T9   tag: perigee burn, new apogee, eccentricity
India's Mars Orbiter Mission (MOM) was launched using the Polar Satellite
Launch Vehicle (PSLV) on 5 November 2013. The dry mass of MOM (body +
instruments) was 500\,kg and it carried fuel of mass 852\,kg. It was initially
placed in an elliptical orbit around the Earth with perigee at a height of
264.1\,km and apogee at a height of 23903.6\,km, above the surface of the
Earth. After raising the orbit six times, MOM was transferred to a trans-Mars
injection orbit (Hohmann orbit).

The first such orbit-raising was performed by firing the engines for a very
short time near the perigee. The engines were fired to change the orbit
without changing the plane of the orbit and without changing its perigee. This
gave a net impulse of $1.73 \times 10^5$\,kg\,m\,s$^{-1}$ to the satellite.
Ignore the change in mass due to burning of fuel.

\begin{description}
\item[(T9.1)] What is the height of the new apogee, $h_{\mathrm{a}}$ above the
  surface of the Earth, after this engine burn? \hfill[14]
\item[(T9.2)] Find the eccentricity ($e$) of the new orbit after the burn and
  the new orbital period ($P$) of MOM in hours. \hfill[6]
\end{description}

\probhead{2.32}{Distance of the Lagrangian point L2 of the Earth-Moon system}{2019}{Keszthely, Hungary}{TH 9}{20}
% source id: 2019_TH_9   tag: Earth-Moon L2 distance
On 3 January 2019, the Chinese spacecraft Chang'e-4 landed on the far side of
the Moon in the area of the von K\'arm\'an crater, which was named after the
world famous Hungarian-born physicist Theodore von K\'arm\'an.

As the Earth remains below the horizon of the spacecraft all the time, a relay
station is also necessary for the communication with mission control on the
Earth. For this purpose, The Chinese Space Agency launched a spacecraft,
Queqiao, which was placed into a halo orbit around the outer Lagrangian point
of the Earth--Moon system, $L_2$, the far side of the Moon.

Calculate the distance ($h$) of this satellite above the surface of the Moon.
The Moon's orbit should be considered as a perfect circle with a radius of
$R = 384\,400$\,km. Neglect perturbations from the Sun and other planets.

\emph{Hint:} You can use the following approximation:
$1/(1+x)^2 \approx 1 - 2x$, if $|x| \ll 1$.

\probhead{2.33}{Occultation of a X-ray Source}{2020}{GeCAA (Estonia, remote)}{TH 6}{20} % sic: "a X-ray" in source
% source id: 2020_TH_6   tag: satellite-orbit occultation timing
Consider a satellite observing x-ray sources, while orbiting the Earth in the
equatorial plane with orbital radius $r$, and orbital time period $P$. Let us
assume that this satellite is pointed to one fixed direction in space for a
given length of time. Take the radius of the earth as $R$.

When the satellite moves `behind' the earth, naturally, the x-ray source is
`occulted' and the measured x-ray flux from the source drops to zero. However,
due to Earth's atmosphere, this drop is gradual. If the line of sight of the
source passes through the atmosphere, the attenuation depends on the air-mass
(i.e.\ length of air column) along the line of sight.

\begin{description}
\item[(a)] Let us assume that pointing towards a fixed source at $0^\circ$
  declination. We consider that the source is occulted when 50\% of the light
  coming from the source gets attenuated due to the atmosphere. Let us say
  that this happens when the minimum height of the line of sight from the
  surface of the Earth is $h$.

  If $\theta_0$ is the angle between the direction to the source and the
  direction to the Earth, as measured from the spacecraft, find an expression
  for $\theta_0$. \hfill[1 point]
\item[(b)] The time duration $\Delta t$ between the source getting attenuated
  from 90\% of pre-occultation flux to 10\% is defined as the `occultation
  time' for the source. Assume the flux attenuates to 90\% when the minimum
  height of the line of sight is $(h + 0.5\Delta h)$ and similarly the flux
  attenuates to 10\% at $(h - 0.5\Delta h)$, where $\Delta h \ll R$.\\
  Find the expression for $\Delta t$ in terms of $r$, $P$, $\Delta h$ and
  $\theta_0$. \hfill[4 points]
\item[(c)] If the satellite was pointing towards a source at declination
  $\beta$ instead ($\beta$ not too large), what will be the expression for
  $\Delta t$? \hfill[15 points]
\end{description}

\textbf{Note:} If the satellite was not in the equatorial plane, then the
problem could have been simply rephrased by assuming the satellite's orbital
plane to be the equatorial plane. In that case, $\beta$ becomes `relative
declination'.

\probhead{2.34}{Minimum velocity of a projectile}{2021}{Bogot\'a (remote)}{TH Q11}{15}
% source id: 2021_TH_Q11   tag: minimum-energy ballistic ellipse
\begin{description}
\item[(11.1)] What is the minimum speed with which a projectile must be
  launched from the Earth's surface at the equator such that the projectile
  reaches the north pole? \hfill[12.0\,pt]
\item[(11.2)] Find the eccentricity of the trajectory described by the
  projectile \hfill[3.0\,pt]
\end{description}

You may ignore the rotation of the Earth. Also assume the earth surface is
spherical.
\ioaafig{2021_TH_Q11-f1.pdf}

\probhead{2.35}{Hodograph}{2021}{Bogot\'a (remote)}{TH Q12}{15}
% source id: 2021_TH_Q12   tag: velocity-space hodographs of orbits
In curvilinear motion of a planet around a star, the direction of the velocity
vector changes continuously. This can be represented by a so-called
``trajectory in velocity space'' and is obtained as follows: for each point on
the spatial trajectory, the corresponding velocity vector is drawn so that its
starting point is at the origin of the velocity space, and its magnitude and
direction is the same as the velocity vector at that point. The tip of this
variable velocity vector generates a curve in velocity space. (The name
`hodograph' was given to this curve by Hamilton in 1846.)

As an example, see figures 1 and 2 below. For a circular orbit (Figure 1), the
magnitude of the velocity is constant and therefore, the hodograph (Figure 2)
of the velocity vector for Keplerian circular motion is also a circle, the
center of which is located at the origin of the velocity space. The radius of
this circle is equal to the constant magnitude of the circular velocity.

\ioaafig{2021_TH_Q12-f1.pdf}
\ioaafig{2021_TH_Q12-f2.pdf}

\begin{description}
\item[(12.1)] Write an expression for the radius of the hodograph in Fig.~2,
  as a function of the mass $M$ of the star, and the radius $R$ of the
  circular orbit of the planet's motion. \hfill[1.0\,pt]
\item[(12.2)] For a planet in a Keplerian trajectory, write the expression for
  centripetal acceleration vector ($\vec{a}$) and the magnitude of angular
  momentum ($L$). For any Keplerian trajectory, it is true that
  \[ |\Delta v| = k\Delta\theta \tag{1} \]
  Where $k$ is a constant for each type of Keplerian trajectory. Find the
  expression for the constant $k$ as a function of the masses $M$ and $m$ of
  the star and the planet, respectively, and the angular momentum, $L$.\\
  (Eq.~1) allows us to conclude that for any Keplerian trajectory, the
  hodograph ($v$ as a function of $\theta$) is a circle, but except for
  circular motion, the centre of the hodograph does not coincide with the
  star. It is not necessary to prove this result, you may simply accept it as
  a given. For the hodograph of uniform circular motion, the compliance with
  (eq.~1) is completely obvious, as evidenced in Fig.~3. \hfill[4.0\,pt]

  \ioaafig{2021_TH_Q12-f3.pdf}

\item[(12.3)] Determine the expression of the constant $k$ for the hodograph
  of circular planetary motion. \hfill[2.0\,pt]
\item[(12.4)] Given that the hodograph of the Keplerian elliptical motion is a
  circle, determine the radius of this hodograph and the distance between the
  center of the hodograph and the position of the star, as a function of the
  velocities at periastron and apoastron. Draw a rough sketch of the hodograph
  in the answer sheet as per the schematic shown in Fig.~4. The black circle
  is the star. \hfill[4.0\,pt]

  \ioaafig{2021_TH_Q12-f4.pdf}

\item[(12.5)] Similarly, for the parabolic Keplerian trajectory, determine the
  radius of the corresponding hodograph and the distance from the center of
  that hodograph circle to the star. Express the radius as a function of the
  velocity at periastron. Draw a rough sketch of the hodograph circle in the
  answer sheet. \hfill[4.0\,pt]
\end{description}

\probhead{2.36}{Formation of the Venus-2}{2021}{Bogot\'a (remote)}{TH Q14}{35}
% source id: 2021_TH_Q14   tag: comet-Venus collision, new orbit
A comet of mass $\alpha m$ is heading (``falls'') radially towards the Sun. It
is known that the total mechanical energy of the comet is zero. The comet
crashes into Venus, whose mass is $m$. We further assume that the orbit of
Venus, before the collision, is circular with radius $R_0$. After the crash,
the comet and Venus form a single object, called ``Venus-2''.

\ioaafig{2021_TH_Q14-f1.pdf}

\begin{description}
\item[(14.1)] Find the expression in terms of $M_{sun}$ and $R_0$ for the
  orbital speed, $v_0$, of Venus before the collision. \hfill[1.0\,pt]
\item[(14.2)] Find an expression for the total mechanical energy of Venus in
  its orbit before colliding with the comet. \hfill[1.0\,pt]
\item[(14.3)] Find an expression for the radial velocity, $v_r$, the angular
  momentum, $L$, of ``Venus-2'' immediately after the collision.
  \hfill[10.0\,pt]
\item[(14.4)] Find an expression for the mechanical energy of the combined
  object ``Venus-2'' and express it in terms of energy before the collision,
  $E_i$, and $\alpha$. \hfill[5.0\,pt]
\item[(14.5)] Show that the post-collision orbit of ``Venus-2'' is elliptical
  and determine the semi-major axis of the orbit. \hfill[5.0\,pt]
\item[(14.6)] Determine if the year for the inhabitants of ``Venus-2'' has
  been shortened or lengthened because of collision with the comet. Write the
  ratio between the period of Venus-2 and Venus. \hfill[3.0\,pt]
\item[(14.7)] What should be the value of $\alpha$ such that the
  post-collision orbit of Venus-2 would make it crash in the Sun? We will call
  this as $\alpha_c$ \hfill[5.0\,pt]
\item[(14.8)] A comet with $\alpha = \alpha_c$ collided with Venus. Calculate
  the percentage change in the magnitude of Venus' velocity ($\delta v$) and
  the change in the direction of the velocity vector ($\delta\theta$)
  immediately after the collision. \hfill[5.0\,pt]
\end{description}

\probhead{2.37}{Mars}{2021}{Bogot\'a (remote)}{TH Q3}{10}
% source id: 2021_TH_Q3   tag: parabolic approach, braking burn
A spacecraft of mass $m = 5.0 \times 10^4$\,kg approaches in a parabolic orbit
$AB$, with respect to Mars. When the spacecraft reaches point $B$ of least
distance to the center of Mars, $r_B = 6.8 \times 10^6$\,m, it undergoes an
instantaneous deceleration using its rockets and goes into a perfectly
calculated orbit so that it will touch the Martian surface exactly at point
$C$, diametrically opposite $B$, as shown in the figure.

\ioaafig{2021_TH_Q3-f1.pdf}

\begin{description}
\item[(3.1)] Determine the speed ($\mathrm{km\,s}^{-1}$) of the spacecraft at
  point $B$ just before the deceleration. \hfill[3.0\,pt]
\item[(3.2)] Calculate the total energy ($J$) of the spacecraft as it is
  moving between points B and C. \hfill[4.0\,pt]
\item[(3.3)] Calculate the speed ($\mathrm{km\,s}^{-1}$) of the spacecraft at
  point $C$. \hfill[3.0\,pt]
\end{description}

\probhead{2.38}{Pluto Satellites}{2021}{Bogot\'a (remote)}{TH Q9}{15}
% source id: 2021_TH_Q9   tag: tidal + centrifugal gravity on Charon
The mass of Charon, the biggest satellite of Pluto, is 1/8th the mass of
Pluto. Both bodies move in a circular orbit around a common center of mass. In
addition, they both are tidally-locked. The distance between the center of the
planet and the center of the satellite is $R = 19\,640$\,km and radius of the
satellite is $r = 593$\,km.

Let $g_0$ be the gravitational acceleration on the surface of Charon due only
to its mass. Let A be the point on Charon surface directly facing Pluto, and B
the point diametrically opposite. Compute the percentage difference between
gravitational acceleration at A and B respect to $g_0$. \hfill[15.0\,pt]

\probhead{2.39}{Co-orbital satellites}{2022}{Kutaisi, Georgia}{TH 12}{50}
% source id: 2022_TH_12   tag: co-orbital horseshoe dynamics
This question applies a method of determining the masses of two approximately
co-orbital satellites developed by Dermott and Murray in 1981.

Suppose that two small satellites of masses $m_1$ and $m_2$ are approximately
co-orbital (moving on very similar orbits) around a large central body of mass
$M$, with $m_1,\,m_2 \ll M$. At any instant, the orbits of the satellites may
be approximated as circular Keplerian orbits with radii $r_1$ and $r_2$
respectively, although $r_1$ and $r_2$ will vary slightly over time due to the
mutual gravitational interaction between the satellites.

Figure 4 depicts the shapes of the orbits in the rotating reference frame with
zero angular momentum, centred on the central body. We denote by $\theta$ the
angle $\measuredangle m_1 M m_2$, while $R$, $x_1$, and $x_2$ denote the mean
orbital radius and radial deviations of the satellites.

Throughout this problem, write all answers in an inertial reference frame.\\
\emph{Hint:} $(1+x)^\alpha \approx 1 + \alpha x$ for $\alpha x \ll 1$

\ioaafig{2022_TH_12-f1.pdf}

First we will determine the value of $\frac{m_1}{m_2}$.

\begin{description}
\item[(a)] Write down the angular momentum $L_i$ of the satellite with mass
  $m_i$ when its circular orbit has radius $r_i$. \hfill(3\,pt)
\item[(b)] The satiletes total angular momentum $L_1 + L_2$ is conserved. Let % sic: "satiletes" in source
  $x_1,\,x_2 \ll R$ be the distances as shown in Figure 4. Find a simple
  relation between the ratios $\frac{m_1}{m_2}$ and $\frac{x_1}{x_2}$.
  \hfill(8\,pt)

  Now, we will try and determine the value of $m_1 + m_2$. For next parts, we
  will use the actual barycenter of the system, which may not be exactly at
  the center of the planet.
\item[(c)] The individual angular momenta of the satellites $m_1$ and $m_2$
  will vary over time due to their gravitational interactions. Show that the
  rate of change of the angular momentum of the second satellite is given by
  \[ \frac{\Delta L_2}{\Delta t} \approx -\frac{G m_1 m_2}{R} h(\theta)
     \quad\text{where}\quad
     h(\theta) = \left[
       \frac{\cos\left(\frac{\theta}{2}\right)}
            {4\sin^2\left(\frac{\theta}{2}\right)} - \sin\theta \right] \]
  \hfill(18\,pt)
\item[(d)] Show that $s = r_2 - r_1$ satisfies
  \[ \frac{\Delta s}{\Delta t} \approx
     -2\sqrt{\frac{G}{MR}}\,(m_1 + m_2)\,h(\theta) \]
  \hfill(8\,pt)
\item[(e)] For the angle $\theta = \measuredangle m_1 M m_2$ as indicated on
  Figure 4, find an expression for
  $\dfrac{\Delta\theta}{\Delta t}$ \hfill(5\,pt)
\item[(f)] Using the results above, find the relation between $\Delta s$ and
  $\Delta\theta$. \hfill(2\,pt)
\item[(g)] After integrating the expression above, we will obtain the result,
  \[ \overline{x}^2 \approx \frac{4R^2}{3}\,\frac{m_1 + m_2}{M}
     \left( \frac{1}{\sin\left(\frac{\theta_{\min}}{2}\right)}
       - 2\cos\theta_{\min} - 3 \right) \]
  where $\overline{x} = \dfrac{x_1 + x_2}{2}$.

  Epimetheus ($m_1$) and Janus ($m_2$) are two approximately co-orbital moons
  of Saturn. Detailed observations of their orbits have been performed by the
  Voyager 1 and Cassini spacecraft, which found that $R = 150\,000$\,km,
  $x_1 = 76$\,km and $x_2 = 21$\,km. The minimum distance between Janus and
  Epimetheus is $13\,000$\,km. The mass of Saturn is known to be
  $5.7 \times 10^{26}$\,kg. Estimate the masses of Epimetheus and Janus.
  \hfill(6\,pt)
\end{description}

\probhead{2.40}{Ring of a planet}{2022}{Kutaisi, Georgia}{TH 8}{20}
% source id: 2022_TH_8   tag: gravity of annular disk, oscillations
We are given a flat disk of mass M with inner radius $r$ and the outer radius
$R$.

\begin{description}
\item[(a)] A point mass $m$ is located on the symmetry axis of the disk at a
  distance $x$ from the plane of the disk (We assume throughout that mass $m$
  stays on the symmetry axis, and is free to move in the vertical direction.).
  What is the gravitational force exerted on the point mass? \hfill(10\,pt)

  \ioaafig{2022_TH_8-f1.pdf}

  \emph{Hint 1:} You may denote surface density by the symbol $\sigma$ and
  calculate force exerted by a small surface element $\Delta S$ subtending a
  solid angle $\Delta\Omega$ with $m$

  \emph{Hint 2:} The area cut by cone with opening angle $2\theta$ on a sphere
  of radius $R$:
  \[ S = 2\pi R^2 (1 - \cos(\theta)) \]

  \ioaafig{2022_TH_8-f2.pdf}

\item[(b)] what will be the frequency of the small oscillations of the system
  ($x \ll r$). \hfill(10\,pt)
\end{description}

\probhead{2.41}{DART}{2023}{Katowice, Poland}{TH Q12}{40}
% source id: 2023_TH_Q12   tag: impact changes Dimorphos orbital period
The Double Asteroid Redirection Test (DART) was a NASA mission to evaluate a
method of planetary defense against near-Earth objects. The spacecraft hit
Dimorphos, a moon of the asteroid Didymos, to study how the impact affected
its orbit.

\begin{description}
\item[(a)] Calculate the expected orbital period change (in minutes), assuming
  that the collision was head-on, central, and perfectly inelastic.

  Assume that before the impact Dimorphos orbited Didymos on a circular orbit
  with a period of $P = 11.92$\,h. The masses of Dimorphos and Didymos are
  $m = 4.3 \times 10^9$\,kg and $M = 5.6 \times 10^{11}$\,kg, respectively.
  The mass and speed of the DART spacecraft relative to Dimorphos at a moment
  of impact were $m_{\mathrm{s}} = 580$\,kg and
  $v_{\mathrm{s}} = 6.1$\,km\,s$^{-1}$. Neglect the gravitational influence of
  other bodies. \hfill(20 points)
\item[(b)] In reality, the orbital period of Dimorphos was observed to be
  changed by $\Delta P_0 = -33$\,min. This is due to the momentum transfer
  associated with the recoil of the ejected debris: the spacecraft was
  absorbed by the asteroid, but the impact excavated some material from the
  asteroid and ejected it into space. Calculate the momentum of the ejected
  debris and express it as a fraction of the momentum of Dimorphos before the
  collision. You can assume that the mass of the ejected material is much
  smaller than the mass of Dimorphos. \hfill(15 points)
\item[(c)] Calculate the velocity change (in mm\,s$^{-1}$) of Dimorphos as a
  result of the impact, taking into account the effect of the ejected debris.
  \hfill(5 points)
\end{description}

\probhead{2.42}{Ground Tracks}{2024}{Rio de Janeiro, Brazil}{TH T11}{75}
% source id: 2024_TH_T11   tag: Sun-synchronous and Tundra orbits
The projection of a satellite's orbit onto the Earth's surface is called its
ground track. At a given instant, one can imagine a radial line drawn outward
from the centre of the Earth to the satellite. The intersection between the
Earth's spherical surface and this radial line is a point on the ground track.

The location of this point is specified by its geocentric latitude and
longitude. The ground track is then essentially the figure traced by this
point as the satellite moves around the Earth.

\textbf{Part I: Sun-Synchronous Orbits (25 points)}

It is particularly interesting to analyse the ground track of a so-called
Sun-Synchronous orbit. This is a nearly polar orbit around a planet where the
satellite passes over any given point on the planet's surface at the same mean
local solar time. This property is especially interesting for satellite
imaging, ensuring similar illumination conditions over different days.

The figure below shows the ground track of a satellite in a Sun-Synchronous
orbit. Its inclination angle ($i$) --- the angle between the satellite orbital
plane and the Earth's equatorial plane --- falls within the range
$90^\circ < i < 180^\circ$. The graph depicts five complete orbits of the
satellite.

\ioaafig{2024_TH_T11-f2.pdf}

For the questions in Part I, assume that the Earth's orbit around the Sun is
circular.

\begin{description}
\item[(a)] (3 points) Determine the nodal precession rate for the orbit in
  rad/s.
\item[(b)] (8 points) Based on the ground track shown in Figure 2, determine
  the inclination of the satellite's orbit (in degrees) and estimate its
  orbital period (in minutes). Consider that the orbital period of the
  satellite is shorter than one sidereal day.
\item[(c)] (2 points) Calculate the semi-major axis $a$ of the orbit in km.
\item[(d)] (1 point) Determine the number of orbits completed by the satellite
  until it returns to the same position on Earth.
\item[(e)] (11 points) As seen in Figure 2, the ground track crosses the
  Brazilian city of Macei\'o
  $(\phi, \lambda) = (9.7^\circ S; 35.7^\circ W)$ and also Chorz\'ow
  $(\phi, \lambda) = (50.3^\circ N; 19.0^\circ E)$, in Poland. Knowing that
  the ground track crosses Macei\'o at noon (local time), determine the local
  time that the satellite track crosses Chorz\'ow. Hint: specifically for this
  task, you may neglect the effects of nodal precession.
\end{description}

\textbf{Part II: Tundra orbits (50 points)}

A Tundra orbit is a type of geosynchronous elliptical orbit characterised by a
high inclination. The apogee is positioned over a specific geographic region,
allowing for prolonged visibility and coverage over that area. This orbit
ensures that a satellite spends the majority of its orbital period over the
northern --- or southern --- hemisphere, making it particularly useful for
communications and weather observation over high-latitude regions.

The image below represents the ground track of a satellite in a Tundra orbit
with an argument of perigee equal to $270^\circ$. The satellite orbits the
Earth in the same direction as its rotation. For the following items, you can
ignore the effects of the Earth's oblateness.

\ioaafig{2024_TH_T11-f3.pdf}

\begin{description}
\item[(f)] (4 points) Based on the graph above, give the inclination of the
  satellite's orbit $i$ (in degrees), its orbital period $T$ (in minutes), and
  its semi-major axis $a$ (in km).
\item[(g)] (12 points) Show that the time a satellite spends in the northern
  hemisphere is given by
  \[ T' = \left( \frac{1}{2} + \frac{\sin^{-1}(e)}{\pi}
     + \frac{e}{\pi}\cdot\sqrt{1 - e^2} \right) T \]
  where $e$ is the eccentricity of the orbit and $T$ is its orbital period.
\item[(h)] (10 points) Estimate numerically the eccentricity $e$ of its orbit.
  You can consider that the eccentricity is so small that
  $\sin(e) \approx e$ and $e^2 \ll 1$.
\item[(i)] (18 points) From the ground track, we can observe that the
  satellite exhibits retrograde motion in both its northern and southern
  hemisphere trajectories. Find the true anomaly (in degrees) of the satellite
  at the beginning and end of its retrograde motion in the southern
  hemisphere.
\item[(j)] (6 points) It is also noticeable that the ground track of a Tundra
  orbit has the shape of a figure-8, similar to an analemma, so that the
  satellite passes over the same point on Earth in a single orbit. Calculate
  the minimum eccentricity the orbit would need to have for this property to
  cease occurring. Use the same orbital inclination as the orbit in Figure 3.
\end{description}

\probhead{2.43}{Asteroid}{2024}{Rio de Janeiro, Brazil}{TH T3}{10}
% source id: 2024_TH_T3   tag: radial-fall time, degenerate Kepler orbit
A peculiar asteroid of mass, $m$, was spotted at a distance, $d$, from a star
with mass, $M$. The magnitude of the asteroid's velocity at the time of the
observation was $v = \sqrt{\frac{GM}{d}}$, where $G$ is the universal
gravitational constant. The distance $d$ is much larger than the radius of the
star.

For both of the following items, express your answers in terms of $M$, $d$,
and physical or mathematical constants.

\begin{description}
\item[(a)] (8 points) If the asteroid is initially moving exactly towards the
  star, how long will it take for it to collide with the star?
\item[(b)] (2 points) If the asteroid is instead initially moving exactly away
  from the star, how long will it now take for it to collide with the star?
\end{description}

\section*{Data-analysis problems}

\probhead{2.44}{Minor Planet}{2020}{GeCAA (Estonia, remote)}{DA 2}{50}
% source id: 2020_DA_2   tag: orbit determination, Kepler elements
Table 1 gives ecliptic longitude ($\lambda$) and parallax ($p$) at different
times ($t$), for a certain hypothetical minor planet. The baseline for the
parallax is the diameter of the earth. The time is expressed in years and for
your reference ecliptic longitudes of the Sun ($\lambda_\odot$) for the same
dates are also given the table. Let us assume that the orbital inclination of
this minor planet, with respect to the ecliptic, is negligible and the
eccentricity of the Earth's orbit is negligible.

\begin{description}
\item[(a)] Calculate the coordinates of the minor planet in the heliocentric
  polar coordinate system and put them in a polar plot. The x-axis in the plot
  should be directed towards the initial position of the minor planet. Draw
  the major axis of the orbit of the minor planet.\\
  Identify erroneous observation(s), if any. \hfill[38 points]
\item[(b)] Assuming the heliocentric orbit of the minor planet to be
  elliptical, determine
  \begin{description}
  \item[(i)] the semimajor axis length $a_p$.
  \item[(ii)] eccentricity $e$.
  \item[(iii)] the period $P$.
  \end{description}
  \hfill[6 points]
\item[(c)] Estimate the errors in the values of $P$, $a_p$ and $e$.
  \hfill[6 points]
\end{description}

\begin{center}
Table 1: Minor planet data

\begin{tabular}{cccc}
\hline
$t$ & $\lambda$ & $\lambda_\odot$ & $p$ \\
{[year]} & {[$^\circ$]} & {[$^\circ$]} & {[$''$]} \\
\hline
2012.3 & 336.73 & 40.95 & 3.82 \\
2012.6 & 3.44 & 134.83 & 7.24 \\
2012.9 & 50.71 & 242.08 & 7.09 \\
2013.4 & 94.52 & 64.84 & 2.40 \\
2013.6 & 121.40 & 134.59 & 2.16 \\
2013.9 & 154.31 & 241.82 & 2.75 \\
2014.2 & 25.33 & 353.29 & 3.16 \\
2014.5 & 148.51 & 99.04 & 1.99 \\
2014.8 & 176.26 & 205.45 & 1.83 \\
2015.0 & 216.33 & 280.19 & 2.03 \\
2015.3 & 187.5 & 28.55 & 2.897 \\
\hline
\end{tabular}
\end{center}

\clearpage
\section*{Solutions to Chapter 2}

\solhead{2.1}{}
% source id: 2007_TH_1.6
According to Kepler's third law we have
\begin{align*}
  (\text{period})^2 &= (\text{constant})(\text{semi-major axis})^3 \\
  T^2 &= (\text{constant})\,a^3
\end{align*}
This constant is $1\,\dfrac{(\text{year})^2}{(\text{A.U.})^3}$ when $T$ is
measured in years and $a$ in A.U.'s.

For this comet we have
\begin{align*}
  a &= \frac{31.5 + 0.5}{2} = 16.0\ \text{A.U.} \\
  T^2 &= (16)^3 = 16 \times 4 \times 4 \times 16 = (64)^2 \\
  T &= 64\ \text{years}
\end{align*}

\solhead{2.2}{}
% source id: 2007_TH_1.7
According to Kepler's second law we have;

The area is swept out at constant rate throughout the orbital motion.
Now, for this comet the area swept out in one orbital period is $\pi ab$ where
$a$ is the semi-major axis and $b$ the semi-minor axis of the orbit.
\begin{align*}
  a &= 16.0\ \text{A.U.} \\
  b^2 &= a^2 - (\text{distance from a focus to centre of ellipse})^2 \\
  &= (16.0)^2 - (16.0 - 0.5)^2 \\
  b &= 3.968\ \text{A.U.} \\
  \pi ab &= 199.5\ (\text{A.U.})^2
\end{align*}
Hence, the area swept out per year is
\[
  \frac{199.5\ (\text{A.U.})^2}{64\ \text{years}}
  = 3.1\ (\text{A.U.})^2/\text{year}
\]

\solhead{2.3}{}
% source id: 2008_TH_S14
The mechanical energy is conserved and has the form \hfill(10\%)
\[
  E = \frac{p^2}{2m} - \frac{G\mathcal{M}m}{r} = \frac{p^2}{2m} - G\mathcal{M}mu
\]
Insert the values for $p$ and $u$ for points A and B and solve for $m$ and $E$.
\hfill(20\%)
\begin{align*}
  E_A &= \frac{p_A^2}{2m} - G\mathcal{M}mu_A \\
  E_B &= \frac{p_B^2}{2m} - G\mathcal{M}mu_B \\
  \frac{p_A^2}{2m} - G\mathcal{M}mu_A &= \frac{p_B^2}{2m} - G\mathcal{M}mu_B \\
  m &= \sqrt{\frac{p_A^2 - p_B^2}{2G\mathcal{M}(u_A - u_B)}} \\
  m &= \sqrt{\frac{(5.2\times 10^{7})^2 - (1.94\times 10^{9})^2}
       {2(6.6726\times 10^{-11})(5.9736\times 10^{24})
        (5.15\times 10^{-8} - 1.94\times 10^{-6})}}
     = 50\ \text{tons.}
\end{align*}
\hfill(20\%)

Thus,
\begin{align*}
  E &= E_A = \frac{p_A^2}{2m} - G\mathcal{M}mu_A \\
  E &= \frac{(5.20\times 10^{7})^2}{2(50000)}
      - (6.6726\times 10^{-11})(5.9736\times 10^{24})(50000)
        (5.15\times 10^{-8}) \\
  &= -1.0\times 10^{12}\ \text{J}
\end{align*}
\hfill(20\%)

To sketch the curve, we use the equation
\[
  E = \frac{p^2}{2m} - G\mathcal{M}mu
  \;\Longrightarrow\;
  u = \frac{1}{2G\mathcal{M}m^2}\,p^2 - \frac{E}{G\mathcal{M}m}
  \;\Longrightarrow\;
  \text{Parabolic curve}
\]
\hfill(10\%)

\ioaafig{2008_TH_S14-s1.pdf}
\hfill(20\%)

\solhead{2.4}{}
% source id: 2008_TH_S8
\ioaafig{2008_TH_S8-s1.pdf}
\hfill(10\%)

D in the figure is the center of the Moon. The Moon's gravitational
acceleration at A and at center of the Earth C are
\begin{align*}
  \vec{g}_A &= -\frac{G\mathcal{M}}{r_A^2}\,\hat{r}_A \\
  \vec{g}_C &= -\frac{G\mathcal{M}}{r_C^2}\,\hat{r}_C
             = -\frac{G\mathcal{M}}{r^2}\,\hat{r}_C
\end{align*}
Their directions are toward the Moon. $\mathcal{M}$ is the mass of the Moon.
$r_{\mathrm{A}} = \mathrm{AD}$ and $r_{\mathrm{C}} = \mathrm{CD} = r$.
$\hat{r}_A$, $\hat{r}_C$ are unit vectors from the center of the Moon toward
A and C, respectively. \hfill(20\%)

The Moon's gravitational acceleration at A according to observers is
\[
  \vec{g}\,'_A = \vec{g}_A - \vec{g}_C
  = -\left( \frac{G\mathcal{M}}{r_A^2}\,\hat{r}_A
          - \frac{G\mathcal{M}}{r^2}\,\hat{r}_C \right)
  = -\left( \frac{G\mathcal{M}}{r_A^3}\,\vec{r}_A
          - \frac{G\mathcal{M}}{r^3}\,\vec{r}_C \right)
\]
Since $\vec{r}_A = \vec{r}_C + \vec{R}$, $\vec{R}$ is the position of A with
respect to the Earth's center of mass
\[
  \vec{g}\,'_A = \vec{g}_A - \vec{g}_C
  = -\frac{G\mathcal{M}}{r_A^3}\,\vec{R}
    - \left( \frac{G\mathcal{M}}{r_A^3}
           - \frac{G\mathcal{M}}{r^3} \right)\vec{r}_C
\]
\hfill(20\%)

The first term is in radial direction. The tangential or horizontal component
of the second term is equal to the horizontal component of the acceleration of
sea water at A:
\begin{align*}
  g'_A(\text{horizontal})
  &= \left( \frac{G\mathcal{M}}{r^3} - \frac{G\mathcal{M}}{r_A^3} \right)
     r\,\sin\varphi \\
  &= G\mathcal{M}\left( \frac{1}{r^3}
     - \frac{1}{(r^2 + R^2 - 2Rr\,\cos\varphi)^{3/2}} \right) r\,\sin\varphi \\
  &= \frac{G\mathcal{M}}{r^2}\left( 1
     - \frac{1}{[\,1 + (R/r)^2 - 2(R/r)\cos\varphi\,]^{3/2}} \right)\sin\varphi
\end{align*}
\hfill(50\%)

\solhead{2.5}{}
% source id: 2009_TH_P12
The rate of change of solar mass could be estimated from solar luminosity:
\hfill(2 points)
\[
  L_\odot = -\frac{\Delta E}{\Delta t} = -\frac{\Delta M c^2}{\Delta t}
          = -\dot{M}c^2
\]
\[
  \dot{M} = -\frac{L_\odot}{c^2}
          = -\frac{3.83\times 10^{26}}{(3.00\times 10^{8})^2}
          = -4.26\times 10^{9}\,(kg\,s^{-1})
\]
Newton second law will give us:
\[
  \frac{v^2}{r} = \frac{GM}{r^2} \;\Longrightarrow\; v^2 = \frac{GM}{r}
\]
\hfill(2 points)

where $v$ and $r$ are orbital velocity and orbital radius of the Earth.
From conservation of angular momentum:
\[
  l = rmv \;\Longrightarrow\; v = \frac{l}{mr}
\]
\hfill(2 points)

Where $m$ and $l$ are the Earth mass and angular momentum which are constant:
\begin{align*}
  \frac{l^2}{m^2r^2} = \frac{GM}{r}
  \Rightarrow r = \frac{l^2}{GMm^2}
  &\Rightarrow \dot{r} = -\frac{l^2}{Gm^2}\,\frac{\dot{M}}{M^2}
   = -\frac{r^2m^2v^2}{GM^2m^2}\,\dot{M} \\
  \Rightarrow \dot{r} = -\frac{(GM/r)r^2}{G}\,\frac{\dot{M}}{M^2}
   = -r\,\frac{\dot{M}}{M}
  &\Rightarrow \frac{\dot{r}}{r} = -\frac{\dot{M}}{M}
   \Rightarrow \dot{r} = -r\,\frac{\dot{M}}{M} \\
  \Rightarrow \dot{r}
   = -\frac{1.50\times 10^{11}\times 4.26\times 10^{9}}{1.99\times 10^{30}}
  &= 3.21\times 10^{-10}\,(m/s)
\end{align*}
\hfill(2 points)
\[
  \Delta r = 3.19\times 10^{-10}\times 100\times 86400\times 365.24
  % sic: source uses 3.19e-10 here although \dot{r} was 3.21e-10 above
  = 1.01\,m \quad (for\ 100\ years)
\]
\hfill(2 points)

\solhead{2.6}{}
% source id: 2009_TH_P14
\ioaafig{2009_TH_P14-s1.pdf}
In the figure, the observer is at $O$ and the satellite is in $S$ the angle
$\angle EOS$ will be
\[
  \angle EOS = 180 - z = 180 - 46.0 = 134^\circ
\]
\hfill(2 points)
\[
  \delta = 180 - \varphi - \angle EOS = 10.4^\circ
\]
\hfill(3 points)

In $EOS$ triangle we have:
\[
  \frac{R_S}{\sin(\angle EOS)} = \frac{R}{\sin\delta}
\]
\hfill(3 points)
\[
  \frac{R_S}{R} = \frac{\sin(\angle EOS)}{\sin\delta} = 3.98
\]
\hfill(2 points)

\solhead{2.7}{High Altitude Projectile}
% source id: 2009_TH_P16
\begin{description}
\item[a)] Total energy of the projectile is
\[
  E = \tfrac{1}{2}mv_\circ^2 - \frac{GMm}{R} = -\frac{GMm}{2R} < 0
\]
$E < 0$ means that orbit might be ellipse or circle. As $\theta > 0$, the
orbit is an ellipse. Total energy for an ellipse is
\[
  E = -\frac{GMm}{2a}
\]
Then
\[
  a = R
\]
\hfill(7 points)

\item[b)] In figure (1) we have
\begin{align*}
  OA + O'A &= 2a \\
  O'A &= a
\end{align*}
\ioaafig{2009_TH_P16-s1.pdf}
In $OAO'$ triangle it is obvious that
\[
  OC = CO'
\]
Then $C$ must be the center of the ellipse with the initial velocity vector
$v_\circ$ parallel to the ellipse major-axis $(LH)$.

In figure (2)
\[
  HM = CH - CM = a - (R - R\sin\theta) = R - R + R\sin\theta
     = R\sin\theta = \frac{R}{2}
\]
\hfill(15 points)

\item[c)] Range of the projectile is $\overset{\frown}{AB}$
\[
  \overset{\frown}{AB} = 2\left(\frac{\pi}{2} - \theta\right)R
  = (\pi - 2\theta)R = \frac{2\pi}{3}R
\]
\hfill(6 points)

\item[d)]
\ioaafig{2009_TH_P16-s2.pdf}
Start with ellipse equation in polar coordinates
\[
  r = \frac{a(1 - e^2)}{1 + e\cos\varphi}
\]
For point A
\[
  R = \frac{R(1 - e^2)}{1 - e\cos(\frac{\pi}{2} + \theta)}
\]
\[
  e = \sin\theta = \frac{1}{2}
\]
\hfill(5 points)

\item[e)] Using Kepler's second law
\ioaafig{2009_TH_P16-s3.pdf}
\begin{align*}
  \frac{\Delta S}{S_0} &= \frac{\Delta T}{T} \\
  \Delta S = S_{AOBH} &= S_{\Delta AOB} + \frac{S_0}{2} \\
  &= 2\times\frac{bae}{2} + \frac{\pi ab}{2}
   = bae + \frac{\pi ab}{2} \\
  \frac{\Delta S}{S} &= \frac{bae + \frac{\pi ab}{2}}{\pi ab}
   = \frac{e + \frac{\pi}{2}}{\pi} = \frac{0.5 + \frac{\pi}{2}}{\pi}
\end{align*}
Kepler's third law
\[
  T = \sqrt{\frac{4\pi^2 R^3}{GM}} = 84.5\ min
\]
\[
  \Delta T = T\times\frac{0.5 + \frac{\pi}{2}}{\pi} = 55.7\ min
\]
\hfill(12 points)
\end{description}

\solhead{2.8}{}
% source id: 2009_TH_P3
We must compare the jumping speed of a normal human with escape velocity of
the planet. A normal human can jump up to 50\,cm then his initial velocity is
\begin{align*}
  v &= \sqrt{2gh} = \sqrt{2\times 9.81\times 0.5} \\
  v &= 3.13\ m\,s^{-1}
\end{align*}
\hfill(3 points)

Comparing this velocity with escape velocity of the planet
\[
  v = \sqrt{\frac{2GM}{R}}
\]
\hfill(3 points)
\begin{align*}
  v^2 &= \frac{2GM}{R} = \frac{2G}{R}\,\frac{4\pi}{3}\,\rho\,R^3 \\
      &= \frac{8\pi G}{3}\,\rho\,R^2
\end{align*}
\hfill(2 points)
\begin{align*}
  R^2 &= \frac{3v^2}{8\pi G\rho} \\
  R &= 2\times 10^{3}\ m
\end{align*}
\hfill(2 points)

\solhead{2.9}{}
% source id: 2009_TH_P6
To solve this problem we must calculate gravitational potential at the center
of the cloud, letting $\phi(\infty) = 0$. For a uniform density and spherical
mass distribution we have

\ioaafig{2009_TH_P6-s1.pdf}
\[
  \frac{1}{2}\,m\,v^2(r=R) - \frac{GMm}{R} = E
  = \frac{1}{2}mv^2(r=0) + \varphi(0)
\]
\hfill(4 points)
\[
  v(r=R) = 0 \qquad \& \qquad v(r=0) = \sqrt{\frac{GM}{R}}
\]
\hfill(1 point)
\[
  \varphi(0) = -\frac{1}{2}mv_0^2 - \frac{GmM}{R}
\]
\hfill(1 point)
\[
  \varphi(0) = \frac{-1}{2}\,m\left(\sqrt{\frac{GM}{R}}\right)^2
             - \frac{GmM}{R}
\]
\hfill(2 points)
\[
  \varphi(0) = \frac{-3}{2}\times\frac{GMm}{R}
\]
To escape from the cloud, the particle should have total energy equal to zero
\[
  E = \tfrac{1}{2}m{v_e}^2 + \phi(r=0)
    = \tfrac{1}{2}m{v_e}^2 - \tfrac{3}{2}\left(\frac{GMm}{R}\right) = 0
\]
\[
  {v_e}^2 = \frac{3GM}{R}
  \qquad\rightarrow\qquad
  v_e = \sqrt{\frac{3GM}{R}}
\]
\hfill(2 points)

\solhead{2.10}{}
% source id: 2010_TH_L16
Solution: (1) 10 points; (2) 20 points

(1) The orbit of the spacecraft is a parabola, this suggests that the
(specific) energy with respect to the Sun is initially
\[
  \varepsilon = 1/2\,v_{\max}^2 + U(r_E) = 0
\]
\hfill(2 points)

and
\[
  v_{\max} = \sqrt{2U} = \sqrt{2k_{sun}/r_E}
\]
The angular momentum is
\[
  l = r_E v_{\max} = \sqrt{2k_{sun}r_E}
\]
\hfill(2 points)

When the spacecraft cross the orbit of the Mars at 1.5\,AU, its total velocity
is
\[
  v = \sqrt{2U} = \sqrt{2k_{sun}r_M} = \sqrt{\frac{2}{3}}\,v_{\max}
  % sic: source prints \sqrt{2k_{sun}r_M}; dimensionally it should be \sqrt{2k_{sun}/r_M}
\]
This velocity can be decomposed into $v_r$ and $v_\theta$, using angular
momentum decomposition,
\[
  r_M v_\theta = l = r_E v_{\max}
\]
\hfill(2 points)

So,
\[
  v_\theta = \frac{r_E}{r_M}\,v_{\max} = \frac{2}{3}\,v_{\max}
\]
\hfill(2 points)

Thus the angle is given by
\[
  \cos\psi = \frac{v_\theta}{v} = \sqrt{\frac{r_E}{r_M}} = \sqrt{\frac{2}{3}}
\]
or
\[
  \psi = 35.26^\circ
\]
\hfill(2 points)

Note: students can arrived at the final answer with conservation of angular
momentum and energy, full mark.

(2) The Mars would be moving on the circular orbit with a velocity
\[
  v_M \equiv \sqrt{\frac{k_{sun}}{r_M}} = \sqrt{\frac{2}{3}}\,v_E
  = 24.32\,km/s
\]
\hfill(3 points)

from the point of view of an observer on Mars, the approaching spacecraft has
a velocity of
\[
  \overrightarrow{v_{rel}} = \vec{v} - \vec{v}_M
\]
\hfill(2 points)

Now
\[
  \vec{v} = v\sin\psi\,\hat{r} + v_\theta\,\hat{\theta}
\]
\hfill(2 points)

with
\[
  \sin\psi = \sqrt{1 - \cos^2\psi} = \frac{1}{\sqrt{3}}
\]
So
\begin{align*}
  \overrightarrow{v_{rel}}
  &= v\sin\psi\,\hat{r} + (v_\theta - v_M)\,\hat{\theta} \\
  &= \frac{1}{\sqrt{3}}\sqrt{\frac{2k_{sun}}{r_M}}\,\hat{r}
   + \Bigl(\frac{2}{3}\sqrt{\frac{2k_{sun}}{r_E}}
   - \sqrt{\frac{k_{sun}}{r_M}}\Bigr)\hat{\theta} \\
  &= \sqrt{\frac{2k_{sun}}{3r_M}}\,\hat{r}
   + \Bigl(\frac{2}{\sqrt{3}} - 1\Bigr)\sqrt{\frac{k_{sun}}{r_M}}\,\hat{\theta} \\
  &= \sqrt{\frac{k_{sun}}{r_M}}\,(0.8165\,\hat{r} + 0.1547\,\hat{\theta})
\end{align*}
\hfill(8 points)

The angle between the approaching spacecraft and Sun seen from Mars is:
\begin{align*}
  \tan\theta &= \frac{0.1547}{0.8165} = 0.1894 \\
  \theta &= 10.72^\circ
\end{align*}
\hfill(3 points)

The approaching velocity is thus
\[
  v_{rel} = \sqrt{\frac{2}{3} + \Bigl(\frac{2}{\sqrt{3}} - 1\Bigr)^2}
            \sqrt{\frac{k_{sun}}{r_M}} = 20.21\,km/s
\]
\hfill(2 points)

\solhead{2.11}{}
% source id: 2010_TH_S6
The mass of the asteroid is
\[
  m_1 = \frac{4}{3}\pi r^3 \rho = 1.23\times 10^{13}\,\mathrm{kg}
\]
\hfill(2 points)

Since $m_2 \ll m_1$, $m_2$ can be omitted,

Then
\[
  v = \sqrt{\frac{Gm_1}{r}} = 0.864\,\mathrm{m/s}
\]
\hfill(3 points)

It is the first cosmological velocity of the asteroid.
If the velocity of the astronaut is greater then $v$, he will escape from the
asteroid.

The astronaut must be at $v_2$ if he wants to complete a circle along the
equator of the asteroid on foot within 2.2 hours, and
\[
  v_2 = \frac{2\pi\times(2200/2)\,\mathrm{m}}{2.2\times 3600\,\mathrm{s}}
      = 0.873\,\mathrm{m/s}
\]
\hfill(3 points)

Obviously $v_2 > v$

So the answer should be ``NO''. \hfill(2 points)

\solhead{2.12}{}
% source id: 2011_TH_S1
Perihelium distance is much smaller than aphelium distance.
Thus, semimajor axis of the comet orbit $a \approx 17500$\,AU. \hfill[1]

Application of the III Kepler's Law to get the comet orbital period \hfill[4]

Actual calculations:
$T = a^{2/3} \approx 2.32\cdot 10^{6}$ ($T$ in years)
% sic: source prints a^{2/3}; Kepler's third law gives T = a^{3/2}
\hfill[4]

The sought solution is equal to one half of $T$, i.e.\ approx.\
$1.16\cdot 10^{6}$ years \hfill[1]

\solhead{2.13}{}
% source id: 2011_TH_S3
Sum of potential and kinetic energy of Voyager is
$E = -G\,m\,M_\odot/r + mv^2/2$ \hfill[3]

Where $m$ -- mass of Voyager, $r = 116.406$\,AU, $v = 17.062$\,km/s.

Substituting numerical values into energy equation, \hfill[3]

we get $E > 0$, therefore the orbit is hyperbolic \hfill[1]

Brightness of the Sun is $m(r) = m_\odot + 5\log(r)$ \hfill[2]

$\approx -16.5$\,mag. \hfill[1]

\solhead{2.14}{}
% source id: 2011_TH_S4
The orbit of Phobos and the rotation of the planet are in the same direction,
so the angular velocity of the moon relative to Mars, $\omega_{P/M}$, is equal
to the difference
\[
  \omega_{P/M} = \omega_M - \omega_P
\]
\hfill[2]

where $\omega_M$ and $\omega_P$ are the angular rotation velocity of Mars and
the angular orbital velocity of Phobos respectively.
For $\omega_M$ we have $\omega_M = 2\pi/T_M$ where $T_M$ is the period \hfill[1]

The time above the horizon, $t$, is found from the relation (see drawing):
\[
  t = 2\alpha/\omega_{P/M}
\]
\hfill[2]

where $\alpha = R_M/R_P$.
% sic: source omits the cosine; the geometry (and the numerical value below) give cos(alpha) = R_M/R_P
\hfill[1]

\ioaafig{2011_TH_S4-s1.pdf}

Find the angular velocity of Phobos $\omega_P$ assuming gravity is equal to
the centripetal force:
\[
  \omega_P = \mathrm{sqrt}(\,GM_M/R^{3/2}_{\;P})
  % sic: source prints sqrt(GM_M / R_P^{3/2}); dimensionally it should be sqrt(GM_M / R_P^3)
\]
\hfill[2]

Substituting:
\begin{align*}
  \omega_M &= 7.0882\cdot 10^{-5}\,\mathrm{s}^{-1}, \\
  \omega_P &= 2.2784\cdot 10^{-4}\,\mathrm{s}^{-1}, \\
  \omega_{P/M} &= -1.5696\cdot 10^{-4}\,\mathrm{s}^{-1}
\end{align*}
where ``$-$'' means that the moon ``leads'' the rotation of the planet;
substitute the absolute value into the equation,

thus
$\alpha = 68^\circ\!.794 = 1.2007$\,rad.
and $t \approx 15300\,\mathrm{s} = 4h\,15m$. \hfill[2]

\solhead{2.15}{}
% source id: 2011_TH_S6
Calculation of the moment of inertia of Earth \hfill[2]

The moment of inertia is
\[ I = \frac{2}{5} m_{\mathrm{E}} R_{\mathrm{E}}^2 \]

Equation of motion for this case
\[ K = I\cdot\varepsilon \]
\hfill[2]

where $\varepsilon$ is angular acceleration and $K$ is moment of force.

Calculation of angular acceleration \hfill[1]

Acceleration is equal to $\varepsilon = K/I$. From the assumption, the
acceleration $\varepsilon$ is constant and we find its numerical value equal
to $6\cdot10^{-22}\,\mathrm{s}^{-2}$. % sic: the minus sign of the exponent is illegible in the scan

Calculation of the change of $\varepsilon$ during assumed period of time % sic: change of omega, not epsilon
\hfill[1]

Multiplying this value by $1.9\cdot10^{16}$\,s [600\,My] we obtain
$\Delta\omega = 1.14\cdot10^{-5}\,\mathrm{s}^{-1}$.

Calculation of the angular frequency in the past \hfill[2]

Using the definition of the angular frequency one can calculate its recent
value $2\pi/86164\,\mathrm{s}^{-1} = 7.29\cdot10^{-5}\,\mathrm{s}^{-1}$.
Subtracting from it calculated above $\Delta\omega$ one gets
$\omega_{\mathrm{past}} = 6.15\cdot10^{-5}\,\mathrm{s}^{-1}$.

Calculation of the day length in the past and number of days in the ancient
year \hfill[2]

Having angular frequency one can get ancient period of Earth rotation.
Dividing sidereal year length by the obtained period one gets number of days
in the ancient year equal to about 420 days.

\solhead{2.16}{}
% source id: 2011_TH_S7
Limiting ourself to the orbits with peri and apogeum in the explosion point,
the angle $\beta$ should be equal to $90^\circ$ or $270^\circ$, while the
angle $\alpha$ should be:

\begin{description}
\item[1).] $\alpha \in \langle 0^\circ;\, 90^\circ)$ \ or \
  $\alpha \in (270^\circ;\, 360^\circ\rangle$, \hfill(2 points)
\item[2).] $\alpha = 90^\circ$ \ or \ $\alpha = 270^\circ$, \hfill(2 points)
\item[3).] $\alpha \in (90^\circ;\, 120^\circ)$ \ or \
  $\alpha \in (240^\circ;\, 270^\circ)$, \hfill(2 points)
\item[4).] $\alpha = 120^\circ$ \ or \ $\alpha = 240^\circ$, \hfill(2 points)
\item[5).] $\alpha \in (120^\circ;\, 180^\circ)$ \ or \
  $\alpha \in (180^\circ;\, 240^\circ)$, \hfill(2 points)
\end{description}
respectively.

Taking into account other kind of orbits (without peri and apogeum) one has to
make the angle $\beta$ free (unfixed).

\solhead{2.17}{}
% source id: 2012_TH_S5
\emph{Fact 1.} $\Delta\phi = 1^\circ = \frac{\pi}{180}$, $L = 111$\,km:
\[ R_\oplus = \frac{L}{\Delta\phi} \]
\hfill(2 pt, relations of fact 1)

\emph{Fact 2.} $g = 9.8\,\mathrm{m/s^2}$:
\[ g = \frac{GM_\oplus}{R_\oplus^2} = \frac{4\pi G \rho_\oplus R_\oplus}{3}
   \quad\therefore\quad
   \rho_\oplus = \frac{3g\,\Delta\phi}{4\pi G L} \]
\hfill(2 pt, relations of fact 2)

\emph{Fact 3.} $T_\oplus = 1$ year:
\[ \frac{T_\oplus^2}{a_\oplus^3} = \frac{4\pi^2}{GM_\odot}
   \quad\therefore\quad
   T_\oplus^2 = \frac{3\pi a_\oplus^3}{G\rho_\odot R_\odot^3},
   \qquad M_\odot = \frac{4}{3}\pi \rho_\odot R_\odot^3 \]
\hfill(2 pt, relations of fact 3)

\emph{Fact 4.} $\alpha = 30'$:
\[ \frac{\alpha}{2} \approx \frac{R_\odot}{a_\oplus}
   \quad\therefore\quad
   T_\oplus^2 = \frac{3\pi a_\oplus^3}{G\rho_\odot R_\odot^3}
   = \frac{24\pi}{G\rho_\odot \alpha^3} \]
\hfill(2 pt, relations of fact 4)

Therefore
\[ \frac{\rho_\oplus}{\rho_\odot}
   = \frac{\dfrac{3g\,\Delta\phi}{4\pi G L}}{\dfrac{24\pi}{G T_\oplus^2 \alpha^3}}
   = \frac{g\,\Delta\phi\,\alpha^3 T_\oplus^2}{32\pi^2 L} = 3.23 \]
\hfill(2 pt, answer)

\solhead{2.18}{}
% source id: 2013_TH_L2
\begin{description}
\item[(1)] For $\Delta t_b = 2\,\mathrm{min} = 120$\,sec, the distance
  travelled by a light pulse is $R_1 = c \times \Delta t$ or
  $R_1 = 0.24$\,AU, while for $\Delta t_a = 39$\,min the maximum distance is
  $R_2 = c \times \Delta t$ or $R_2 = 4.67$\,AU. \hfill(4 points)

  Since the orbits do not intersect and the $R_2$ exceeds by far 1\,AU, the
  orbit of Seneca is exterior to that of the Earth. $R_1$ corresponds to the
  minimum relative distance of the two bodies (i.e.\ at opposition), while
  $R_2$ corresponds to the maximum relative distance (i.e.\ at conjunction).

  If $q = a\times(1-e)$ is the perihelion and $Q = a\times(1+e)$ the aphelion
  distance of Seneca, then $R_1 = q - 1$\,AU (the minimum distance of Seneca
  from the Sun minus the semi-major axis of the Earth's orbit), while
  $R_2 = Q + 1$\,AU (the maximum distance of Seneca from the Sun plus the
  semi-major axis of the Earth's orbit). \hfill(6 points)

  Thus,
  \begin{align*}
    a_{\mathrm{Sen}}(1 - e_{\mathrm{Sen}}) &= 1 + R_1 = 1.24\ \mathrm{AU} \\
    a_{\mathrm{Sen}}(1 + e_{\mathrm{Sen}}) &= R_2 - 1 = 3.67\ \mathrm{AU}
  \end{align*}
  from which one finds $a_{\mathrm{Sen}} \sim 2.46$\,AU \hfill(2 point)

  and $e_{\mathrm{Sen}} \sim 0.49$ \hfill(2 points)

  (figures rounded to 2 dec.\ digits)

\item[(2)] The period, $T_{\mathrm{Sen}}$, of Seneca's orbit can be found by
  using Kepler's 3rd law for Seneca and the Earth (ignoring their small
  masses, compared to the Sun's):
  \[ a_{\mathrm{Sen}}^{3}/T_{\mathrm{Sen}}^{2}
     = a_{\mathrm{Earth}}^{3}/T_{\mathrm{Earth}}^{2} = 1 \]
  \hfill(6 points)

  from which one finds $T_{\mathrm{Sen}} \sim 3.87$\,yr \hfill(2 point)

  (Alternatively one can use Kepler's law for Seneca only and use natural
  units for the mass of the Sun and $G$, i.e.\ $[\mathrm{mass}] =
  1\,M_{\mathrm{Sun}}$, $[t] = 1$\,yr and $[r] = 1$\,AU, in which case
  $T_{\mathrm{Sen}} = a_{\mathrm{Sen}}^{3/2} = 3.86$\,yr) \hfill(8 points)

  Assuming non-retrogate orbit, % sic: "non-retrogate" in source
  the synodic period, $T_{\mathrm{syn}}$ of Seneca and Earth is given by
  \[ 1/T_{\mathrm{syn}} = 1/T_{\mathrm{Earth}} - 1/T_{\mathrm{Sen}}
     \quad (\text{Seneca is superior}) \]
  \hfill(6 points)

  which gives $T_{\mathrm{syn}} = 1.35$\,yr \hfill(2 point)

\item[(3)] From Kepler's law, we can calculate the mass of the Sun
  \[ 4\pi^2\, a_{\mathrm{Sen}}^{3}/T_{\mathrm{Sen}}^{2} = G\,M_{\mathrm{Sun}},
     \quad\text{i.e.}\quad
     M_{\mathrm{Sun}} \sim 1.988\times10^{30}\ \mathrm{kg}
     \quad(\text{depending on accuracy}) \]
  \hfill(4 points)

  Then, using Kepler's 3rd law for Seneca and Jupiter, one gets
  \[ \left(a_{\mathrm{Jup}}^{3}/T_{\mathrm{Jup}}^{2}\right) \big/
     \left(a_{\mathrm{Sen}}^{3}/T_{\mathrm{Sen}}^{2}\right)
     = \left(M_{\mathrm{Sun}} + M_{\mathrm{Jup}}\right)/M_{\mathrm{Sun}}
     = 1 + x \]
  \hfill(4 points)

  where $x$ is the mass ratio of Jupiter to the Sun. Then, solving for $x$ we
  get (depending on the accuracy used in determining the elements of Seneca's
  orbit)
  \[ x = 0.0016, \quad\text{thus}\quad
     M_{\mathrm{Jup}} = 3.2\times10^{27}\ \mathrm{kg} \]
  \hfill(4 points)
\end{description}

\solhead{2.19}{Cosmic Pendulum}
% source id: 2014_TH_L2
\begin{description}
\item[a)] In figure 1 are represented the forces in the system.

  $\vec F_1$ -- the gravitational attraction force acting on the shuttle due
  to the Earth;

  $\vec F_2$ -- the gravitational attraction force acting on the satelite due % sic: "satelite"
  to the Earth;

  $\vec F$ -- the tension force in the suspention rod. % sic: "suspention"

\ioaafig{2014_TH_L2-s1.pdf}

  For a value of $\alpha \neq 0$ the movement equations are:

  For the shuttle on circular orbit around the Earth:
  \[
    m_1\omega^2r = K\frac{m_1M}{r^2} + F\cos\alpha;
  \]
  For the satelite on circular orbit around the Earth:
  \[
    m_2\omega^2\left(r - l\cos\alpha\right)
      = K\frac{m_2M}{\left(r - l\cos\alpha\right)^2} - \frac{F}{\cos\alpha}.
  \]
  If the gravitational interaction between the shuttle and the satellite is
  negligible results: % sic: "is negligible results"
  \begin{align*}
    F\cos\alpha &\ll K\frac{m_1M}{r^2}; \\
    m_1\omega^2r \approx K\frac{m_1M}{r^2};\quad
      \omega^2 &= \frac{KM}{r^3} = \frac{4\pi^2}{T^2}; \\
    T &= 2\pi\sqrt{\frac{r^3}{KM}},
  \end{align*}
  The revolution time of the shuttle around the Earth:
  \begin{align*}
    m_2\frac{KM}{r^3}\left(r - l\cos\alpha\right)
      &= K\frac{m_2M}{\left(r - l\cos\alpha\right)^2}
       - \frac{F}{\cos\alpha}; \\
    m_2\frac{KM}{r^3}\left(r - l\cos\alpha\right)^3
      &= Km_2M - \frac{F}{\cos\alpha}\left(r - l\cos\alpha\right)^2; \\
    m_2\frac{KM}{r^3}r^3\left(1 - \frac{l}{r}\cos\alpha\right)^3
      &= Km_2M - \frac{F}{\cos\alpha}r^2
        \left(1 - \frac{l}{r}\cos\alpha\right)^2; \\
    m_2KM\left(1 - \frac{l}{r}\cos\alpha\right)^3
      &= Km_2M - \frac{F}{\cos\alpha}r^2
        \left(1 - \frac{l}{r}\cos\alpha\right)^2; \\
    m_2KM\left(1 - 3\frac{l}{r}\cos\alpha\right)
      &= Km_2M - \frac{F}{\cos\alpha}r^2
        \left(1 - 2\frac{l}{r}\cos\alpha\right); \\
    3Km_2M\frac{l}{r}\cos\alpha &\approx \frac{F}{\cos\alpha}r^2; \\
    F &\approx 3Km_2M\frac{l}{r^3}\cos^2\alpha,
  \end{align*}
  The tension force in the rod.

  The satellite movement regardless to the shuttle is non uniform circular
  one, described by the equation:
  \[
    m_2\vec a_{\mathrm{t}} = \vec F'';
  \]
  \[
    m_2a_{\mathrm{t}} = m_2\varepsilon l
      = m_2\frac{\mathrm{d}^2\alpha}{\mathrm{d}t^2}l
      = -F'' = -F\tan\alpha
      = -3Km_2M\frac{l}{r^3}\sin\alpha\cos\alpha;
  \]
  \[
    \frac{\mathrm{d}^2\alpha}{\mathrm{d}t^2}
      + 3\frac{KM}{r^3}\sin\alpha\cos\alpha = 0,
  \]
  The solutions of this equation is $\alpha(t)$. % sic: "solutions ... is"

  If during the system evolution the configuration of the system remains the
  same results:
  \begin{gather*}
    \alpha = \text{constant};\quad
      \frac{\mathrm{d}^2\alpha}{\mathrm{d}t^2} = 0; \\
    \sin\alpha\cdot\cos\alpha = 0; \\
    \sin\alpha = 0;\ \alpha_1 = 0;\ \alpha_2 = \pi;\
      \text{(echilibru stabil);} \\ % sic: Romanian "stable equilibrium"
    \cos\alpha = 0;\ \alpha_3 = \pi/2;\ \alpha_4 = 3\pi/2;\
      \text{(echilibru instabil),} % sic: Romanian "unstable equilibrium"
  \end{gather*}
  As seen in the pictures.

\ioaafig{2014_TH_L2-s2.pdf}

\item[b)] Coresponding to the position with $\alpha = 0$, the movement % sic: "Coresponding"
  equation of the satellite displaced from equilibrium position with a very
  small angle $\beta$:
  \begin{gather*}
    \frac{\mathrm{d}^2\beta}{\mathrm{d}t^2}
      + 3\frac{KM}{r^3}\sin\beta\cdot\cos\beta = 0; \\
    \sin\beta \approx \beta;\ \cos\beta \approx 1; \\
    \frac{\mathrm{d}^2\beta}{\mathrm{d}t^2} + 3\frac{KM}{r^3}\beta = 0,
  \end{gather*}
  Which represents the linear harmonic oscillator;
  \[
    \frac{\mathrm{d}^2\beta}{\mathrm{d}t^2} + \omega_0^2\beta = 0;
  \]
  Thus the period of the small oscillations will be
  \[
    \omega_0^2 = 3\frac{KM}{r^3} = \frac{4\pi^2}{T_0^2};\quad
    T_0 = 2\pi\sqrt{\frac{1}{3}\frac{r^3}{KM}} = \frac{T}{\sqrt{3}},
  \]

\item[c)]
  \begin{align*}
    m_1\omega^2r &= K\frac{m_1M}{r^2} + F\cos\alpha; \\
    m_2\omega^2\left(r - l\cos\alpha\right)
      &= K\frac{m_2M}{\left(r - l\cos\alpha\right)^2}
       - \frac{F}{\cos\alpha}; \\
    \omega^2 &= \frac{KM}{r^3} + \frac{F\cos\alpha}{m_1r}; \\
    m_2\omega^2\left(r - l\cos\alpha\right)^3
      &= Km_2M - \frac{F}{\cos\alpha}\left(r - l\cos\alpha\right)^2; \\
    m_2\omega^2r^3\left(1 - \frac{l}{r}\cos\alpha\right)^3
      &= Km_2M - \frac{F}{\cos\alpha}r^2
        \left(1 - \frac{l}{r}\cos\alpha\right)^2; \\
    m_2\omega^2r^3\left(1 - 3\frac{l}{r}\cos\alpha\right)
      &= Km_2M - \frac{F}{\cos\alpha}r^2
        \left(1 - 2\frac{l}{r}\cos\alpha\right); \\
    m_2\left(\frac{KM}{r^3} + \frac{F\cos\alpha}{m_1r}\right)
      r^3\left(1 - 3\frac{l}{r}\cos\alpha\right)
      &= Km_2M - \frac{F}{\cos\alpha}r^2
        \left(1 - 2\frac{l}{r}\cos\alpha\right); \\
    m_2\left(KM + \frac{F\cos\alpha}{m_1}r^2\right)\cdot
      \left(1 - 3\frac{l}{r}\cos\alpha\right)
      &= Km_2M - \frac{F}{\cos\alpha}r^2
        \left(1 - 2\frac{l}{r}\cos\alpha\right);
  \end{align*}
  \[
    F = \frac{3Km_2M\dfrac{l}{r}\cos\alpha}
      {\left(\dfrac{m_2}{m_1}\cos\alpha
        + \dfrac{1}{\cos\alpha}\right)r^2
        - rl\left(\dfrac{m_2}{m_1}\cos^2\alpha + 1\right)};
  \]
  \[
    \omega^2 = \frac{KM}{r^3} + \frac{F\cos\alpha}{m_1r};
  \]
  \[
    \omega^2 = \frac{KM}{r^3}
      + \frac{3K\dfrac{m_2}{m_1}M\dfrac{l}{r^2}\cos^2\alpha}
      {\left(\dfrac{m_2}{m_1}\cos\alpha
        + \dfrac{1}{\cos\alpha}\right)r^2
        - rl\left(\dfrac{m_2}{m_1}\cos^2\alpha + 1\right)}.
  \]
  \textit{Observation}: If $m_2 \ll m_1$ and $l \ll r$, the results are
  already found:
  \[
    \omega^2 = \frac{KM}{r^3} = \omega_0^2;
  \]
  \[
    F = \frac{3Km_2M\dfrac{l}{r}\cos\alpha}
      {\left(\dfrac{1}{\cos\alpha}\right)r^2}
      = 3Km_2M\frac{l}{r^3}\cos^2\alpha.
  \]
  Corresponding to the initial circulae orbit with radius $r$, the total % sic: "circulae"
  mechanical energy of the system shuttle -- Earth is:
  \[
    E_0 = \frac{m_1\mathrm{v}_0^2}{2} - K\frac{m_1M}{r} = -K\frac{m_1M}{2r}.
  \]
  After one complete rotation, the tangential component of the tension in the
  rod $\vec F_{\mathrm{t}}$, acting on the shuttle will determine the change
  of the radius of the orbit i.e.\ $(r - \Delta r)$. The total energy of the
  system will be:
  \begin{gather*}
    E = \frac{m_1\mathrm{v}^2}{2} - K\frac{m_1M}{r - \Delta r}
      = -K\frac{m_1M}{2(r - \Delta r)}. \\
    \Delta E = E - E_0
      = -K\frac{m_1M}{2(r - \Delta r)} + K\frac{m_1M}{2r}
      = -K\frac{m_1M}{2}
        \left(\frac{1}{r - \Delta r} - \frac{1}{r}\right); \\
    \Delta E = -K\frac{m_1M}{2}\cdot\frac{r - r + \Delta r}{r(r - \Delta r)}
      = -K\frac{m_1M}{2}\cdot\frac{\Delta r}{r(r - \Delta r)}; \\
    r - \Delta r \approx r;
  \end{gather*}
  The variation of the mechanical energy after a complete rotation
  \[
    \Delta E \approx -K\frac{m_1M\cdot\Delta r}{2r^2} < 0,
  \]
  Thus
  \[
    \Delta E = L_{\mathrm{t}} = -2\pi r\cdot F_{\mathrm{t}};\quad
    F_{\mathrm{t}} = F\sin\alpha;
  \]
  \[
    -K\frac{m_1M\cdot\Delta r}{2r^2} = -2\pi r\cdot F_{\mathrm{t}}
  \]
  The variation of the altitude due to the action of the satellite will be:
  \[
    \Delta r = \frac{4\pi r^3F_{\mathrm{t}}}{Km_1M}
      = \frac{4\pi r^3F\sin\alpha}{Km_1M},
  \]
  The potentialenergy v % sic: "The potentialenergy v" as printed
  \begin{gather*}
    \Delta E_{\mathrm{p}} = -K\frac{m_1M}{r - \Delta r}
      - \left(-K\frac{m_1M}{r}\right)
      = -Km_1M\left(\frac{1}{r - \Delta r} - \frac{1}{r}\right); \\
    \Delta E_{\mathrm{p}} = -Km_1M\cdot\frac{r - r + \Delta r}{r(r - \Delta r)}
      = -Km_1M\frac{\Delta r}{r(r - \Delta r)}; \\
    r - \Delta r \approx r; \\
    \Delta E_{\mathrm{p}} \approx -K\frac{m_1M\cdot\Delta r}{r^2} < 0,
  \end{gather*}
  Thus
  \begin{gather*}
    \Delta E_{\mathrm{p}} \approx -K\frac{m_1M\cdot\Delta r}{2r^2}\cdot2
      = \Delta E\cdot2 = -2\pi rF_{\mathrm{t}}\cdot2
      = -4\pi rF_{\mathrm{t}}; \\
    \Delta E_{\mathrm{c}} = \frac{m_1\mathrm{v}^2}{2}
      - \frac{m_1\mathrm{v}_0^2}{2}; \\
    \frac{m_1\mathrm{v}_0^2}{r} = K\frac{m_1M}{r^2};\quad
      m_1\mathrm{v}_0^2 = K\frac{m_1M}{r}; \\
    \frac{m_1\mathrm{v}^2}{r - \Delta r}
      = K\frac{m_1M}{(r - \Delta r)^2};\quad
      m_1\mathrm{v}^2 = K\frac{m_1M}{r - \Delta r}; \\
    \Delta E_{\mathrm{c}} = K\frac{m_1M}{2(r - \Delta r)}
      - K\frac{m_1M}{2r}
      = K\frac{m_1M}{2}\left(\frac{1}{r - \Delta r} - \frac{1}{r}\right); \\
    \Delta E_{\mathrm{c}} = K\frac{m_1M}{2}
      \left(\frac{1}{r - \Delta r} - \frac{1}{r}\right)
      = K\frac{m_1M}{2}
      \left(\frac{1}{r - \Delta r} - \frac{1}{r}\right); \\ % sic: identical expression repeated on both sides
    \Delta E_{\mathrm{c}} = K\frac{m_1M}{2}
      \left(\frac{1}{r - \Delta r} - \frac{1}{r}\right)
      = K\frac{m_1M}{2}\cdot\frac{r - r + \Delta r}{r(r - \Delta r)}; \\
    \Delta E_{\mathrm{c}} = K\frac{m_1M}{2}\cdot
      \frac{\Delta r}{r(r - \Delta r)}; \\
    r - \Delta r \approx r; \\
    \Delta E_{\mathrm{c}} \approx K\frac{m_1M\cdot\Delta r}{2r^2} > 0,
  \end{gather*}
  reprezent\^and varia\c{t}ia energiei cinetice a navei, dup\u{a} o rota\c{t}ie
  complet\u{a} \^in jurul P\u{a}m\^antului; % sic: the official solution switches to Romanian here
  \[
    \Delta E_{\mathrm{c}} > 0;\ \mathrm{v} > \mathrm{v}_0,
  \]
  rezultat care constituie ``paradoxul sateli\c{t}ilor'', adic\u{a}, atunci
  c\^and altitudinea orbitei navei scade, viteza navei cre\c{s}te;
  \begin{gather*}
    \Delta E = \Delta E_{\mathrm{c}} + \Delta E_{\mathrm{p}}; \\
    \Delta E_{\mathrm{c}} = \Delta E - \Delta E_{\mathrm{p}}
      = -2\pi rF_{\mathrm{t}} - \left(-4\pi rF\ \right)
      = 2\pi rF_{\mathrm{t}} > 0; \\ % sic: F for F_t inside the bracket
    \Delta E_{\mathrm{c}} = \frac{m_1\mathrm{v}^2}{2}
      - \frac{m_1\mathrm{v}_0^2}{2}
      = \frac{m_1}{2}\left(\mathrm{v}^2 - \mathrm{v}_0^2\right)
      = \frac{m_1}{2}\left(\mathrm{v} - \mathrm{v}_0\right)
        \left(\mathrm{v} + \mathrm{v}_0\right); \\
    \mathrm{v} = \mathrm{v}_0 + \Delta\mathrm{v};\quad
      \mathrm{v} - \mathrm{v}_0 = \Delta\mathrm{v};\quad
      \mathrm{v} + \mathrm{v}_0 = 2\mathrm{v}_0 + \Delta\mathrm{v}
      \approx 2\mathrm{v}_0; \\
    \Delta E_{\mathrm{c}} \approx \frac{m_1}{2}2\mathrm{v}_0\cdot
      \Delta\mathrm{v} = m_1\mathrm{v}_0\Delta\mathrm{v}
      = 2\pi rF_{\mathrm{t}}; \\
    \Delta\mathrm{v} = \frac{2\pi rF_{\mathrm{t}}}{m_1\mathrm{v}_0};\quad
      \mathrm{v}_0 = \sqrt{K\frac{M}{r}}; \\
    \Delta\mathrm{v} = \frac{2\pi rF_{\mathrm{t}}}{m_1}\cdot
      \sqrt{\frac{r}{KM}}.
  \end{gather*}

\item[d)] Manevra tehnic\u{a} propus\u{a} este reprezentat\u{a} \^in figura
  al\u{a}turat\u{a}: o for\c{t}\u{a} reactiv\u{a}, egal\u{a} \^in modul
  \c{s}i de sens contrar cu tensiunea din tij\u{a}: % sic: part d) of the official solution is in Romanian
  \[
    \vec F_{\mathrm{r}} = -\vec F,
  \]
  astfel \^inc\^at for\c{t}a rezultant\u{a} care ac\c{t}ioneaz\u{a} asupra
  navei s\u{a} r\u{a}m\^an\u{a} for\c{t}a de atrac\c{t}ie
  gravita\c{t}ional\u{a} din partea P\u{a}m\^antului:
  \[
    \vec F_{\mathrm{Nava}} = \vec F_1 + \vec F + \vec F_{\mathrm{r}}
      = \vec F_1.
  \]

\ioaafig{2014_TH_L2-s3.pdf}
\end{description}

\solhead{2.20}{From Romania to Antipod! A ballistic messenger}
% source id: 2014_TH_L3
\begin{description}
\item[a)] The two places are represented in the figure.

\ioaafig{2014_TH_L3-s1.pdf}

  \begin{gather*}
    \varphi_{\mathrm{Y}} = \varphi_{\mathrm{Y,Sud}}
      = \varphi_{\mathrm{X,Nord}} = \varphi_{\mathrm{X}}; \\
    \lambda_{\mathrm{Y,Vest}} + \lambda_{\mathrm{X,Eest}} = 180^\circ;\quad % sic: "Eest"
      \lambda_{\mathrm{Y}} + \lambda_{\mathrm{X}} = 180^\circ. \\
    \varphi_{\mathrm{Romania}} = 43^\circ\ \text{Nord};\quad % sic: statement gives 44 deg North
      \lambda_{\mathrm{Romania}} = 30^\circ\ \text{Est},
  \end{gather*}
  The landing point is
  \[
    \varphi_{\mathrm{Antipod}} = 43^\circ\ \text{Sud};\quad
    \lambda_{\mathrm{Antipod}} = 150^\circ\ \text{Vest},
  \]
  Somewhere South EEst from Tasmania (South from Australia). % sic: "EEst"

\ioaafig{2014_TH_L3-s2.pdf}
\ioaafig[1.0]{2014_TH_L3-s3.pdf}

\item[b)] The schetch of the trajectory % sic: "schetch"

\ioaafig{2014_TH_L3-s4.pdf}

  In order to hit the point the trajectory of the missile has to be an elypse % sic: "elypse"
  with the Earth center in the center of the Earth. Se \c{s}tie c\u{a}: % sic: Romanian "it is known that"
  \[
    \mathrm{F_2B} = 2\cdot\mathrm{F_1B};\quad
    \mathrm{F_1F_2} = 3\cdot\mathrm{F_1B}.
  \]
  Rezult\u{a}: % sic: Romanian "it results"
  \begin{gather*}
    \tan2\alpha = \frac{\mathrm{F_1F_2}}{R};\quad
      \mathrm{F_1F_2} = R\cdot\tan2\alpha; \\
    \tan\alpha = \frac{\mathrm{F_1B}}{R};\quad
      \mathrm{F_1B} = R\cdot\tan\alpha; \\
    R\cdot\tan2\alpha = \mathrm{F_1B} = R\cdot\tan\alpha; \\ % sic: missing factor 3 in the middle member
    \tan2\alpha = 3\cdot\tan\alpha; \\
    \frac{\sin2\alpha}{\cos2\alpha} = 3\frac{\sin\alpha}{\cos\alpha}; \\
    \frac{2\sin\alpha\cdot\cos\alpha}{\cos2\alpha}
      = 3\frac{\sin\alpha}{\cos\alpha}; \\
    2\cos^2\alpha = 3\cos2\alpha;\quad
      2\cos^2\alpha = 3\left(\cos^2\alpha - \sin^2\alpha\right); \\
    3\sin^2\alpha = \cos^2\alpha;\quad \tan^2\alpha = \frac{1}{3}; \\
    \tan\alpha = \frac{\sqrt{3}}{3};\quad \alpha = 30^\circ;\quad
      2\alpha = 60^\circ;\quad
      \beta = 90^\circ - 2\alpha = 30^\circ;\quad 2\beta = 60^\circ; \\
    \Delta(\mathrm{RF_2A}) \rightarrow \text{triunghi echilateral;} \\ % sic: Romanian "equilateral triangle"
    \mathrm{RF_2} = \mathrm{AF_2} = \mathrm{RA} = 2R; \\
    \mathrm{RF_2} + \mathrm{RF_1} = 2a = 3R; \\
    a = \frac{3}{2}R;
  \end{gather*}
  \begin{gather*}
    \mathrm{v}_0 = \sqrt{KM\left(\frac{2}{r} - \frac{1}{a}\right)};\quad
      r = R;\ g_0 = K\frac{M}{R^2}; \\
    \mathrm{v}_0 = \sqrt{\frac{KM}{R^2}\cdot R^2
      \left(\frac{2}{R} - \frac{2}{3R}\right)}
      = 2\sqrt{\frac{g_0R}{3}}.
  \end{gather*}

\item[c)]
  \[
    \mathrm{v}_{\mathrm{Antipod}} = \mathrm{v}_0.
  \]

\item[d)]
  \begin{gather*}
    \mathrm{F_1F_2} = R\cdot\tan2\alpha = 2c;\quad
      c = \frac{R}{2}\cdot\tan2\alpha = \frac{R}{2}\cdot\tan60^\circ
      = \frac{\sqrt{3}}{2}R; \\
    b = \sqrt{a^2 - c^2} = \sqrt{\frac{3}{2}}R; \\
    2a = 2r_{\min} + 2c;\quad
      r_{\min} = a - c = \frac{1}{2}\left(3 - \sqrt{3}\right)R; \\
    r_{\max} = 2a - r_{\min} = \frac{1}{2}\left(3 + \sqrt{3}\right)R; \\
    \mathrm{v}_{\min} = \sqrt{KM\left(\frac{2}{r_{\max}}
      - \frac{1}{a}\right)};\quad
      r_{\max} = \frac{1}{2}\left(3 + \sqrt{3}\right)R; \\
    \mathrm{v}_{\min} = \sqrt{\frac{KM}{R^2}\cdot R^2
      \left(\frac{4}{\left(3 + \sqrt{3}\right)R} - \frac{2}{3R}\right)}
      = \sqrt{\frac{2g_0R}{3}\cdot\frac{3 - \sqrt{3}}{3 + \sqrt{3}}}.
  \end{gather*}

\item[e)] Accordin to Kepplers laws: % sic: "Accordin to Kepplers"

\ioaafig{2014_TH_L3-s5.pdf}

  \begin{gather*}
    \Omega = \frac{\mathrm{d}S}{\mathrm{d}t} = \text{constant}; \\
    \frac{S_0}{T} = \frac{2\dfrac{S_0}{4} + 2S_1}{\Delta t};\quad
    \frac{S_0}{T} = \frac{\dfrac{S_0}{2} + 2S_1}{\Delta t};\quad
    \frac{S_0}{T} = \frac{S_0 + 4S_1}{2\cdot\Delta t}; \\
    S_0 = \pi ab;\quad
    S_1 = \frac{ab}{2}\left[\sqrt{1 - \frac{b^2}{a^2}}\cdot\frac{b}{a}
      + \arcsin\sqrt{1 - \frac{b^2}{a^2}}\right]; \\
    \Delta t = \frac{S_0 + 4S_1}{2S_0}\cdot T
      = \left(\frac{1}{2} + 2\frac{S_1}{S_0}\right)\cdot T; \\
    T = 2\pi\sqrt{\frac{a^3}{KM}};\quad
      T = \frac{2\pi}{R}\sqrt{\frac{a^3}{g_0}}; \\
    \frac{2S_1}{S_0} = \frac{1}{\pi}\left(\frac{b}{a}\cdot
      \sqrt{1 - \frac{b^2}{a^2}}
      + \arcsin\sqrt{1 - \frac{b^2}{a^2}}\right); \\
    \sqrt{1 - \frac{b^2}{a^2}} = e;\quad
      \frac{2S_1}{S_0} = \frac{1}{\pi}\left(\frac{b}{a}\cdot e
      + \arcsin e\right); \\
    \Delta t = \left(\frac{1}{2} + \frac{eb}{\pi a}
      + \frac{\arcsin e}{\pi}\right)\cdot T.
  \end{gather*}

\item[f)] The integral luminosity of Sun:
  \[
    L_{\mathrm{S}} = \frac{E_{\mathrm{emis,Soare}}}{t}
      = 3.86\cdot10^{26}\ \mathrm{W},
  \]
  Dac\u{a} For a circumsolar surface $\Sigma$ with radius % sic: "Daca For", mixed Romanian/English
  $r_{\mathrm{PS}}$, see picture bellow the solar radiation enegy passing % sic: "bellow", "enegy"
  through the surface in one second is $L_{\mathrm{S}}$.

\ioaafig{2014_TH_L3-s6.pdf}

  \textit{Density of solar flux}
  \[
    \phi_{\mathrm{Soare},r_{\mathrm{PS}}}
      = \frac{E_{\mathrm{emis,Soare}}}{St}
      = \frac{\dfrac{E_{\mathrm{emis,Soare}}}{t}}{S}
      = \frac{L_{\mathrm{S}}}{S}
      = \frac{L_{\mathrm{S}}}{4\pi r_{\mathrm{PS}}^2}
      = \text{constant}.
  \]
  \[
    F_{\mathrm{incident},FullMoon}
      = \phi_{\mathrm{S}un,r_{\mathrm{PS}}}\cdot\pi R_{\mathrm{L}}^2.
  \]
  Dac\u{a} $\alpha_{\mathrm{L}}$ este albedoul Lunii, rezult\u{a}: % sic: the remainder of f) is largely in Romanian
  \[
    \alpha_{\mathrm{L}} = \frac{F_{\mathrm{reflectat},FullMoon}}
      {F_{\mathrm{incident},FullMoon}},
  \]
  unde $F_{\mathrm{reflectat,Luna\,Plina}}$ -- fluxul energetic al
  radia\c{t}iilor reflectate de Luna Plin\u{a} spre observatorul de pe
  P\u{a}m\^ant;
  \[
    F_{\mathrm{reflectat},FullMoon}
      = \alpha_{\mathrm{L}}\cdot F_{\mathrm{incident},FullMoon}
      = \alpha_{\mathrm{L}}\cdot
        \phi_{\mathrm{Soa,,re},r_{\mathrm{PS}}}\cdot % sic: garbled subscript "Soa,,re" as printed
        \pi R_{\mathrm{L}}^2.
  \]
  \^In consecin\c{t}\u{a}, densitatea fluxului energetic ajuns la observator,
  dup\u{a} reflexia pe suprafa\c{t}a Lunii, este:
  \[
    \phi_{moon,\ \mathrm{observator}}
      = \frac{F_{\mathrm{reflectat},\ FullMoon}}{2\pi r_{\mathrm{PL}}^2}
      = \alpha_{\mathrm{L}}\cdot\phi_{\mathrm{Soare},r_{\mathrm{PS}}}\cdot
        \frac{\pi R_{\mathrm{L}}^2}{2\pi r_{\mathrm{PL}}^2}.
  \]
  Symilarly % sic: "Symilarly"
  \[
    \phi_{\mathrm{proiectil,\ observator}}
      = \frac{F_{\mathrm{reflectat,\ proiectil}}}
        {4\pi r_{\mathrm{D,proiectil}}^2}
      = \alpha_{\mathrm{proiectil}}\cdot
        \phi_{\mathrm{Soare},r_{\mathrm{PS}}}\cdot
        \frac{\pi R_{\mathrm{proiectil}}^2}{4\pi r_{\mathrm{D,proiectil}}^2}.
  \]
  \^In expresia anterioar\u{a} s-a avut \^in vedere faptul c\u{a} densitatea
  fluxului energetic al proiectilului la observator rezult\u{a} din
  distribuirea prin suprafa\c{t}a sferei cu raza $r_{\mathrm{P,proiectil}}$.

  Utiliz\^and formula lui Pogson, vom compara magnitudinea aparent\u{a}
  vizual\u{a} a Lunii Pline cu magnitudinea aparent\u{a} vizual\u{a} a
  proiectilului balistic:
  \begin{gather*}
    \log\frac{\phi_{\mathrm{Luna,observator}}}
      {\phi_{\mathrm{proiectil,observator}}}
      = -0.4\left(m_{\mathrm{Luna\,Plina}}
        - m_{\mathrm{proiectil}}\right); \\
    \log\frac{\phi_{\mathrm{Luna,observator}}}
      {\phi_{\mathrm{proiectil,observator}}}
      = \log\frac{\alpha_{\mathrm{L}}\cdot
        \phi_{\mathrm{Soare},r_{\mathrm{PS}}}\cdot
        \dfrac{\pi R_{\mathrm{L}}^2}{2\pi r_{\mathrm{PL}}^2}}
        {\alpha_{\mathrm{proiectil}}\cdot
        \phi_{\mathrm{Soare},r_{\mathrm{PS}}}\cdot
        \dfrac{\pi R_{\mathrm{proiectil}}^2}
          {4\pi r_{\mathrm{D,proiectil}}^2}}
      = \log\frac{\alpha_{\mathrm{L}}\cdot
        \dfrac{R_{\mathrm{L}}^2}{r_{\mathrm{PL}}^2}}
        {\alpha_{\mathrm{proiectil}}\cdot
        \dfrac{R_{\mathrm{proiectil}}^2}{2r_{\mathrm{D,proiectil}}^2}}; \\
    \log\frac{\alpha_{\mathrm{L}}\cdot
      \dfrac{R_{\mathrm{L}}^2}{r_{\mathrm{PL}}^2}}
      {\alpha_{\mathrm{proiectil}}\cdot
      \dfrac{R_{\mathrm{proiectil}}^2}{2r_{\mathrm{D,proiectil}}^2}}
      = -0.4\left(m_{\mathrm{L}} - m_{\mathrm{proiectil}}\right); \\
    \log\frac{\alpha_{\mathrm{L}}}{\alpha_{\mathrm{proiectil}}}\cdot
      \left(\frac{R_{\mathrm{L}}}{R_{\mathrm{proiectil}}}\right)^2\cdot
      2\cdot\left(\frac{r_{\mathrm{D,proiectil}}}{r_{\mathrm{PL}}}\right)^2
      = -0.4\left(m_{\mathrm{L}} - m_{\mathrm{proiectil}}\right); \\
    \alpha_{\mathrm{L}} = 0.12;\quad \alpha_{\mathrm{proiectil}} = 1; \\
    R_{\mathrm{L}} = 1738\,\mathrm{km};\quad
      R_{\mathrm{proiectil}} = 400\,\mathrm{mm}; \\
    r_{\mathrm{D,proiectil}}
      = r_{\mathrm{max,observator-proiectil}} = h_{\max}
      = r_{\max} - R;\quad
      r_{\max} = \frac{1}{2}\left(3 + \sqrt{3}\right)R; \\
    h_{\max} = \frac{1}{2}\left(3 + \sqrt{3}\right)R - R
      = \frac{1}{2}\left(1 + \sqrt{3}\right)R \approx 8700\ \mathrm{km}; \\
    r_{\mathrm{PL}} = r_{\mathrm{observator,Luna}} = 384400\,\mathrm{km};\quad
      m_{\mathrm{L}} = -12.7^{\mathrm{m}}; \\
    \log\frac{\alpha_{\mathrm{L}}}{\alpha_{\mathrm{proiectil}}}
      + 2\log\frac{R_{\mathrm{L}}}{R_{\mathrm{proiectil}}} + \log2
      + 2\log\frac{r_{\mathrm{D-proiectil}}}{r_{\mathrm{PL}}}
      = -0.4\left(m_{\mathrm{L}} - m_{\mathrm{proiectil}}\right); \\
    \log(0.12) + 2\log\frac{1738000\ \mathrm{m}}{0.400\ \mathrm{m}}
      + \log2 + 2\log\frac{8700\ \mathrm{km}}{384400\ \mathrm{km}}
      = -0.4\left(m_{\mathrm{L}} - m_{\mathrm{proiectil}}\right); \\
    \log(0.12) + 2\log\frac{1738000}{0.400} + \log2
      + 2\log\frac{8700}{384400}
      = -0.4\left(m_{\mathrm{L}} - m_{\mathrm{proiectil}}\right); \\
    -0.920818754 + 13.27597956 + 0.301029995 - 3.290528253
      = -0.4\left(m_{\mathrm{L}} - m_{\mathrm{proiectil}}\right); \\
    23.4^{\mathrm{m}} = 12.7^{\mathrm{m}} + m_{\mathrm{proiectil}}; \\
    m_{\mathrm{proiectil}} = 10.7^{\mathrm{m}}; \\
    m_{\max} \approx 6^{\mathrm{m}};\quad
      m_{\mathrm{proiectil}} > m_{\max},
  \end{gather*}
  The projectile wasn't seen when it was at its apogee
\end{description}

\solhead{2.21}{Lagrange Points}
% source id: 2014_TH_P1
According to the notations in fig.\,1.1 and fig.\,2.1 % sic: the figures are actually labelled "Figure 1a" and "Figure 1b" in the source
\ioaafig{2014_TH_P1-s1.pdf}
\ioaafig{2014_TH_P1-s2.pdf}

\begin{align*}
  \vec{F}_{\mathrm{S}} + \vec{F}_{\mathrm{P}}
    &= \vec{F}_{\mathrm{cp}} = m\vec{a}_{\mathrm{cp}} \\
  F_{\mathrm{S}} + F_{\mathrm{P}}
    &= ma_{\mathrm{cp}} = m\omega^2\left(r_{\mathrm{PS}} \pm w\right);
\end{align*}
\hfill(2 points)

The sign ``$+$'' for position $\mathrm{L}_3'$ and ``$-$'' for $\mathrm{L}_3''$
\[
  K\frac{mM_{\mathrm{S}}}{\left(r_{\mathrm{PS}} \pm w\right)^2}
  + K\frac{mM_{\mathrm{P}}}{\left(2r_{\mathrm{PS}} \pm w\right)^2}
  = m\omega^2\left(r_{\mathrm{PS}} \pm w\right);
\]
Using the assumption that $w \ll r_{\mathrm{PS}}$
\[
  \left(1 \pm \frac{w}{r_{\mathrm{PS}}}\right)^{-2}
  \approx 1 \mp 2\frac{w}{r_{\mathrm{PS}}}, \qquad
  \left(1 \pm \frac{w}{2r_{\mathrm{PS}}}\right)^{-2}
  \approx 1 \mp \frac{w}{r_{\mathrm{PS}}};
\]
\hfill(2 points)

The rotation speed
\[ \omega^2 = \frac{KM_{\mathrm{S}}}{r_{\mathrm{PS}}^3} \]
\hfill(2 p)

The final relation
\[ w = \mp \frac{M_{\mathrm{P}}\, r_{PS}}{\left(12M_{\mathrm{S}} + M_{P}\right)} \]

The value has to be positive, thus the $\mathrm{L}_3''$ is the position of one
Lagrange point \hfill(2 p)

\[ w \approx 37.54\ \mathrm{Km} \]
\hfill(2 p)

\solhead{2.22}{Sun gravitational catastrophe!}
% source id: 2014_TH_P2
\begin{description}
\item[(A)] The orbit of Earth is elliptical, so the shape of the orbit after
  the solar catastrophe will depend on the moment when the decrees % sic: "decrees" (decrease) in source
  of the mass of the Sun will occur.

  Initial analysis of the problem:
  \begin{description}
  \item[a)] In 3rd July the Earth is at the aphelion. The speed of the Earth
    is smaller than the speed of Earth on a circular orbit with radius
    $r_{0,\mathrm{max}} = a_0(1 + e_0)$.
  \item[b)] In 3rd January the Earth is at perihelion. The speed of the Earth
    is bigger than the speed of Earth on a circular orbit with radius
    $r_{0,\mathrm{max}} = a_0(1 - e_0)$. % sic: source writes r_{0,max} here; should be r_{0,min}
  \end{description}

  Conclusion (A): the period should be calculated only for situation a). The
  expected trajectory in this case is an elliptic one. The possibility that
  Earth hit the Sun is available too.

\item[(B)] Calculations:

  In 3rd July the distance from Sun is maximum: fig.~2.1,
  \[ r_{0,\mathrm{max}} = a_0(1 + e_0). \]
\ioaafig{2014_TH_P2-s1.pdf}

  \textbf{Before the catastrophe:}
  \begin{itemize}
  \item $v_{0,\mathrm{aph}}$ -- the speed of Earth on aphelion,
  \item $a_0$ -- big Earth's elliptical orbit semi axis,
  \item $v_0$ -- the speed of Earth if its orbit is circular with radius
    $r_0 = a_0$.
  \end{itemize}

  According to Keppler's % sic: "Keppler's" in source
  second law and the law of energy conservation (see figure 2.1) the
  following relations can be written:
  \begin{align*}
    & v_{0,\mathrm{per}}\, r_{0,\mathrm{per}}
      = v_{0,\mathrm{aph}}\, r_{0,\mathrm{aph}}; \\
    & \frac{v_{0,\mathrm{per}}^2}{2} - K\frac{M_0}{r_{0,\mathrm{per}}}
      = \frac{v_{0,\mathrm{aph}}^2}{2} - K\frac{M_0}{r_{0,\mathrm{aph}}} \\
    & r_{0,\mathrm{min}} = r_{0,\mathrm{per}} = a_0(1 - e_0) \\
    & r_{0,\mathrm{max}} = r_{0,\mathrm{aph}} = a_0(1 + e_0) \\
    & KM_0 = v_0^2\, r_0 = v_0^2\, a_0 \\
    & v_0 = \sqrt{K\frac{M_0}{r_0}} = \sqrt{K\frac{M_0}{a_0}} \\
    & v_{0,\mathrm{per}} = v_0 \sqrt{\frac{1 + e_0}{1 - e_0}} > v_0 \\
    & v_{0,\mathrm{per}} > v_0 \tag*{(1)} \\
    & v_{0,\mathrm{aph}} = v_0 \sqrt{\frac{1 - e_0}{1 + e_0}} \\
    & v_{0,\mathrm{aph}} < v_0 \tag*{(2)}
  \end{align*}

  Conclusion -- According to the relations (1) and (2) the new orbit of the
  Earth could be an elliptic one.

  For the new elliptical Earth orbit:
  \begin{align*}
    & r_{\mathrm{per}} = r_{0,\mathrm{aph}}; \\
    & r_{\mathrm{min}} = r_{\mathrm{per}} = a(1 - e); \\
    & a_0(1 + e_0) = a(1 - e); \quad a = a_0\,\frac{1 + e_0}{1 - e}; \\
    & v_{\mathrm{per}} = v_{0,\mathrm{aph}}; \\
    & v_{\mathrm{per}} = v\sqrt{\frac{1 + e}{1 - e}},
  \end{align*}
  where $v$ este % sic: Romanian "este" (is) left in source
  the Earth's speed on a circular orbit with the radius $r = a$, when the
  mass of the Sun becomes $M = M_0/2$;
  \begin{align*}
    & v\sqrt{\frac{1 + e}{1 - e}} = v_0\sqrt{\frac{1 - e_0}{1 + e_0}}; \\
    & v = \sqrt{K\frac{M}{r}} = \sqrt{K\frac{M_0}{2a}}
      = \sqrt{K\frac{M_0}{a_0}}\sqrt{\frac{a_0}{2a}}
      = v_0\sqrt{\frac{a_0}{2a}}; \\
    & e = 1 - 2e_0; \quad a = a_0\,\frac{1 + e_0}{2e_0}.
  \end{align*}

  Conclusion:
  \begin{align*}
    T_0 &= \frac{2\pi r_0}{v_0} = \frac{2\pi a_0}{v_0} \\
    T &= \frac{2\pi r}{v} = \frac{2\pi a}{v}; \\
    \frac{T}{T_0} &= \frac{a}{a_0}\,\frac{v_0}{v}
      = \frac{1 + e_0}{2e_0}\sqrt{\frac{2a}{a_0}}
      = \frac{1 + e_0}{2e_0}\,\sqrt{2}\,\sqrt{\frac{1 + e_0}{2e_0}}; \\
    T &= T_0\sqrt{2}\left(\frac{1 + e_0}{2e_0}\right)^{3/2}
      \approx 230\ \text{years}
  \end{align*}

  b) In 3rd of January the Earth is at perihelion. In that moment the Erath % sic: "Erath" in source
  speed is larger than the speed necessary for an Earth's circular orbit.
  Thus the trajectory of the Earth after the catastrophe will be an open
  trajectory, i.e.\ an hyperbolic or parabolic orbit.

  Conclusion: it is not necessary to calculate the period of revolution or it
  could be issued as infinite.
\end{description}

\solhead{2.23}{The Astronaut saved by \ldots\ ice from a can!}
% source id: 2014_TH_P5
\textbf{Theoretical considerations:}

The incident particle flux on a wall (i.e.\ a certain direction on a surface)
is:
\[ \Phi = m_0 \cdot \Omega
   = \frac{1}{6}\cdot m_0 \cdot n \cdot S \cdot \bar{v} \tag*{(3)} \]
\[ m_0 \cdot n = m_0 \cdot \frac{N}{a^3} = \frac{m_0 N}{a^3}
   = \frac{m}{V} = \rho, \]
where: $m_0$ is the mass of one molecule; $m$ -- mass of the gas in the cube;
$V$ -- volume of the cube; $\rho$ -- the density of the gas;
\[ \rho = \frac{\mu p}{RT}, \]
where $p$ -- the pressure of the gas in the cube. The relation (3) -- the mass
flux relation becomes:
\[ \Phi = \frac{1}{6}\cdot\rho\cdot S\cdot\bar{v}
   = \frac{1}{6}\cdot\frac{\mu p}{RT}\cdot S\cdot\sqrt{\frac{3RT}{\mu}}
   = \frac{1}{6}\cdot p\cdot S\cdot\sqrt{\frac{3\mu}{RT}}. \tag*{(4)} \]
\[ \Phi = \frac{1}{6}\cdot\rho\cdot S\cdot\bar{v}
   = \frac{1}{6}\cdot\frac{\mu p}{RT}\cdot S\cdot\sqrt{\frac{3RT}{\mu}}
   = \frac{1}{6}\cdot p\cdot S\cdot\sqrt{\frac{3\mu}{RT}}. \tag*{(5)} \]
% sic: relation (5) is a verbatim repeat of relation (4) in the source
\hfill(3 points)

\textbf{B. Reasoning}

Because the cylindrical can is not full of ice, in the empty part of it there
are saturated vapors, i.e.\ the mass flux of the molecule which sublimate is
equal with mass flux of gass % sic: "gass" in source
which transform into ice. Thus the pressure in the can is the saturated vapor
pressure $p_{\mathrm{s}}$ and it has the corresponding maximum density
$\rho_{\mathrm{s}}$. See figure 6.2.
\ioaafig{2014_TH_P5-s1.pdf}

\begin{align*}
  \Phi_1 &= \Phi_{\mathrm{sublimation}}
    = \frac{1}{6}\cdot\rho_{\mathrm{s}}\cdot S\cdot\bar{v}
    = \frac{1}{6}\cdot p_{\mathrm{s}}\cdot S\cdot\sqrt{\frac{3\mu}{RT}}; \\
  \Phi_2 &= \Phi_{\mathrm{solidification}}
    = \frac{1}{6}\cdot\rho_{\mathrm{s}}\cdot S\cdot\bar{v}
    = \frac{1}{6}\cdot p_{\mathrm{s}}\cdot S\cdot\sqrt{\frac{3\mu}{RT}}; \\
  \Phi_1 &= \Phi_2.
\end{align*}

After the can was opened, there be no molecules which sublimate thus the mass
flux of the molecules which gather the ice become null. So the pressure
becomes $\left(p_{\mathrm{s}}/2\right)$.

Thus the force acting on the astronaut will be

\textbf{C. Calculations according to the reasoning}, and/or as support for
reasoning \hfill(2 points)
\[ F = \frac{p_{\mathrm{s}}}{2}\cdot S, \]
Opening the can the astronaut will be accelearated % sic: "accelearated" in source
with:
\[ a = \frac{F}{M} = \frac{p_{\mathrm{s}}\cdot S}{2M}
   = \frac{550\,\mathrm{N\,m^{-2}} \cdot 30\cdot10^{-4}\,\mathrm{m^2}}
          {2\cdot10^{2}\,\mathrm{kg}}
   = 0.00825\,\mathrm{m\,s^{-2}}. \]

The total time of the acceleration movement will be the total time of ice
sublimation:
\[ \tau = \frac{m}{\Phi_1}
   = \frac{m}{\dfrac{1}{6}\cdot p_{\mathrm{s}}\cdot S
     \cdot\sqrt{\dfrac{3\mu}{RT}}}
   = \frac{6m}{p_{\mathrm{s}}\cdot S}\cdot\sqrt{\frac{RT}{3\mu}}
   \approx 150\,\mathrm{s}. \]

\textbf{D. Correct result}

The travel distance in this time will be:
\[ L = \frac{a\tau^2}{2}
   = \frac{0.00825\,\mathrm{m\,s^{-2}} \cdot 225\cdot10^{2}\,\mathrm{s^2}}{2}
   \approx 93\,\mathrm{m}, \]
The astronaut could arrive safely in at the shuttle if he didn't lose to much
time by solving the problem. % sic: "in at the shuttle", "to much" in source

\solhead{2.24}{}
% source id: 2015_TH_L1
\subsubsection*{1. Polar version}

\begin{description}
\item[(a)] \textbf{Rewrite the orbit equation using polar coordinate}
  \begin{align*}
    x &= R(\alpha)\cos\alpha \\
    y &= R(\alpha)\sin\alpha
  \end{align*}
  Thus, if we let
  \begin{align*}
    x^* &= R(\alpha)\cos\left(\alpha - \frac{\pi}{3}\right)
      = \frac{x}{2} + \frac{y\sqrt{3}}{2} \\
    y^* &= R(\alpha)\sin\left(\alpha - \frac{\pi}{3}\right)
      = -\frac{x\sqrt{3}}{2} + \frac{y}{2}
  \end{align*}
  \hfill(10 points)

  the equation of the orbit can be written as
  \begin{align*}
    9\left(R(\alpha)\cos\left(\alpha - \frac{\pi}{3}\right) - 4\right)^2
      + 25\left(R(\alpha)\sin\left(\alpha - \frac{\pi}{3}\right)\right)^2
      &= 225 \\
    9(x^* - 4)^2 + 25(y^*)^2 &= 225
  \end{align*}
  \hfill(10 points)

  \textbf{Determine semi major and semi minor axis of the orbit}

  Since
  \[
    9(x^* - 4)^2 + 25(y^*)^2 = 225
  \]
  is equivalent to
  \[
    \frac{(x^* - 4)^2}{25} + \frac{(y^*)^2}{9} = 1
  \]
  then, the semi major axis is 5, and \hfill(5 points) \\
  the semi minor axis is 3 \hfill(5 points)

\item[(b)] \textbf{Determine the zenith angle at perigee} \hfill(20 points)

  At the perigee, $\alpha - \frac{\pi}{3} = \pi$. Thus, the zenith angle at
  the perigee is
  \[
    \alpha - \frac{\pi}{2} = \pi + \frac{\pi}{3} - \frac{\pi}{2}
      = \frac{5\pi}{6}
  \]

\ioaafig{2015_TH_L1-s1.pdf}

\item[(c)] \textbf{Characterize the coordinates of the moon when it looks
  largest to the observer.}

  Let $P(x_0, y_0)$ be the point where it looks largest to the observer. Since
  the distance between the moon and the observer decreases as it moves from
  apogee towards perigee, then the moon looks largest to the observer at the
  planet when it closest to the planet while the whole of the moon still can % sic: "when it closest"
  be seen in full. Hence, $y_0 = r$.
  \begin{align*}
    y_0 &= r = R(\pi - \alpha_0\alpha_0)\sin(\pi - \alpha_0)
      = R_0\sin(\pi - \alpha_0) \\ % sic: "R(pi - alpha_0 alpha_0)" as printed
    x_0 &= R_0\cos(\pi - \alpha_0)
  \end{align*}
  \hfill(10 points)

\ioaafig{2015_TH_L1-s2.pdf}

  \textbf{Determine the coordinates of the point.}

  Let the coordinate be $(x_0, r)$.
  \[
    9\left(\frac{x_0}{2} + \frac{\sqrt{3}r}{2} - 4\right)^2
      + 25\left(-\frac{\sqrt{3}x_0}{2} + \frac{r}{2}\right)^2 = 225
  \]
  Rewrite it as a quadratic equation in $x_0$.
  \begin{align*}
    9\left(x_0 + r\sqrt{3} - 8\right)^2
      + 25\left(-x_0\sqrt{3} + r\right)^2 &= 900 \\
    84(x_0)^2 + \left(-32\sqrt{3}r - 144\right)x_0
      + \left(52r^2 - 144\sqrt{3}r + 576 = 900\ \right) % sic: stray "= 900" inside the bracket as printed
  \end{align*}
  \hfill(20 points)

  Solving for $x$, we get
  \[
    x_0 = \frac{4r\sqrt{3} + 18}{21}
      \pm \frac{15\sqrt{-r^2 + 4r\sqrt{3} + 9}}{21}
  \]
  Then choose the smaller one
  \[
    x_0 = \frac{4r\sqrt{3} + 18}{21}
      - \frac{15\sqrt{-r^2 + 4r\sqrt{3} + 9}}{21}
  \]
  \hfill(10 points)

  \textbf{Find $\tan\frac{\theta}{2}$.} Hence
  \[
    \tan\frac{\theta_0}{2} = \tan\alpha_0 = \frac{r}{x_0}
      = \frac{21r}{4r\sqrt{3} + 18 - 15\sqrt{-r^2 + 4r\sqrt{3} + 9}}
  \]
  \hfill(10 points)
\end{description}

\subsubsection*{2. Cartesian version}

\begin{description}
\item[(a)] \textbf{Notice the standard version of the orbits}

  Simplify the equation by introducing new variables
  \begin{align*}
    x^* &= \frac{x}{2} + \frac{\sqrt{3}y}{2}
      = x\cos\frac{\pi}{3} + y\sin\frac{\pi}{3} \\
    y^* &= -\frac{\sqrt{3}x}{2} + \frac{y}{2}
      = x\sin\left(-\frac{\pi}{3}\right) + y\cos\frac{\pi}{3}
  \end{align*}
  \hfill(10 points)

  Write the orbit equation in terms new variables % sic: "in terms new variables"
  \[
    9(x^* - 4)^2 + 25(y^*) = 225 % sic: exponent missing on (y*) as printed
  \]
  \hfill(10 points)

\ioaafig{2015_TH_L1-s3.pdf}

  \textbf{Determine semi major and semi minor axis of the orbit}

  Since
  \[
    9(x^* - 4)^2 + 25(y^*)^2 = 225
  \]
  is equivalent to
  \[
    \frac{(x^* - 4)^2}{25} + \frac{(y^*)^2}{9} = % sic: right-hand side left blank as printed
  \]
  The semi major axis is 5 \hfill(5 points) \\
  The semi major axis is 3 \hfill(5 points) % sic: "semi major" twice as printed

\item[(b)] \textbf{Determine the zenith angle at perigee} \hfill(20 points)

  Since the original ellipse may be obtained from standard ellipse
  \[
    \frac{(x^* - 4)^2}{25} + \frac{(y^*)^2}{9} = 1
  \]
  by $\pi/3$ counterclockwise rotation respect to the origin, the the zenith % sic: "the the"
  angle at perigee is equal to
  \[
    \pi + \frac{\pi}{3} - \frac{\pi}{2} = \frac{5\pi}{6}
  \]

\item[(c)] \textbf{Characterize the coordinates of the moon when it looks
  largest to the observer.}

  Let $P(x_0, y_0)$ be the point where it looks largest to the observer. Since
  the distance between the moon and the observer decreases as it moves from
  apogee towards perigee, then the moon looks largest to the observer at the
  planet when it closest to the planet while the whole of the moon still can
  be seen in full. Hence, $y_0 = r$. \hfill(10 points)

  \textbf{Find $x_0$.}
  \[
    9\left(\frac{x_0}{2} + \frac{\sqrt{3}r}{2} - 4\right)^2
      + 25\left(-\frac{\sqrt{3}x_0}{2} + \frac{r}{2}\right)^2 = 225
  \]
  Rewrite it as a quadratic equation in $x$.
  \begin{align*}
    9\left(x_0 + r\sqrt{3} - 8\right)^2
      + 25\left(-x_0\sqrt{3} + r\right)^2 &= 900 \\
    \left(9x_0^2 + \left(18\sqrt{3}r - 144\right)x_0 + 27r^2
      - 144\sqrt{3}r + 576\right)
      + \left(75x_0^2 - 50\sqrt{3}rx_0 + 25r^2\right) &= 900 \\
    84(x_0)^2 + \left(-32\sqrt{3}r - 144\right)x_0
      + \left(52r^2 - 144\sqrt{3}r + 576\ \right) &= 900
  \end{align*}
  Solving it for $x$ we get
  \[
    x_0 = \frac{4r\sqrt{3} + 18}{21}
      \pm \frac{15\sqrt{-r^2 + 4r\sqrt{3} + 9}}{21}
  \]
  \hfill(20 points)

  Then choose the smaller one
  \[
    x_0 = \frac{4r\sqrt{3} + 18}{21}
      - \frac{15\sqrt{-r^2 + 4r\sqrt{3} + 9}}{21}
  \]
  \hfill(10 points)

  \textbf{Find $\tan\frac{\theta}{2}$.} Hence
  \[
    \tan\frac{\theta}{2} = \tan\alpha_0 = \frac{r}{x_0}
      = \frac{21r}{4r\sqrt{3} + 18 - 15\sqrt{-r^2 + 4r\sqrt{3} + 9}}
  \]
  \hfill(10 points)
\end{description}

\subsubsection*{3. Cartesian version} % sic: both the second and third versions are titled "Cartesian version"

\begin{description}
\item[(a)] \textbf{Notice the standard version of the orbits}

  Simplify the equation by introducing new variables
  \begin{align*}
    x^* &= \frac{x}{2} + \frac{\sqrt{3}y}{2}
      = x\cos\frac{\pi}{3} + y\sin\frac{\pi}{3} \\
    y^* &= -\frac{\sqrt{3}x}{2} + \frac{y}{2}
      = x\sin\left(-\frac{\pi}{3}\right) + y\cos\frac{\pi}{3}
  \end{align*}
  \hfill(10 points)

  Write the orbit equation in terms new variables
  \[
    9(x^* - 4)^2 + 25(y^*) = 225
  \]
  \hfill(10 points)

  \textbf{Determine semi major and semi minor axis of the orbit}

  Since
  \[
    9(x^* - 4)^2 + 25(y^*)^2 = 225
  \]
  is equivalent to
  \[
    \frac{(x^* - 4)^2}{25} + \frac{(y^*)^2}{9} =
  \]
  The semi major axis is 5 \hfill(5 points) \\
  The semi major axis is 3 \hfill(5 points) % sic: "semi major" twice as printed

\item[(b)] \textbf{Notice the standard version of the orbits} \hfill(20 points)
  % sic: heading repeated from (a) as printed

  Since the original ellipse may be obtained from standard ellipse
  \[
    \frac{(x^* - 4)^2}{25} + \frac{(y^*)^2}{9} = 1
  \]
  by counterclockwise $\pi/3$ rotation respect to the origin, the zenith angle
  at perigee is
  \[
    \pi + \frac{\pi}{3} - \frac{\pi}{2} = \frac{5\pi}{6}
  \]

\item[(c)] \textbf{Identify and characterized the coordinates when the planet
  looks largest.} % sic: "characterized"

  Since the distance between the moon and the observer decreases as it moves
  from apogee towards perigee, then the moon looks largest to the observer at
  the planet when it closest to the planet while the whole of the moon still
  can be seen in full. Hence, it is at necessary that $y = r$. Let the % sic: "it is at necessary"
  coordinates be $(x_0, r)$. \hfill(10 points)

\ioaafig{2015_TH_L1-s4.pdf}

  \textbf{Obtain expression for calculating distance from a point on an
  ellipse to the foci to set up an equation}

  Consider standard ellipse
  \[
    \frac{(x^* - 4)^2}{5^2} + \frac{(y^*)^2}{3^2} = 1
  \]
  Choose any point $(x^*, y^*)$ on the ellipse. Let $d_1$ and $d_2$ be the
  distances from any point $(x^*, y^*)$ on the ellipse to the foci $(0.0)$
  and $d(8.0)$; $d_1^2 = (x^*)^2 + (y^*)^2$ and
  $d_2^2 = (x^* - 8)^2 + (y^*)^2$

  Thus,
  \[
    d_1^2 = (x^*)^2 + (y^*)^2
      = (x^*)^2 + \left(9 - \frac{9(x^* - 4)^2}{5^2}\right)
      = \frac{(16(x^*)^2 + 72x^* + 81)}{25}
  \]
  Therefore
  \[
    d_1^2 = \frac{(16(x^*)^2 + 72x^* + 81)}{25} = x^2 + r^2
  \]
  \hfill(10 points)

  where $x^* = \frac{x}{2} + \frac{r\sqrt{3}}{2}$. Therefore, we can obtain
  the value of $x$ by solving the above equation for $x$.

  \textbf{Solve equation
  $\frac{(16(x^*)^2 + 72x^* + 81)}{25} = x^2 + r^2$ to obtain value of $x$.}

  Substituting $x^* = \frac{x}{2} + \frac{r\sqrt{3}}{2}$ to the equation, we
  have
  \[
    25d_1^2 = 25x^2 + 25r^2
      = 16\left(\frac{x}{2} + \frac{r\sqrt{3}}{2}\right)^2
      + 72\left(\frac{x}{2} + \frac{r\sqrt{3}}{2}\right) + 81 = % sic: trailing "=" as printed
  \]
  or
  \begin{align*}
    25x^2 + 25r^2 &= 4x^2 + \left(8r\sqrt{3} + 36\right)x
      + \left(12r^2 + 36\sqrt{3}r + 81\right) \\
    21x^2 - \left(8r\sqrt{3} + 36\right)x
      + \left(13r^2 - 36r\sqrt{3} - 81\right) &= 0
  \end{align*}
  Then use quadratic formula to obtain $x$ in term of $r$. Choose smaller
  value of $x$ to get larger value of
  $\tan\frac{\theta}{2} = \frac{r}{x}$. The smaller one is
  \[
    x = \frac{4r\sqrt{3} + 18}{21}
      - \frac{15\sqrt{-r^2 + 4r\sqrt{3} + 9}}{21}
  \]
  \hfill(20 points)

  \textbf{Find the expression of $\tan\frac{\theta}{2}$.}

  Hence
  \[
    \tan\frac{\theta}{2} = \frac{r}{x}
      = \frac{21r}{4r\sqrt{3} + 18 - 15\sqrt{-r^2 + 4r\sqrt{3} + 9}}
  \]
  \hfill(10 points)
\end{description}

\solhead{2.25}{}
% source id: 2015_TH_S1
\textbf{Determine planet periods.}
Synodic period of planet 1 and 2 (assume $P_1 < P_2$) is
\[ \frac{1}{P_S} = \frac{1}{P_1} - \frac{1}{P_2}
   \quad\Rightarrow\quad
   P_S = \frac{P_1 P_2}{P_2 - P_1} \]
Resonant orbit happens if
\[ P_S = mP_1
   \quad\Rightarrow\quad
   \frac{P_1 P_2}{P_2 - P_1} = mP_1
   \quad\Rightarrow\quad
   m = \frac{P_2}{P_2 - P_1} > 1,\ \text{integer} \]

Using the data of the Gliesean system we have
\hfill(50 points)

\begin{center}
\begin{tabular}{|l|c|c|c|}
\hline
\textbf{Gliese System} & \textbf{Mass (kg)} & \textbf{Semi Major Axis (m)}
  & \textbf{Period (Earth, days)} $P = 2\pi\sqrt{\dfrac{a^3}{GM_G}}$ \\
\hline
Gliese 876 star & $6.64359 \times 10^{29}$ & & \\
\hline
Gliese 876 b & $4.31938 \times 10^{27}$ & $3.1161 \times 10^{10}$ & 60.07568 \\
\hline
Gliese 876 c & $1.35564 \times 10^{27}$ & $1.9388 \times 10^{10}$ & 29.48304 \\
\hline
Gliese 876 d & $4.079 \times 10^{25}$ & $3.1116 \times 10^{9}$ & 1.8956 \\
\hline
Gliese 876 e & $8.7194 \times 10^{25}$ & $5.0011 \times 10^{10}$ & 122.1432 \\
\hline
\end{tabular}
\end{center}

\textbf{Calculate resonant orbit for each planet.}
From the results above we see that
\begin{itemize}
\item Gliese 876 c and Gliese 876 b have resonant orbit
  \[ m = \frac{P_b}{P_b - P_c} = 1.96 \approx 2 \]
  \hfill(25 points)
\item Gliese 876 b and Gliese 876 e have resonant orbit
  \[ m = \frac{P_e}{P_e - P_b} = 1.9679 \approx 2 \]
  \hfill(25 points)
\end{itemize}

\solhead{2.26}{}
% source id: 2015_TH_S11
\[
  \frac{1}{2}mv_0^2 - \frac{GmM}{R} = 0
  \rightarrow R = \frac{2GM}{v_0^2}
\]
\hfill(20 points)

\[
  \Delta m(r) = \rho4\pi r^2\Delta r = 4\pi\alpha r\Delta r
  \rightarrow M = 2\pi\alpha R^2
\]
\[
  \left(dm(r) = \rho4\pi r^2dr = 4\pi\alpha r\,dr
  \rightarrow M = 2\pi\alpha R^2\right)
\]
\hfill(30 points)

\[
  M = \alpha2\pi R^2 = \alpha2\pi\frac{4G^2M^2}{v_0^4}
  \rightarrow M = \frac{v_0^4}{8\pi\alpha G^2}
\]
\hfill(30 points)

\[
  M = \frac{v_0^4}{8\pi\alpha G^2}
    = \frac{(1.5\times10^4\ \mathrm{m\,s^{-1}})^4}
      {8\pi(5\times10^{13}\ \mathrm{kg\,m^{-2}})
       (6.67\times10^{-11}\ \mathrm{kg^{-1}m^3s^{-2}})^2}
    = 9.1\times10^{21}\,\mathrm{kg}
\]
\hfill(20 points)

\solhead{2.27}{}
% source id: 2015_TH_S13
\textbf{Calculate the tidal force}

The tidal force (assuming the point mass of hemisphere is located at
$R_{Io}/2$):
\begin{align*}
  F_{\mathrm{tidal}} &= F(r - R_{Io}) - F(r + R_{Io}) \\ % sic: arguments written with R_Io, the fractions below use R_Io/2
  &= -\frac{GM_J\frac{m_{Io}}{2}}{\left(r - \frac{R_{Io}}{2}\right)^2}
     + \frac{GM_J\frac{m_{Io}}{2}}{\left(r + \frac{R_{Io}}{2}\right)^2}
   = -GM_J\frac{m_{Io}}{2}
     \left[\frac{1}{\left(r - \frac{R_{Io}}{2}\right)^2}
     - \frac{1}{\left(r + \frac{R_{Io}}{2}\right)^2}\right] \\
  &= -\frac{GM_J\frac{m_{Io}}{2}}{r^2}
     \left[\frac{\left(\frac{2R_{Io}}{r}\right)}
     {\left[1 - \left(\frac{R_{Io}}{r}\right)^2\right]^2}\right] % sic: exact denominator would be (1-(R_Io/2r)^2)^2
\end{align*}
\hfill(40 points)

The term $\left(\dfrac{R_{Io}}{r}\right)^2$ is too small and it can be
ignored, so that
\[
  F_{\mathrm{tidal}} = -\frac{GM_J m_{Io} R_{Io}}{r^3}
\]
\hfill(10 points)

\medskip
\textbf{Calculate the difference of the tidal force}

The difference in the tidal force experienced at perijove (closest to Jupiter)
and apojove (furthest from Jupiter)
\[
  \Delta F = \left|F_{\mathrm{tidal}}(r_{\mathrm{peri}})
    - F_{\mathrm{tidal}}(r_{\mathrm{apo}})\right|
\]
\hfill(15 points)
\[
  = GM_J m_{Io} R_{Io}
    \left[\frac{1}{r_{\mathrm{peri}}^3} - \frac{1}{r_{\mathrm{apo}}^3}\right]
  = 6.65\times10^{18}\ \mathrm{N}
\]
\hfill(15 points)

\medskip
\textbf{Calculate the average power of work done}

Hence, the average power of the work done on the rock during one-half orbit
is
\[
  \bar P = \frac{\Delta F\,d}{t_{Io}/2}
\]
\hfill(10 points)
\[
  \bar P = 4.35\times10^{15}\,W
\]
\hfill(10 points)

\bigskip
\textbf{Alternate solution}

If assume the mass of the hemisphere is located at the center of mass of a % sic: "If assume"
solid hemisphere, $3R_{Io}/8$, then

\medskip
\textbf{Calculate the tidal force}

The tidal force (assuming the point mass of hemisphere is located at
$3R_{Io}/8$):
\begin{align*}
  F_{\mathrm{tidal}} &= F(r - R_{Io}) - F(r + R_{Io}) \\
  &= -\frac{GM_J\frac{m_{Io}}{2}}{\left(r - \frac{3R_{Io}}{8}\right)^2}
     + \frac{GM_J\frac{m_{Io}}{2}}{\left(r + \frac{3R_{Io}}{8}\right)^2}
   = -GM_J\frac{m_{Io}}{2}
     \left[\frac{1}{\left(r - \frac{3R_{Io}}{8}\right)^2}
     - \frac{1}{\left(r + \frac{3R_{Io}}{8}\right)^2}\right] \\
  &= -\frac{GM_J\frac{m_{Io}}{2}}{r^2}
     \left[\frac{\left(\frac{3R_{Io}}{2r}\right)}
     {\left[1 - \left(\frac{6R_{Io}}{8r}\right)^2\right]^2}\right] % sic: 6R_Io/8r in source; exact would be 3R_Io/8r
\end{align*}
\hfill(40 points)

The term $\left(\dfrac{R_{Io}}{r}\right)^2$ is too small and it can be
ignored, so that
\[
  F_{\mathrm{tidal}} = -\frac{3GM_J m_{Io} R_{Io}}{4r^3}
\]
\hfill(10 points)

\medskip
\textbf{Calculate the difference of the tidal force}

The difference in the tidal force experienced at perijove (closest to Jupiter)
and apojove (furthest from Jupiter)
\[
  \Delta F = \left|F_{\mathrm{tidal}}(r_{\mathrm{peri}})
    - F_{\mathrm{tidal}}(r_{\mathrm{apo}})\right|
\]
\hfill(15 points)
\[
  = \frac{3}{4}GM_J m_{Io} R_{Io}
    \left[\frac{1}{r_{\mathrm{peri}}^3} - \frac{1}{r_{\mathrm{apo}}^3}\right]
  = 4.97\times10^{18}\ \mathrm{N}
\]
\hfill(15 points)

\medskip
\textbf{Calculate the average power of work done}

Hence, the average power of the work done on the rock during one-half orbit
is
\[
  \bar P = \frac{\Delta F\,d}{t_{Io}/2}
\]
\hfill(10 points)
\[
  \bar P = 3.25\times10^{15}\,W
\]
\hfill(10 points)

\solhead{2.28}{}
% source id: 2015_TH_S2
\textbf{First Version}

Kepler third Law for satelite % sic: "satelite" in source
and Moon: \hfill(30 points)
\begin{align*}
  \frac{P_S^2}{a_S^3} &= \frac{P_S^2}{(15.3\,R_P)^3} = \frac{4\pi^2}{GM_P}; \\
  \frac{P_M^2}{a_M^3} &= \frac{P_M^2}{(60.3\,R_E)^3} = \frac{4\pi^2}{GM_E};
\end{align*}

Express mass in term of mass density: \hfill(30 points)
\begin{align*}
  \frac{P_S^2}{(15.3\,R_P)^3} &= \frac{4\pi^2}{GM_P}
    = \frac{4\pi^2}{G\rho_P \frac{4}{3}\pi R_P^3} \\
  \frac{P_M^2}{(60.3\,R_E)^3} &= \frac{4\pi^2}{GM_E}
    = \frac{4\pi^2}{G\rho_E \frac{4}{3}\pi R_E^3}
\end{align*}

Evaluate ratio of mass-density: \hfill(30 points)
\[ \frac{\rho_E}{\rho_P}
   = \frac{60.3^3 \times P_S^2}{15.3^3 \times P_M^2} \]
% sic: the source labels this ratio (and the next) rho_E/rho_P; the value
% 0.238 is actually rho_P/rho_E, as the Second Alternate Version confirms
\[ \frac{\rho_E}{\rho_P} = 0.238 \approx 0.24 \]
\hfill(10 points)

\textbf{Second Alternate Version}

\begin{itemize}
\item The period of the Satellite:
  $P_S = 7\ \text{days},\ 3\ \text{hours},\ 43\ \text{minutes}
   = 0.0196\ \text{years}$,
\item Mean orbital radius of the Satellite: $a_S = 15.3\,R_P$; $R_P$: radius
  of the Planet.
\item The period of the Moon:
  $P_M = 27\ \text{days},\ 7\ \text{hours},\ 43\ \text{minutes}
   = 0.0749\ \text{years}$,
\item Mean orbital radius of the Moon: $a_M = 60.3\,R_E$; $R_E$: radius of
  the Earth.
\item Mass of the Sun ($M$)
\item Kepler's Third Law: $G(M + m) = 4\pi^2 \dfrac{a^3}{T^2}$
  $\rightarrow$
  \begin{itemize}
  \item Mass of the Planet: $m_P$
  \item Radius of the Planet: $R_P$
  \item Semi major axis of the Planet $=$ Mean orbital radius of the
    Plant: % sic: "Plant" in source
    $a_P$
  \item Semi major axis of the Satellite $=$ Mean orbital radius of the
    Satellite: $a_S$
  \item Semi major axis of the Moon $=$ Mean orbital radius of the Moon:
    $a_M$
  \item Mass of the Earth: $m_E$
  \item Radius of the Earth: $R_E$
  \item Semi major axis of the Earth $=$ Mean orbital radius of the Earth:
    $a_E$
  \end{itemize}
\end{itemize}

\begin{description}
\item[Planet--Satellite:]
  $G(m_P + m_S) = 4\pi^2 \dfrac{a_S^3}{P_S^2}$, $m_P \gg m_S$
  $\rightarrow$ $G(m_P) = 4\pi^2 \dfrac{a_S^3}{P_S^2}$ \hfill(1)\quad(5 points)
\item[The Earth--the Moon:]
  $G(m_E + m_M) = 4\pi^2 \dfrac{a_M^3}{P_M^2}$, $m_E \gg m_M$
  $\rightarrow$ $G(m_E) = 4\pi^2 \dfrac{a_M^3}{P_M^2}$ \hfill(2)\quad(5 points)
\item[Planet--Sun:]
  $G(M + m_P) = 4\pi^2 \dfrac{a_P^3}{P_P^2}$, $M \gg m_P$
  $\rightarrow$ $G(M) = 4\pi^2 \dfrac{a_P^3}{P_P^2}$ \hfill(3)\quad(5 points)
\item[Earth--Sun:]
  $G(M + m_E) = 4\pi^2 \dfrac{a_E^3}{P_E^2}$, $M \gg m_E$
  $\rightarrow$ $G(M) = 4\pi^2 \dfrac{a_E^3}{P_E^2}$ \hfill(4)\quad(5 points)
\end{description}

From equation (1) and (3):
\[ \frac{m_P}{M}
   = \frac{4\pi^2 \frac{a_S^3}{P_S^2}}{4\pi^2 \frac{a_P^3}{P_P^2}}
   = \frac{(a_S)^3}{(a_P)^3}\,\frac{(P_P)^2}{(P_S)^2}
   = \frac{(15.3R_P)^3}{(a_P)^3}\,\frac{(P_P)^2}{(P_S)^2}
   = \frac{(15.3R_P)^3}{(a_P)^3}\,\frac{(P_P)^2}{(0.0196)^2} \]
\hfill(20 points)

From equation (2) and (4):
\[ \frac{m_E}{M}
   = \frac{4\pi^2 \frac{a_M^3}{P_M^2}}{4\pi^2 \frac{a_E^3}{P_E^2}}
   = \frac{(a_M)^3}{(a_E)^3}\,\frac{(P_E)^2}{(P_M)^2}
   = \frac{(60.3R_E)^3}{(a_E)^3}\,\frac{(P_E)^2}{(P_E)^2} % sic: denominator printed as (P_E)^2 in source; should be (P_M)^2
   = \frac{(60.3R_E)^3}{(a_E)^3}\,\frac{(P_E)^2}{(0.0749)^2}
   = \frac{(60.3R_E)^3}{(1)^3}\,\frac{(1)^2}{(0.0749)^2} \]
\hfill(20 points)

$\rightarrow$
\[ \frac{m_E}{m_P}
   = \frac{\dfrac{(15.3R_P)^3}{(a_P)^3}\,\dfrac{(P_P)^2}{(0.0196)^2}}
          {\dfrac{(60.3R_E)^3}{(1)^3}\,\dfrac{(1)^2}{(0.0749)^2}}
   = \frac{(15.3R_P)^3/(0.0196)^2}{(60.3R_E)^3/(0.0749)^2} \]
% sic: the source writes m_E/m_P for what is actually m_P/m_E
\hfill(20 points)

$\rightarrow$
\[ \frac{\rho_P}{\rho_E}
   = \frac{m_P/\mathrm{Vol}_P}{m_E/\mathrm{Vol}_E}
   = \frac{m_P/\frac{4}{3}\pi(R_P)^3}{m_E/\frac{4}{3}\pi(R_E)^3}
   = \frac{(15.3R_P)^3/(0.0196)^2}{(60.3R_E)^3/(0.0749)^2}
     \cdot \frac{\frac{4}{3}\pi(R_E)^3}{\frac{4}{3}\pi(R_P)^3}
   = \frac{\dfrac{(15.3)^3}{(0.0196)^2}}{\dfrac{(60.3)^3}{(0.0749)^2}}
   = 0.24 \]
\hfill(20 points)

\solhead{2.29}{}
% source id: 2015_TH_S5
\textbf{Identify expression for $r(\theta)$.} \hfill(30 points)

We know that
$r(\theta) = \dfrac{a\left(1 - e^2\right)}{1 - e\cos\theta}$. Hence % sic: sign convention with perihelion at theta = pi
\begin{align*}
  r(0) &= \frac{a\left(1 - e^2\right)}{1 - e} = a(1 + e) \\
  r\left(\frac{\pi}{2}\right)
    &= \frac{a\left(1 - e^2\right)}{1 - e\times0} = a\left(1 - e^2\right)
\end{align*}

\ioaafig{2015_TH_S5-s1.pdf}
\ioaafig{2015_TH_S5-s2.pdf}

\medskip
\textbf{Estimate area of sectors using area of triangles.} \hfill(20 points)

Estimate the area of swept out sectors as the area of sectors of circles
\begin{align*}
  A(S_1) &\approx 0.5\times\Delta\theta_1\times\left(r(0)\right)^2
    = \Delta\theta_1\left(a(1 + e)\right)^2 \\ % sic: the 0.5 factor is dropped on the right-hand sides
  A(S_2) &\approx 0.5\times\Delta\theta_2
    \times\left(r\left(\tfrac{\pi}{2}\right)\right)^2
    = \Delta\theta_2\left(a(1 - e^2)\right)^2
\end{align*}

\medskip
\textbf{Use Kepler's second law to find a relation between the ratio
$\frac{\Delta\theta_1}{\Delta\theta_2}$ and the eccentricity $e$.}
\hfill(30 points)

Use Kepler's second law to obtain that
\begin{align*}
  A(S_1) &= A(S_2) \\
  \Delta\theta_1\times(a(1 + e))^2 &= \Delta\theta_2\times(a(1 - e^2))^2
\end{align*}
Thus,
\[
  \frac{\Delta\theta_1}{\Delta\theta_2}
    = \left(\frac{1 - e^2}{1 + e}\right)^2 = (1 - e)^2
    = \frac{2'44''}{21'17''} \approx 0.12843
\]

\medskip
\textbf{Obtain an estimate value of the eccentricity.} \hfill(20 points)

Thus, the eccentricity is $e \approx 1 - \sqrt{0.12843} = 0.64$.

\solhead{2.30}{True or False}
% source id: 2016_TH_T1
\begin{description}
\item[(T1.1)] \textbf{True.} \hfill(2 points)

  The colour of the clear sky during night is the same as during daytime,
  since the spectrum of sunlight reflected by the Moon is almost the same as
  the spectrum of sunlight. Only the intensity is lower.

\item[(T1.2)] \textbf{True.} \hfill(2 points)

  If the Earth's axis were perpendicular to its orbital plane, the celestial
  equator will coincide with ecliptic and the Sun will remain along the
  celestial equator every day. However, as the Earth's orbit is elliptical,
  the true sun would still lead or lag mean sun by a few minutes on different
  days of year.

\item[(T1.3)] \textbf{False.} \hfill(2 points)

  The semi-major axis of the orbit of the body will be less than that of
  Uranus. However the minor body's orbit may have a high eccentricity, in
  which case it may go outside that of Uranus.

\item[(T1.4)] \textbf{False.} \hfill(2 points)

  The centre of mass of Sun--Jupiter pair is just outside the Sun. Thus, if
  all gas giants are on same side of the Sun, the centre of mass of Solar
  system is definitely outside the Sun.

\item[(T1.5)] \textbf{True.} \hfill(2 points)

  For photons the wavelength increases when the Universe expands.
\end{description}

\solhead{2.31}{Mars Orbiter Mission}
% source id: 2016_TH_T9
\begin{description}
\item[(T9.1)]
Let apogee and perigee distances be $r_{\mathrm{a}}$ and $r_{\mathrm{p}}$
respectively.
\begin{align*}
r_{\mathrm{p}} &= R_\oplus + h^i_{\mathrm{p}} = (6371 + 264)\,\mathrm{km}
  = 6635\,\mathrm{km} \\
r_{\mathrm{a}} &= R_\oplus + h^i_{\mathrm{a}} = (6371 + 23904)\,\mathrm{km}
  = 30\,275\,\mathrm{km}
\end{align*}
Conservation of energy and angular momentum gives
\[ E = -\frac{GmM_\oplus}{r_{\mathrm{p}} + r_{\mathrm{a}}} \]
Total energy at perigee
\begin{align*}
\frac{1}{2}m{v_{\mathrm{p}}}^2 - \frac{GmM_\oplus}{r_{\mathrm{p}}}
  &= E = -\frac{GmM_\oplus}{r_{\mathrm{p}} + r_{\mathrm{a}}} \\
\frac{1}{2}{v_{\mathrm{p}}}^2
  &= \frac{GM_\oplus}{r_{\mathrm{p}}}
     \left(1 - \frac{r_{\mathrm{p}}}{r_{\mathrm{p}} + r_{\mathrm{a}}}\right) \\
\therefore v_{\mathrm{p}}
  &= \sqrt{2\,\frac{GM_\oplus}{r_{\mathrm{a}} + r_{\mathrm{p}}}\,
     \frac{r_{\mathrm{a}}}{r_{\mathrm{p}}}} \\
  &= \sqrt{\frac{2 \times 6.674 \times 10^{-11} \times 5.972 \times 10^{24}
     \times 3.0275 \times 10^{7}}
     {6.635 \times 10^{6} \times (3.0275 + 6.635 \times 10^{6})}} \\
% sic: the source prints the denominator as "(3.0275 + 6.635 x 10^6)";
% it should read (3.0275 x 10^7 + 6.635 x 10^6).
  &= 9.929\,\mathrm{km\,s^{-1}}
\end{align*}
As the engine burn is just 41.6\,s, we assume that the entire impulse is
applied instantaneously at perigee. The impulse is
$J = 1.73 \times 10^{5}\,\mathrm{kg\,m\,s^{-1}}$. Note that the total mass of
MOM must include the fuel, so we have to use
\[ m = 500 + 852 = 1352\,\mathrm{kg}. \]
Change in velocity due to impulse at perigee is
\[ \Delta v = \frac{J}{m} = \frac{1.73 \times 10^{5}}{1352}
   = 128.0\,\mathrm{m\,s^{-1}} \]
The new velocity will be given by (we use $'$ symbol to denote quantities
after the first orbit-raising maneuvre)
\[ v'_{\mathrm{p}} = v_{\mathrm{p}} + \Delta v = 10.06\,\mathrm{km\,s^{-1}} \]
The perigee remains unchanged. So we get $r'_{\mathrm{p}} = r_{\mathrm{p}}$.
Since the satellite is moving faster, the new apogee will be higher.
\begin{align*}
v'_{\mathrm{p}}
  &= \sqrt{2GM_\oplus \frac{r'_{\mathrm{a}}}
     {r'_{\mathrm{p}}(r'_{\mathrm{a}} + r'_{\mathrm{p}})}} \\
\therefore 1 + \frac{r'_{\mathrm{p}}}{r'_{\mathrm{a}}}
  &= \frac{2GM_\oplus}{(v'_{\mathrm{p}})^2 \times r'_{\mathrm{p}}} \\
  &= \frac{2 \times 6.674 \times 10^{-11} \times 5.972 \times 10^{24}}
     {(10.06 \times 10^{3})^2 \times 6.635 \times 10^{6}} = 1.188 \\
r'_{\mathrm{a}} &= \frac{6635}{0.188} = 35\,380\,\mathrm{km} \\
h_{\mathrm{a}} &= 35380 - 6371 \\
  &= \boxed{29\,009\,\mathrm{km}}
\end{align*}
Acceptable range: $\pm 150$\,km

\item[(T9.2)]
As seen above,
\begin{align*}
\frac{r'_{\mathrm{p}}}{r'_{\mathrm{a}}} &= 0.188 = \frac{1-e}{1+e} \\
\therefore e &= \frac{1 - 0.188}{1 + 0.188} = \boxed{0.683}
\end{align*}
Acceptable range: $\pm 0.002$

The new orbital semi-major axis and orbital period will be,
\begin{align*}
a' &= \frac{r'_{\mathrm{a}} + r'_{\mathrm{p}}}{2} \\
   &= \frac{35380 + 6635}{2} = 20\,933\,\mathrm{km} \\
P &= 2\pi\sqrt{\frac{a'^3}{GM_\oplus}} \\
  &= 2\pi\sqrt{\frac{(2.0933 \times 10^{7})^3}
     {6.674 \times 10^{-11} \times 5.972 \times 10^{24}}}
   = \boxed{30\,136\,\mathrm{s} = 8.37\,\mathrm{h}}
\end{align*}
Acceptable range: $\pm 0.1$\,h
\end{description}

\solhead{2.32}{Distance of the Lagrangian point L2 of the Earth-Moon system}
% source id: 2019_TH_9
The relay satellite revolves on a circular orbit, due to the gravitational
forces of the Earth and the Moon. Its net acceleration is: \hfill(4 p)
\begin{equation*}
\left(r + \frac{M}{m + M}R\right)\omega^2
  = \frac{GM}{(R + r)^2} + \frac{Gm}{r^2}, \tag{9.1}
\end{equation*}
where $R$ and $r$ denote the (constant) distances between Earth and Moon, and
between Moon and the satellite, respectively, while $M$ and $m$ stand for the
masses of Earth and Moon, respectively. Using Kepler's third law, the angular
velocity of the satellite, which is equal to the angular velocity of the Moon
can be written as \hfill(2 p)
\begin{equation*}
\omega^2 = \frac{4\pi^2}{P^2} = \frac{G(M + m)}{R^3}. \tag{9.2}
\end{equation*}
Substituting Eq.\ (9.2) into Eq.\ (9.1), eliminating the gravitational
constant ($G$), and introducing the dimensionless quantity $x = r/R$, we
obtain \hfill(2 p + 2 p)
\begin{gather*}
\frac{(M + m)x + M}{R^2} = \frac{M}{R^2(1 + x)^2} + \frac{m}{R^2 x^2}
  \tag{9.3} \\
(M + m)x + M = \frac{M}{(1 + x)^2} + \frac{m}{x^2}. \tag{9.4}
\end{gather*}
Assuming that $x = r/R \ll 1$ we can apply the approximation
$(1 + x)^{-2} \approx 1 - 2x$, which leads to the following expression:
\hfill(2 p)
\begin{equation*}
(M + m)x + M = M(1 - 2x) + \frac{m}{x^2}, \tag{9.5}
\end{equation*}
from which we get \hfill(2 p)
\begin{equation*}
x^3 = \frac{m}{3M + m}. \tag{9.6}
\end{equation*}
Substituting the mass ratio $m/M = 0.0123$ of the Earth--Moon system one can
obtain \hfill(2 p)
\begin{equation*}
x = 0.159\,83, \tag{9.7}
\end{equation*}
and therefore \hfill(2 p)
\begin{equation*}
r = xR_{\mathrm{Moon}} = 0.159\,83 \times 384\,400\,\mathrm{km}
  = 61\,440\,\mathrm{km} \tag{9.8}
\end{equation*}
% sic: the source writes "xR_Moon" although 384 400 km is the Earth--Moon
% distance R, not the radius of the Moon.
From this result we should subtract the radius of the Moon, and such a way we
obtain that \hfill(2 p)
\begin{equation*}
h = r - R_{\mathrm{Moon}} = 61\,440\,\mathrm{km} - 1737\,\mathrm{km}
  = \boxed{59\,703\,\mathrm{km}} \tag{9.9}
\end{equation*}

Here is another solution, which is erroneous in principle (or say,
``approximative'') but numerically gives an almost equivalent result.

One might say, that $m \ll M$ and therefore $m$ can be neglected either in
the left hand side of Eq.\ (9.1) and in the r.h.s.\ of Eq.\ (9.2), i.e.\ one
can use the following approximation:
\begin{equation*}
M + m \approx M \tag{9.10}
\end{equation*}
In this case instead of Eq.\ (9.6) one can arrive at
\begin{equation*}
x^3 = \frac{m}{3M}, \tag{9.11}
\end{equation*}
which leads to
\begin{equation*}
x = 0.160\,05, \tag{9.12}
\end{equation*}
and therefore,
\[ r = xR_{\mathrm{Moon}} = 0.160\,05 \times 384\,400\,\mathrm{km}
   = 61\,524\,\mathrm{km} \]
and, finally
\[ h = r - R_{\mathrm{Moon}} = 61\,524\,\mathrm{km} - 1737\,\mathrm{km}
   = 59\,787\,\mathrm{km} \]
This latter should be accepted as a full point solution, too. If, however,
$m$ is dropped out only from one of the two equations, then at least 2 points
should be subtracted.

\solhead{2.33}{Occultation of a X-ray Source}
% source id: 2020_TH_6
\begin{description}
\item[(a)] The solution for first part is obvious from the figure below.
\ioaafig{2020_TH_6-s1.pdf}
\begin{align*}
\sin\theta_0 &= \left(\frac{R + h}{r}\right) \\
\therefore \theta_0 &= \sin^{-1}\left(\frac{R + h}{r}\right)
\end{align*}

\item[(b)] Let us say the satelite % sic
has to move in the orbit by $\Delta\theta/2$ for $0.5\Delta h/2$ change in
height.
\begin{align*}
\sin\left(\theta_0 \pm \frac{\Delta\theta}{2}\right)
  &= \left(\frac{R + h \pm 0.5\Delta h}{r}\right) \\
\sin\left(\theta_0 + \frac{\Delta\theta}{2}\right)
  - \sin\left(\theta_0 - \frac{\Delta\theta}{2}\right)
  &= \left(\frac{\Delta h}{r}\right) \\
\therefore 2\cos\theta_0 \sin\left(\frac{\Delta\theta}{2}\right)
  &= \left(\frac{\Delta h}{r}\right) \\
2 \times \frac{\Delta\theta}{2}\cos\theta_0
  &= \left(\frac{\Delta h}{r}\right) \\
\Delta\theta &= \left(\frac{\Delta h}{r\cos\theta_0}\right)
  = \frac{2\pi\Delta t}{P} \\
\therefore \Delta t &= \left(\frac{P\Delta h}{2\pi r\cos\theta_0}\right)
\end{align*}

\item[(c)] Let $\theta$ be the new angle for 50\% attenuation. In the diagram
below, The line of sight from the satelite % sic
($S$) to the source is at minimum distance from the centre of the earth ($C$)
at point $A$. The points $S$, $C$ and $B$ are in the equatorial plane with
line $BA$ normal to the equatorial plane. One also realises that the line
$BS$ is normal to the plane defined by $A$, $B$ and $C$.
\ioaafig{2020_TH_6-s2.pdf}
\begin{align*}
a &= r\cos\theta \\
c &= r\sin\theta \\
b &= a\tan\beta \\
(R + h)^2 &= b^2 + c^2 \\
  &= (r\cos\theta\tan\beta)^2 + (r\sin\theta)^2 \\
\therefore \frac{(R + h)^2}{r^2}
  &= \cos^2\theta\tan^2\beta + \sin^2\theta \\
  &= (1 - \sin^2\theta)\tan^2\beta + \sin^2\theta \\
  &= \tan^2\beta + \sin^2\theta(1 - \tan^2\beta) \\
\therefore \sin^2\theta
  &= \left(\frac{(R + h)^2}{r^2} - \tan^2\beta\right)
     \left(\frac{1}{(1 - \tan^2\beta)}\right)
\end{align*}
Let us say $\tan^2\beta = k$
\[
\therefore \sin\theta
  = \sqrt{\left(\frac{(R + h)^2}{r^2} - k\right)
    \left(\frac{1}{(1 - k)}\right)}
\]
Using similar analysis as the previous part,
\begin{align*}
\Delta\theta\cos\theta
  &= \sqrt{\frac{1}{(1 - k)}}
     \left[\sqrt{\left(\frac{(R + h + \frac{\Delta h}{2})^2}{r^2} - k\right)}
     - \sqrt{\left(\frac{(R + h - \frac{\Delta h}{2})^2}{r^2} - k\right)}
     \right] \\
  &= \sqrt{\frac{1}{(1 - k)}}
     \left(\frac{\Delta h (R + h)}{r^2}\right)
     \left(\frac{(R + h)^2}{r^2} - k\right)^{-0.5} \\
  &= \sqrt{\frac{1}{(1 - k)}
     \left(\frac{(R + h)^2}{r^2} - k\right)}
     \left(\frac{\Delta h}{r} \times \frac{(R + h)}{r}\right)
     \left(\left(\frac{(R + h)}{r}\right)^{2} - k\right)^{-1} \\
  &= \sin\theta\left(\frac{\Delta h\sin\theta_0}{r}\right)
     \left(\sin^2\theta_0 - \tan^2\beta\right)^{-1} \\
\therefore \Delta t
  &= \left(\frac{P\Delta h\tan\theta}{2\pi r}\right)
     \left(\frac{\sin\theta_0}{\sin^2\theta_0 - \tan^2\beta}\right)
\end{align*}
\end{description}

\solhead{2.34}{Minimum velocity of a projectile}
% source id: 2021_TH_Q11
\begin{description}
\item[(11.1)]
\ioaafig{2021_TH_Q11-s1.pdf}
\begin{equation*}
EM = -\frac{GMm}{R} + \tfrac{1}{2}mv^2 \tag*{[1 point]}
\end{equation*}
At points A and B, the $v$ is the same; so, if $v$ is minimum, EM is minimum:

We know that the mechanical energy for an elliptical orbit is given by:
\begin{equation*}
EM = -\frac{GMm}{2a} \tag*{[1 point]}
\end{equation*}
where $2a$ is the length of the major axis.

Since EM must be minimal, then $2a$ must be minimal \hfill[1 point]
\ioaafig{2021_TH_Q11-s2.pdf}

\textbf{O:} Center of the Earth and one of the focus of the ellipse.

\textbf{F:} the other focus of the ellipse
\begin{equation*}
OB + BF = 2a \tag*{[1 point]}
\end{equation*}
$OB + BF = 2a$, must be minimal.

F takes an arbitrary position, so to minimize, BF must be perpendicular to
OF:
\ioaafig{2021_TH_Q11-s3.pdf}
\begin{gather*}
OB + BF = R + \frac{R}{\sqrt{2}} \tag*{[1 point]} \\
2a = R + \frac{R}{\sqrt{2}}
\end{gather*}
This is the minimum value it can take

Correspondingly, the expression for the minimum velocity is
\begin{equation*}
v = \sqrt{\frac{2GM}{(1 + \sqrt{2}\,)R}} \tag*{[1 point]}
\end{equation*}
Using the M and R values for the Earth, the minimum velocity is
\begin{equation*}
v = 7195\ \frac{\mathrm{m}}{\mathrm{s}} \tag*{[1 point]}
\end{equation*}

\item[(11.2)] Furthermore, $OF = 2\,e{\cdot}a$, where $e$ is the eccentricity
\begin{gather*}
\frac{R}{\sqrt{2}} = 2\,e\,\frac{R}{2}\left(1 + \frac{1}{\sqrt{2}}\right)
  \tag*{[2 points]} \\
e = \frac{1}{1 + \sqrt{2}} = 0.414 \tag*{[1 point]}
\end{gather*}
\end{description}

\solhead{2.35}{Hodograph}
% source id: 2021_TH_Q12
\begin{description}
\item[(12.1)] For a satellite that describes a uniform circular motion of
radius $r$ around the star and mass M, we have:
\begin{equation*}
v = \left[R_{Hodograph}\right] = \sqrt{\frac{GM}{r}} \tag*{[1 point]}
\end{equation*}

\item[(12.2)]
\begin{gather*}
\vec{a} = -\frac{Gm}{r^2}\hat{r} \tag*{[1 point]} \\
% sic: the source prints "Gm" here rather than "GM".
L = mr^2\omega \tag*{[1 point]} \\
\rightarrow\ \vec{a} = -\frac{GMm\omega}{L}\hat{r}
  = \frac{\Delta\vec{v}}{\Delta t} \\
\frac{GMm}{L}\,\frac{\Delta\theta}{\Delta t}
  = \left|\frac{\Delta v}{\Delta t}\right| \\
\Delta v = \pm\frac{GMm}{L}\Delta\theta \\
k = \frac{GMm}{L} \tag*{[2 points]}
\end{gather*}

\item[(12.3)]
\begin{gather*}
L = mv_c R \tag*{[1 point]} \\
k_c = \frac{GMm}{L} = v_c = \sqrt{\frac{GM}{R}} \tag*{[1 point]}
\end{gather*}

\item[(12.4)] In the following scheme, the hodographic circumference has
been constructed as follows:
\begin{itemize}
\item For $\theta = 0$, the velocity vector is drawn in arbitrary units (3,
  for instance) corresponding to the velocity in the periastron, $v_P$. Its
  head determines the point $A$.
\item For $\theta = \pi$, the velocity vector is drawn diametrically
  opposite in the apastron, $v_A$, also in arbitrary units. Its head
  determines point $B$. The hodograph must pass through points $A$ and $B$.
  Therefore, the radius of the hodograph must be equal to:
\begin{equation*}
R_{Hodograph} = \frac{v_P + v_A}{2} \tag*{[1 point]}
\end{equation*}
  And the ``distance'' $d$ from the star to the center of the circumference
  will be:
\begin{equation*}
d = \frac{v_P - v_A}{2} \tag*{[1 point]}
\end{equation*}
\end{itemize}
\ioaafig{2021_TH_Q12-s1.pdf}

\item[(12.5)] In the following scheme, the hodographic circumference has
been constructed as follows:
\begin{itemize}
\item For $\theta = 0$, the velocity vector is drawn in arbitrary units (4,
  for instance) corresponding to the velocity in the periastron, that is
  $v_P$. Its head determines the point $A$. Then the velocity vector is
  drawn diametrically opposite in the apastron, that is, $v_A = 0$. Its head
  determines point B that coincides with the star. The hodograph must pass
  through points A and B. Therefore, the radius of the hodograph must be
  equal to:
\begin{equation*}
R_{Hodograph} = \frac{v_P}{2} \tag*{[1 point]}
\end{equation*}
  And the distance $d$ from the star to the center of the hodographic
  circumference
\begin{equation*}
d = \frac{v_P}{2} \tag*{[1 point]}
\end{equation*}
\end{itemize}
\ioaafig{2021_TH_Q12-s2.pdf}
\end{description}

\solhead{2.36}{Formation of the Venus-2}
% source id: 2021_TH_Q14
\ioaafig{2021_TH_Q14-s1.pdf}

\begin{description}
\item[14.1] Let $M$ be the mass of the Sun, then
  \begin{align*}
    \frac{GMm}{R_0^2} &= \frac{mv_0^2}{R_0} \\
    v_0^2 &= \frac{GM}{R_0}
  \end{align*}
  \hfill[1 point]

\item[14.2] The mechanical energy of Venus (before collision) is:
  \[
    E_i = -\frac{GMm}{R_0} + \frac{1}{2}mv_0^2 = -\frac{GMm}{2R_0}
  \]
  \hfill[1 point]

\item[14.3] Since the comet moves radially towards the Sun, it has no angular
  momentum (with respect to the origin in the Sun). Then, the angular momentum
  of Venus is the same as that of ``Venus-2''
  \[
    L = R_0 m v_0
  \]
  \hfill[2 point] % sic: "[2 point]"

  This allows us to find the angular component of the velocity of Venus,
  $v_\theta$, just after the collision. The angular momentum just after the
  collision is
  \[
    L = R_0 m v_0 = R_0(m + \alpha m)v_\theta
  \]
  \hfill[1 points] % sic: "[1 points]"

  Since the angular momentum is conserved it follows that
  \[
    v_\theta = \frac{v_0}{1 + \alpha}
  \]
  The comet has zero energy, so
  \[
    K = -U = +\frac{GM\alpha m}{R_0} = \frac{1}{2}\alpha m v_c^2
  \]
  \hfill[2 points]

  where $v_c$ is the velocity of the comet just before the collision. It
  follows that
  \[
    v_c^2 = \frac{2GM}{R_0} = 2v_0^2
  \]
  \hfill[1 point]

  The collision between the comet and Venus is inelastic, thus the total
  energy is not conserved. Also notice that the mass of ``Venus-2'' is
  $(1 + \alpha)m$. \hfill[2 point] % sic: "[2 point]"

  We now apply the conservation of linear momentum. Along the radial direction
  we get
  \[
    \alpha m v_c = (m + \alpha m)v_r
  \]
  \hfill[1 point]

  where $v_r$ is the velocity of ``Venus-2'' just after the collision. It
  follows that
  \[
    v_r = \frac{\alpha}{1 + \alpha}v_c = \frac{\sqrt{2}\,\alpha}{1 + \alpha}v_0
  \]
  \hfill[1 point]

\item[14.4] Venus-2 mechanical energy (we evaluate it just after the
  collision) is:
  \[
    E_f = U + K = -\frac{GMm(1 + \alpha)}{R_0}
      + \frac{1}{2}m(1 + \alpha)\left(v_\theta^2 + v_r^2\right)
  \]
  \hfill[4 points]
  \begin{align*}
    &= -(1 + \alpha)\frac{GMm}{R_0} + \frac{1}{2}(1 + \alpha)
      \left[2\left(\frac{\alpha}{1 + \alpha}\right)^2
      + \frac{1}{(1 + \alpha)^2}\right]\frac{GMm}{R_0} \\
    &= -\frac{GMm}{2R_0}\left(\frac{1 + 4\alpha}{1 + \alpha}\right)
      = E_i\left(\frac{1 + 4\alpha}{1 + \alpha}\right)
  \end{align*}
  \hfill[1 point]

\item[14.5] ``Venus-2'' orbit is no longer a circle. Since the energy is
  negative it must, therefore, be elliptical. It follows that
  \hfill[1 points] % sic: "[1 points]"
  \[
    E_f = -G\frac{M(1 + \alpha)m}{2a_f}
  \]
  \hfill[2 points]

  where $a_f$ is the semi-major axis of the orbit of Venus-2. Now, solving for
  $a_f$:
  \begin{align*}
    -G\frac{M(1 + \alpha)m}{2a_f}
      &= E_i\left(\frac{1 + 4\alpha}{1 + \alpha}\right)
       = -G\frac{Mm}{2R_0}\left(\frac{1 + 4\alpha}{1 + \alpha}\right) \\
    \frac{(1 + \alpha)}{a_f}
      &= \frac{1}{R_0}\left(\frac{1 + 4\alpha}{1 + \alpha}\right) \\
    a_f &= R_0\,\frac{(1 + \alpha)^2}{1 + 4\alpha}
  \end{align*}
  \hfill[2 points]

\item[14.6] Using Kepler's third law, we have that
  \[
    \frac{T_f}{T_0} = \left(\frac{a_f}{R_0}\right)^{3/2}
  \]
  \hfill[2 points]
  \[
    \frac{T_f}{T_0}
      = \left(\frac{(1 + \alpha)^2}{1 + 4\alpha}\right)^{3/2}
      = \frac{(1 + \alpha)^3}{(1 + 4\alpha)^{3/2}}
  \]
  \hfill[1 points] % sic: "[1 points]"

  For $\alpha < 2$, the year is shorter, for $\alpha > 2$, the year is longer,
  while for $\alpha = 2$ the duration of the year does not change.

\item[14.7] For the `Venus-2' to collide with the Sun, the perihelion radius
  of the post-collision orbit should be $R_\odot$. \hfill[1 point]
  \[
    r_p = R_\odot = 6.955\times10^{8}\,m
      = \frac{6.955\times10^{8}}{0.723\times1.496\times10^{11}}R_0
      = 0.00643R_0
  \]
  Let $x = \dfrac{R_\odot}{R_0} = 0.00643$.

  Now, we apply conservation of angular momentum right after the collision and
  at the moment ``Venus-2'' is at the perihelion:
  \[
    L = mv_0R_0 = (1 + \alpha)mv_pR_\odot
  \]
  where $v_p$ is the speed of Venus-2 at the perihelion. Solving for $v_p$
  \[
    v_p = \frac{v_0}{x(1 + \alpha_c)}
  \]
  Applying conservation of energy between the two same instants we obtain:
  \begin{align*}
    E_f &= -\frac{GM(1 + \alpha_c)m}{R_\odot}
      + \frac{1}{2}m(1 + \alpha_c)v_p^2 \\
    -\left(\frac{1 + 4\alpha_c}{1 + \alpha_c}\right)\frac{GMm}{2R_0}
      &= -\frac{GM(1 + \alpha_c)m}{xR_0}
      + \frac{1}{2}m(1 + \alpha_c)
        \left(\frac{v_0}{x(1 + \alpha_c)}\right)^2 \\
    -\left(\frac{1 + 4\alpha_c}{1 + \alpha_c}\right)\frac{GMm}{2R_0}
      &= -\frac{(1 + \alpha_c)}{x}\frac{GMm}{R_0}
      + \frac{1}{2}(1 + \alpha_c)
        \left(\frac{1}{x(1 + \alpha_c)}\right)^2\frac{GMm}{R_0} \\
    -\frac{1}{2}\left(\frac{1 + 4\alpha_c}{1 + \alpha_c}\right)
      &= -\frac{(1 + \alpha_c)}{x}
      + \frac{1}{2}(1 + \alpha_c)
        \left(\frac{1}{x(1 + \alpha_c)}\right)^2 \\
    (1 + 4\alpha_c) &= \frac{2(1 + \alpha_c)^2}{x} - \frac{1}{x^2}
  \end{align*}
  \hfill[3 points]
  \begin{align*}
    2x\alpha_c^2 + 4x(1 - x)\alpha_c - (1 - x)^2 &= 0 \\
    \alpha_c = (1 - x)\left(\sqrt{1 + \frac{1}{2x}} - 1\right) & \\
    \alpha_c = 7.824 &
  \end{align*}
  \hfill[1 point]

\item[14.8] For post-collision orbit,
  \begin{align*}
    v_\theta &= \frac{1}{1 + \alpha_c}v_0 \\
    v_r &= \frac{\sqrt{2}\,\alpha_c}{1 + \alpha_c}v_0 \\
    v_f &= \sqrt{v_r^2 + v_\theta^2}
  \end{align*}
  \hfill[1 point]
  \[
    v_f = \frac{\sqrt{2\alpha_c^2 + 1}}{1 + \alpha_c}v_0
  \]
  \hfill[3 point] % sic: "[3 point]"

  Thus,
  \begin{align*}
    \delta v &= v_0 - v_f \\
      &= \left(1 - \frac{\sqrt{2\alpha_c^2 + 1}}{1 + \alpha_c}\right)v_0 \\
      &= -0.259\,v_0
  \end{align*}
  \[
    \frac{\delta v}{v_0} = 25.9\%
  \]
  \hfill[0.5 points]
  \begin{align*}
    \tan\delta\theta &= \left(\frac{v_r}{v_\theta}\right) \\
    \tan\delta\theta &= \left(\sqrt{2}\,\alpha_c\right) \\
    \delta\theta &= 1.480\ \mathrm{rad}
  \end{align*}
  \hfill[0.5 points]
  \[
    \delta\theta = 84.84^\circ
  \]
\end{description}

\solhead{2.37}{Mars}
% source id: 2021_TH_Q3
% note: the official 2021 T3 solution PDF is a single page which ends after
% "Solution A" of part 3.2; the announced alternative method ("next") and
% part 3.3 are not present in the source file.
% printed.
\begin{description}
\item[(3.1)] Since the path between $A$ and $B$ is parabolic, the total
energy of the spacecraft is zero,
\[ E_{AB} = 0, \qquad \leftrightarrow \qquad \varepsilon = 1. \]
So, when you get to point B, then
\begin{equation*}
E = \tfrac{1}{2}mv_B^2 - \frac{GMm}{r_B} = 0; \tag*{[1 point]}
\end{equation*}
From where
\begin{equation*}
v_B = \sqrt{\frac{2GM}{r_B}} \tag*{[1 point]}
\end{equation*}
Replacing with available values
\begin{equation*}
v_B = \sqrt{\frac{2 \times (6.67 \times 10^{-11}) \times
  (6.4 \times 10^{23})}{6.8 \times 10^{6}}}
  \approx 3.55\ \mathrm{km\,s^{-1}} \tag*{[1 point]}
\end{equation*}

\item[(3.2)] There are two possible solutions

\textbf{Solution A}

Immediately after breaking % sic: "braking"
the total energy and the angular momentum change since the braking force is
tangential, but along the elliptical path BC the values with which it passes
through are conserved and take on during the entire journey. C; then applying
conservation of energy between points B$'$ (immediately after braking) and C
as well as conservation of the angular momentum between B$'$ and C, we have
\begin{gather*}
\tfrac{1}{2}mv_{B'}^2 - \frac{GMm}{r_B}
  = \tfrac{1}{2}mv_C^2 - \frac{GMm}{R_M} \quad Eq.~1
  \tag*{[1 point]} \\
L = mv_{B'}r_B = mv_C R_M \quad Eq.~2 \tag*{[1 point]}
\end{gather*}
From these two equations we solve for $v_C$ and it's result is used to
calculate the total energy in C which gives us:
\begin{equation*}
E_{elipse} = -\frac{GMm}{R_M + r_B} \approx -2.10 \times 10^{11}\,J
  \quad Eq.~3a \tag*{[2 points]}
\end{equation*}
In both methods, this and next, two points for getting the equation and two
points for calculation + negative sign + unit.
\end{description}

\solhead{2.38}{Pluto Satellites}
% source id: 2021_TH_Q9
\ioaafig{2021_TH_Q9-s1.pdf}
Let P and C be the centers of Pluto and Charon respectively, and CM their
center of mass. Clearly $\overline{CM - P} = \frac{1}{9}R$ and
$\overline{CM - C} = \frac{8}{9}R$. The effective gravitational field at
points A and B are given by:
\begin{center}
[By recognizing the vectorial nature of gravities/forces/fields]
\hfill [4 points]
\end{center}
\begin{gather*}
g_A = \frac{GM_P}{(R-r)^2}\hat{r} - \frac{GM_C}{r^2}\hat{r}
  - \omega^2\left(\tfrac{8}{9}R - r\right)\hat{r} \tag*{[1 point]} \\
g_B = \frac{GM_P}{(R+r)^2}\hat{r} + \frac{GM_C}{r^2}\hat{r}
  - \omega^2\left(\tfrac{8}{9}R + r\right)\hat{r} \tag*{[1 point]}
\end{gather*}
where $\omega^2 = 9\dfrac{GM_C}{R^3}$ is the angular speed common to both
Pluto and Charon. Thus, the gravitational accelerations at these points are:
\begin{center}
[by finding angular speed] [2 points]
\end{center}
\begin{gather*}
g_A = \frac{GM_P}{(R-r)^2} - \frac{GM_C}{r^2}
  - \frac{GM_C}{R^2}\left(8 - 9\frac{r}{R}\right) \\
g_B = \frac{GM_P}{(R+r)^2} + \frac{GM_C}{r^2}
  - \frac{GM_C}{R^2}\left(8 + 9\frac{r}{R}\right)
\end{gather*}
Factorizing the expression to obtain formulas proportional to Charon's mass
[1 point]:
\begin{gather*}
g_A = -\frac{GM_C}{r^2}\left[1
  - \frac{M_P}{M_C}\left(\frac{r}{R}\right)^2
    \frac{1}{\left(1 - \frac{r}{R}\right)^2}
  + \left(\frac{r}{R}\right)^2\left(8 - 9\frac{r}{R}\right)\right] \\
g_B = \frac{GM_C}{r^2}\left[1
  + \frac{M_P}{M_C}\left(\frac{r}{R}\right)^2
    \frac{1}{\left(1 + \frac{r}{R}\right)^2}
  - \left(\frac{r}{R}\right)^2\left(8 + 9\frac{r}{R}\right)\right]
\end{gather*}
Replacing the numerical value of the ratio of the two masses:
\begin{gather*}
g_A = -\frac{GM_C}{r^2}\left[1
  - \left(\frac{r}{R}\right)^2\frac{8}{\left(1 - \frac{r}{R}\right)^2}
  + \left(\frac{r}{R}\right)^2\left(8 - 9\frac{r}{R}\right)\right] \\
g_B = \frac{GM_C}{r^2}\left[1
  + \left(\frac{r}{R}\right)^2\frac{8}{\left(1 + \frac{r}{R}\right)^2}
  - \left(\frac{r}{R}\right)^2\left(8 + 9\frac{r}{R}\right)\right]
\end{gather*}
Noting that the expressions in square brackets are dimensionless, it is
helpful to replace the given values of $r$ and $R$, resulting [2 points by
expressions of $g_A$ and $g_B$]:
\begin{gather*}
g_A = -\frac{GM_C}{r^2}\,(0.99929) \\
g_B = \frac{GM_C}{r^2}\,(0.99933)
\end{gather*}
The percentage difference with respect to Charon's normal gravity
$g_0 = \dfrac{GM_C}{r^2}$ is
\begin{equation*}
\frac{\big||g_A| - |g_B|\big|}{g_0}\cdot 100\% = 0.004\%
  = 4 \times 10^{-3}\,\% \tag*{[1 point]}
\end{equation*}
% kept as printed.

\solhead{2.39}{Co-orbital satellites}
% source id: 2022_TH_12
\begin{description}
\item[(a)] By Kepler's third law we have
  \begin{align*}
    T_i^2 &= \frac{4\pi^2}{GM}r_i^3 \\
    \omega_i &= \frac{2\pi}{T_i} = \sqrt{\frac{GM}{r_i^3}}
  \end{align*}
  \hfill(1.5 points)
  \begin{align*}
    L_i &= m_i\omega_i r_i^2 \\
    \therefore L_i &= m_i\sqrt{GMr_i}
  \end{align*}
  \hfill(1.5 points)

\item[(b)] When the satellites are opposing positions (i.e.\ $\theta = \pi$),
  we have % sic: "are opposing positions"
  \[
    r_1 = R \pm \frac{x_1}{2}, \quad r_2 = R \mp \frac{x_2}{2}
  \]
  \hfill(2.0 points)

  Thus, using the conservation of angular momentum
  \begin{align*}
    L_{tot} &= m_1\sqrt{GM\left(R + \frac{x_1}{2}\right)}
             + m_2\sqrt{GM\left(R - \frac{x_2}{2}\right)} \\
            &= m_1\sqrt{GM\left(R - \frac{x_1}{2}\right)}
             + m_2\sqrt{GM\left(R + \frac{x_2}{2}\right)}
  \end{align*}
  \hfill(4.0 points)
  \begin{align*}
    \therefore\ m_1\left(1 + \frac{x_1}{4R}\right)
      + m_2\left(1 - \frac{x_2}{4R}\right)
      &= m_1\left(1 - \frac{x_1}{4R}\right)
      + m_2\left(1 + \frac{x_2}{4R}\right) \\
    m_1\frac{x_1}{2R} &= m_2\frac{x_2}{2R} \\
    \therefore\ \frac{m_1}{m_2} &= \frac{x_2}{x_1}
  \end{align*}
  \hfill(2.0 points)

\item[(c)] Let the centre of mass of the ternary system be denoted by C. The
  total gravitational force on satellite $m_2$ due to the primary mass $M$,
  and the satellite $m_1$ may be resolved into a radial component, parallel to
  the line between $m_2$ and C and a tangential component, perpendicular to
  the line between $m_2$ and C. Only the tangential component of the force
  will act to change the angular momentum of the satellite $m_2$.

  The tangential component of the force is given by
  \[
    F_{2_\perp} = \frac{GMm_2}{r_2^2}\sin(\delta\beta)
      - \frac{Gm_1m_2}{d^2}\sin\beta
  \]
  \hfill(2.0 points)

  where $d = \mathit{dist}(m_1, m_2)$, $\beta = \measuredangle m_1 m_2 C$, and
  $\delta\beta = \measuredangle M m_2 C$. There are two ways to simplify this
  expression.

  The first way straightforward. By immediately exploiting the fact that % sic: "way straightforward"
  $m_1, m_2 \ll M$ and hence $\delta\beta \to 0$, we can say,
  \begin{align*}
    \frac{d}{\sin\theta} &= \frac{r_1}{\sin(\beta + \delta\beta)} \\
    \sin\beta &= \left(\frac{r_1}{d}\right)\sin\theta
  \end{align*}
  \hfill(3.0 points)
  \begin{align*}
    \text{Also,}\quad
    \frac{\left(\dfrac{r_1 m_1}{M + m_1}\right)}{\sin(\delta\beta)}
      &= \frac{r_2}{\sin(180 - \theta - \delta\beta)} \\
    \therefore\ \sin(\delta\beta)
      &= \left(\frac{m_1 r_1}{r_2(m_1 + M)}\right)\sin\theta
  \end{align*}
  \hfill(4.0 points)

  Another way is using the sine and cosine rules for triangles, we have the
  expressions
  \begin{align*}
    \frac{r_2}{\sin(180 - \gamma)}
      &= \frac{\left[\left(\frac{r_1 m_1}{M + m_1}\right)^2 + r_2^2
         - \frac{2m_1 r_1 r_2}{M + m_1}\cos\theta\right]^{\frac{1}{2}}}
         {\sin\theta} \\
    \frac{r_2\sin\theta}{\sin\gamma}
      &= \left[r_2^2 + \left(\frac{r_1 m_1}{M + m_1}\right)^2
         - \frac{2m_1 r_1 r_2}{M + m_1}\cos\theta\right]^{\frac{1}{2}} \\
    \frac{d}{\sin\gamma} &= \frac{\frac{Mr_1}{m_1 + M}}{\sin\beta} \\
    \sin\beta &= \frac{Mr_1}{m_1 + M}\,\frac{r_2}{d}\sin\theta
       \left[r_2^2 + \left(\frac{m_1 r_1}{m_1 + M}\right)^2
       - \frac{2m_1 r_1 r_2}{m_1 + M}\cos\theta\right]^{-\frac{1}{2}}
       \approx \frac{r_1}{d}\sin\theta \\
    \sin(\delta\beta) &= \frac{m_1 r_1}{m_1 + M}\sin\theta
       \left[r_2^2 + \left(\frac{m_1 r_1}{m_1 + M}\right)^2
       - \frac{2m_1 r_1 r_2}{m_1 + M}\cos\theta\right]^{-\frac{1}{2}}
       \approx \frac{m_1 r_1}{r_2(m_1 + M)}\sin\theta
  \end{align*}

  In both cases, finally we obtain
  \begin{align*}
    F_{2_\perp} &= \frac{GMm_2}{r_2^2}\,\frac{m_1 r_1}{r_2(m_1 + M)}\sin\theta
      - \frac{Gm_1 m_2}{d^2}\,\frac{r_1}{d}\sin\theta \\
    \intertext{As, $M + m_1 \approx M$}
    \therefore\ F_{2_\perp} &= Gm_1 m_2 r_1\sin\theta\,(r_2^{-3} - d^{-3})
  \end{align*}
  \hfill(3.0 points)

\ioaafig{2022_TH_12-s1.pdf}

  The distance between satellites
  \[
    d \approx 2R\sin\left(\frac{\theta}{2}\right)
  \]
  \hfill(1.0 points)

  The tangential component of the gravitational force then becomes
  \begin{align*}
    F_{2_\perp} &\approx \frac{Gm_1 m_2}{R^2}\sin\theta
      \left(1 - \frac{1}{8\sin^3\left(\frac{\theta}{2}\right)}\right) \\
    F_{2_\perp} &= \frac{Gm_1 m_2}{R^2}
      \left(\sin\theta - \frac{\cos\left(\frac{\theta}{2}\right)}
        {4\sin^2\left(\frac{\theta}{2}\right)}\right)
  \end{align*}
  \hfill(2.0 points)

  Finally, the torque generated by the gravitational force is then
  \[
    \frac{\Delta L_2}{\Delta t} = F_{2_\perp} r
      \approx -\frac{Gm_1 m_2}{R}
      \left(\frac{\cos\left(\frac{\theta}{2}\right)}
        {4\sin^2\left(\frac{\theta}{2}\right)} - \sin\theta\right)
  \]
  \hfill(3.0 points)

\item[(d)] We can see that
  \begin{align*}
    \Delta L_i &= m_i\sqrt{GM(r_i + \Delta r_i)} - m_i\sqrt{GMr_i} \\
      &= m_i\sqrt{GMr_i}
         \left[\left(1 + \frac{\Delta r_i}{r_i}\right)^{\frac{1}{2}} - 1\right] \\
      &= m_i\sqrt{GMr_i}\left(\frac{\Delta r_i}{2r_i}\right) \\
    \therefore\ \frac{\Delta L_i}{\Delta t}
      &= \frac{1}{2}m_i\sqrt{\frac{GM}{r_i}}\,\frac{\Delta r_i}{\Delta t}
  \end{align*}
  \hfill(2.0 points)
  \begin{align*}
    \frac{\Delta r_i}{\Delta t}
      &= \frac{2}{m_i}\sqrt{\frac{r_i}{GM}}\,\frac{\Delta L_i}{\Delta t} \\
    \therefore\ \frac{\Delta r_2}{\Delta t}
      &\approx -\frac{2}{m_2}\sqrt{\frac{r_2}{GM}}\,
        \frac{Gm_1 m_2}{R}\,h(\theta) \\
    \frac{\Delta r_2}{\Delta t}
      &\approx -2m_1\sqrt{\frac{G}{MR}}\,h(\theta)
  \end{align*}
  \hfill(1.0 points)

  Again, since $r1 \approx r2 \approx R$, we can say by symmetry/the % sic: upright r1, r2 in source
  conservation of angular momentum
  \[
    \frac{\Delta r_1}{\Delta t} \approx 2m_2\sqrt{\frac{G}{MR}}\,h(\theta)
  \]
  \hfill(3.0 points)
  \begin{align*}
    \therefore\ \frac{\Delta s}{\Delta t}
      &= \frac{\Delta r_2}{\Delta t} - \frac{\Delta r_1}{\Delta t} \\
    \frac{\Delta s}{\Delta t}
      &= -2\sqrt{\frac{G}{MR}}\,(m_1 + m_2)\,h(\theta)
  \end{align*}
  \hfill(2.0 points)

\item[(e)] Since $\theta$ is the angle between the two satellites, we have
  \[
    \frac{\Delta\theta}{\Delta t} = \omega_2 - \omega_1
      = \sqrt{\frac{GM}{r_2^3}} - \sqrt{\frac{GM}{r_1^3}}
  \]
  \hfill(2.0 points)
  \begin{align*}
    &= \sqrt{\frac{GM}{R^3}}
       \left[\left(\frac{r_2}{R}\right)^{-\frac{3}{2}}
       - \left(\frac{r_1}{R}\right)^{-\frac{3}{2}}\right] \\
    &= \sqrt{\frac{GM}{R^3}}
       \left[\left(1 + \frac{r_2 - R}{R}\right)^{-\frac{3}{2}}
       - \left(1 + \frac{r_1 - R}{R}\right)^{-\frac{3}{2}}\right] \\
    &= \sqrt{\frac{GM}{R^3}}
       \left[1 - \frac{3}{2}\frac{(r_2 - R)}{R} - 1
       + \frac{3}{2}\frac{(r_1 - R)}{R}\right] \\
    &\approx \sqrt{\frac{GM}{R^3}}
       \left[\frac{3}{2}\frac{r_1 - r_2}{R}\right] \\
    \frac{\Delta\theta}{\Delta t}
    &\approx -\frac{3}{2}\sqrt{\frac{GM}{R^3}}\,\frac{s}{R}
  \end{align*}
  \hfill(3.0 points)

\item[(f)] Using the expressions in the previous two parts,
  \[
    \frac{\Delta s}{\Delta\theta}
      = \frac{\Delta s}{\Delta t}\cdot\frac{\Delta t}{\Delta\theta}
  \]
  \hfill(1.0 points)
  \begin{align*}
    &= 2\sqrt{\frac{G}{MR}}\,(m_1 + m_2)\,h(\theta)\cdot
       \frac{2}{3}\sqrt{\frac{R^3}{GM}}\,\frac{R}{s} \\
    &= \frac{4}{3}\frac{R^2}{M}(m_1 + m_2)\,h(\theta)
       \left(\frac{1}{s}\right) \\
    \therefore\ s\,\Delta s
    &= \frac{4R^2}{3}\,\frac{(m_1 + m_2)}{M}\,h(\theta)\,\Delta\theta
  \end{align*}
  \hfill(1.0 points)

\item[(g)] The minimum distance of 13\,000\,km corresponds to a minimum angle
  of
  \[
    \theta_{min} \approx \frac{13000}{150000}
      \approx 0.0868\,\mathrm{rad} = 4^\circ 58'
  \]
  \hfill(2.0 points)

  Then we substitute the given values the into given expression and the result % sic: "the into given"
  of (b) which gives
  \[
    \frac{m_2}{m_1} \approx 3.6
  \]
  \hfill(1.0 points)

  and
  \[
    m_1 + m_2 \approx 2.5\times10^{18}\,\mathrm{kg}
  \]
  \hfill(2.0 points)

  Finally
  \[
    m_1 \approx 5.3\times10^{17}\,\mathrm{kg} \quad
    m_2 \approx 1.9\times10^{18}\,\mathrm{kg}
  \]
  \hfill(1.0 points)
\end{description}

\solhead{2.40}{Ring of a planet}
% source id: 2022_TH_8
\begin{description}
\item[(a)] Consider small surface element ($\Delta S$) of the disk. The
  gravitational field ($\Delta E$) generated by this surface element at at a
  point $O$ located at a distance $a$ from $\Delta S$ (see the figure below): % sic: "at at"
  \[
    \Delta E = G\frac{\sigma\,\Delta S}{a^2}
  \]
  where $\sigma$ is a surface density of the mass.

\ioaafig{2022_TH_8-s1.pdf}

  where $\alpha$ is the angle between surface normal and the vector $\vec r$
  and $\Delta\Omega$ is infinitesimal solid angle corresponding to the area
  $\Delta S$.

  For a ring, the surface density will be given by
  $\sigma = M/\pi(R^2 - r^2)$.

  As the horizontal component will get cancelled due to symmetry
  considerations, we only consider the normal component:
  \[
    \Delta E_\perp = \Delta E\cos(\alpha)
      = \frac{G\sigma}{a^2}\,\Delta S\cos(\alpha)
      = \frac{G\sigma}{a^2}\,\Delta S_\perp
      = G\sigma\,\Delta\Omega
  \]
  \hfill(5.0 points)

  By Summation, we can obtain total normal field of the surface:
  \[
    E_\perp = G\sigma\Omega
  \]
  \hfill(1.0 points)

  here $\Omega$ is the total solid angle subtended by the disk. In our case,
  when we observe disk from the point of view of the test particle:
  \begin{align*}
    \Omega &= 2\pi\left(1 - \cos\theta_{out}\right)
              - 2\pi\left(1 - \cos\theta_{in}\right) \\
           &= 2\pi\left(1 - \frac{x}{\sqrt{x^2 + R^2}}\right)
              - 2\pi\left(1 - \frac{x}{\sqrt{x^2 + r^2}}\right) \\
    \Omega &= 2\pi x\left(\frac{1}{\sqrt{x^2 + r^2}}
              - \frac{1}{\sqrt{x^2 + R^2}}\right)
  \end{align*}
  \hfill(3.0 points)
  \begin{align*}
    \therefore F_\perp &= G\sigma\Omega m \\
    F_\perp &= \frac{2GMmx}{R^2 - r^2}
       \left(\frac{1}{\sqrt{x^2 + r^2}} - \frac{1}{\sqrt{x^2 + R^2}}\right)
  \end{align*}
  \hfill(1.0 points)

\item[(b)] In small displacement approximation:
  \[
    F_\perp \approx \frac{2GMmx}{R^2 - r^2}
      \left(\frac{1}{r} - \frac{1}{R}\right)
      = \left(\frac{2GMm}{Rr(R + r)}\right)x = kx
  \]
  \hfill(3.0 points)

  when the distance between the disk and the point mass is $x$, the distance
  of the center of mass of the system from the center of the disk will be:
  \[
    x_1 = \frac{mx}{m + M}
  \]
  \hfill(2.5 points)

  If we choose the center of mass of the system as an origin of our coordinate
  system, then equation of motion of the disk will have a form:
  \begin{align*}
    Ma_1 &= -kx \\
    Ma_1 + k\,\frac{m + M}{m}\,x_1 &= 0
  \end{align*}
  \hfill(2.5 points)

  For this simple harmonic oscillator, the period of small oscillations will
  be:
  \[
    T = 2\pi\sqrt{\frac{Rr(R + r)}{2G(M + m)}}
  \]
  \hfill(2.0 points)
\end{description}

\solhead{2.41}{DART}
% source id: 2023_TH_Q12
\paragraph{Part (a)}

Didymos' mass is much larger than Dimorphos' mass. Therefore, the radius of
the orbit of Dimorphos before impact, $a$, can be calculated from Kepler's 3rd
law:
\[
  \frac{GM}{4\pi^2} = \frac{a^3}{P^2},
\]
and therefore:
\[
  a = \left(\frac{GMP^2}{4\pi^2}\right)^{1/3} = 1.2\,\mathrm{km}.
\]
The orbital velocity of Dimorphos before impact was:
\[
  v_0 = \frac{2\pi a}{P} = 0.176\,\mathrm{m/s}.
\]
\hfill(2 points)

Let $v'$ be the Dimorphos velocity right after the collision. Using the law of
conservation of momentum, we have:
\[
  mv_0 - m_s v_s = (m + m_s)v',
\]
and:
\[
  v' = \frac{mv_0 - m_s v_s}{m + m_s} \approx v_0 - \frac{m_s}{m}v_s,
\]
where we used the fact that the mass of the spacecraft mass is much smaller % sic: "the mass of the spacecraft mass"
than the mass of Dimorphos.

We then use the vis-viva equation to calculate the semi-major axis of the
orbit after collision $a'$:
\[
  (v')^2 = GM\left(\frac{2}{a} - \frac{1}{a'}\right)
    = \frac{GM}{a}\left(2 - \frac{a}{a'}\right)
    = v_0^2\left(2 - \frac{a}{a'}\right),
\]
so
\[
  2 - \frac{a}{a'} = \left(\frac{v'}{v_0}\right)^2
    = \left(1 - \frac{m_s}{m}\frac{v_s}{v_0}\right)^2
    \approx 1 - \frac{2m_s}{m}\frac{v_s}{v_0}.
\]
Thus, the semi-major axis changed by:
\[
  \frac{\Delta a}{a} = \frac{a' - a}{a} = -\frac{2m_s}{m}\frac{v_s}{v_0}.
\]
\hfill(8 points)

If the semi-major axis changes from $a$ to $a + \Delta a$, then the orbital
period changes from $P$ to $P + \Delta P$, and the mass of the spacecraft can
be neglected. Then:
\[
  \frac{a^3}{P^2} = \frac{(a + \Delta a)^3}{(P + \Delta P)^2}
    = \frac{a^3(1 + \Delta a/a)^3}{P^2(1 + \Delta P/P)^2}
    = \frac{a^3}{P^2}
      \left(1 + \frac{3\Delta a}{a} - \frac{2\Delta P}{P}\right),
\]
hence:
\[
  \frac{\Delta P}{P} = \frac{3}{2}\frac{\Delta a}{a}.
\]
Thus, the orbital period changes by:
\[
  \frac{\Delta P}{P} = \frac{3}{2}\frac{\Delta a}{a}
    = -\frac{3m_s}{m}\frac{v_s}{v_0}
\]
\hfill(8 points)

We therefore expect that the orbital period of Dimorphos should decrease by
1.4\%, that is, 10 minutes. \hfill(2 points)

\medskip
\textit{Alternative solution (requires better numerical precision)}
\begin{align*}
  v_0 &= 0.17622\ \mathrm{m/s} \\
  v' &= 0.17662 - 0.00082 = 0.17540\ \mathrm{m/s} \\ % sic: 0.17662 for 0.17622
  a' &= 0.99080a = 1192.5\ \mathrm{m} \\
  P' &= 705.4\ \mathrm{min}
\end{align*}
Hence $\Delta P = 705.4 - 715.2 = -9.8 \approx -10$\,min. I propose to grant
full points ONLY if the first four significant figures match the solution. If
the first three significant figures match the solution, grant 80\% of the
points. Otherwise (if the method is correct), grant HALF of the points.

\paragraph{Part (b)}

Let $\Delta p$ be the momentum of the ejected debris. Then, the momentum
conservation equation becomes:
\[
  mv_0 - m_s v_s = (m + m_s)v' + \Delta p,
\]
\hfill(4 points)

so:
\[
  v' = \frac{mv_0 - m_s v_s - \Delta p}{m + m_s}
    \approx v_0 - \frac{m_s}{m}v_s - \frac{\Delta p}{m}.
\]
Using similar calculations as in point a), we get:
\[
  \frac{\Delta a}{a}
    = -2\left(\frac{m_s}{m}\frac{v_s}{v_0} + \frac{\Delta p}{mv_0}\right),
\]
\hfill(8 points)

hence:
\[
  \frac{\Delta P_0}{P} = \frac{3}{2}\frac{\Delta a}{a}
    = -3\left(\frac{m_s}{m}\frac{v_s}{v_0} + \frac{\Delta p}{mv_0}\right).
\]
Thus:
\[
  \frac{\Delta p}{mv_0} = -\frac{\Delta P_0}{3P}
    - \frac{m_s}{m}\frac{v_s}{v_0} = 0.011.
\]
\hfill(3 points)

\medskip
\textit{Alternative solution (requires better numerical precision)}

The orbital period after the collision is
$P' = 11.92 - 33/60 = 11.37$\,h. Thus, the semi-major axis of the orbit (after
the collision) is
\[
  a' = \left(\frac{GMP'^2}{4\pi^2}\right)^{1/3} = 1166\,\mathrm{m}.
\]
The velocity of Dimorphos right after the collision is
$v' = v_0\sqrt{2 - a/a'} = 0.1734\,\mathrm{m\,s^{-1}}$, so
$\Delta p/mv_0 = 0.011$. I propose to grant full points ONLY if the first four
significant figures match the solution. If the first three significant figures
match the solution, grant 80\% of the points. Otherwise (if the method is
correct), grant HALF of the points.

\paragraph{Part (c)}
\[
  \Delta v = v' - v_0
    = -\frac{m_s}{m}v_s - \frac{\Delta p}{m}
    = -\frac{m_s}{m}v_s + \frac{\Delta P_0}{3P}v_0 + \frac{m_s}{m}v_s
    = \frac{\Delta P_0}{3P}v_0 = -2.7\,\mathrm{mm/s}.
\]
\hfill(5 points)

\solhead{2.42}{Ground Tracks}
% source id: 2024_TH_T11
\subsubsection*{Part I: Sun-Synchronous Orbits}

\begin{description}
\item[(a)] Since the orbit is Sun-Synchronous, this means that the satellite's
  line of nodes must drift $360^\circ$ along one sidereal year. This ensures
  that the satellite will always pass over a given meridian at the same local
  mean solar time. Hence, \hfill(1.0 points)
  \[
    \dot\Omega = \frac{360^\circ}{T_s}
      = \frac{2\pi}{365.2564\cdot24\cdot60\cdot60}
      = 1.991\cdot10^{-7}\ \mathrm{rad/s}
  \]
  \hfill(2.0 points)

\item[(b)] After one orbit, the difference in longitude perceived by the
  satellite on the ground track is associated with the combined effects of the
  rotation of the Earth and nodal precession. This is expressed by
  \[
    \lambda_2 = \lambda_1 - \omega_\oplus\cdot T + \dot\Omega\cdot T
  \]
  \hfill(2.0 points)

  From the Ground Track, $\lambda_1 = 90^\circ$ and
  $\lambda_2 \approx -135^\circ$ for five orbits. Therefore, knowing the value
  of $\omega_\oplus$, the angular velocity of the rotation of our planet,
  \[
    T = \frac{1}{5}\cdot\frac{\lambda_1 - \lambda_2}{\omega_\oplus - \dot\Omega}
      = \frac{1}{5}\cdot
        \frac{225^\circ\cdot\pi/180^\circ}
             {7.292\cdot10^{-5} - 1.991\cdot10^{-7}}
      \approx 10800\ \mathrm{s} = 180\ \mathrm{min}
  \]
  \hfill(2.0 points)

  Although $\dot\Omega$ is significantly smaller than $\omega_\oplus$, it is
  conceptually important to include the nodal precession rate in the formula.

  The inclination of the orbit can be obtained directly from the Ground
  Track. Taking any of the two points of the orbit with the largest (absolute)
  value of latitude, one finds
  \[
    |\phi_{max}| \approx 55^\circ
  \]
  \hfill(2.0 points)

  Since $i > 90^\circ$, $|\phi_{max}|$ is the supplementary angle of the
  inclination:
  \[
    i = 180^\circ - 55^\circ = 125^\circ
  \]
  \hfill(2.0 points)

\item[(c)] Having calculated the period associated with the orbit, the
  semi-major axis is directly obtained from
  \[
    \frac{T^2}{a^3} = \frac{4\pi^2}{GM_\oplus}
    \implies
    a = \left(\frac{T^2 GM_\oplus}{4\pi^2}\right)^{1/3}
      = 1.06\cdot10^4\ \mathrm{km}
  \]
  \hfill(2.0 points)

\item[(d)] Since the orbit is sun-synchronous, the satellite completes an
  integer number of orbits during one solar day. Hence,
  \[
    N_S = \frac{24\cdot60}{180} = 8\ \text{orbits}
  \]
  \hfill(1.0 points)

\item[(e)] The satellite takes more than 2 orbits and less than 3 orbits to go
  from Poland to Brazil. Using as reference the point $P_o$ ahead of Macei\'o
  that is exactly 3 orbits after Chorz\'ow,
  \[
    \lambda_o = \lambda_C - 3\cdot45^\circ
  \]
  \[
    \phi_0 = \phi_C
  \]
  Using sine law in the following spherical triangle,

\ioaafig{2024_TH_T11-s1.pdf}

  \[
    \sin\theta = \frac{\sin\phi}{\sin(180^\circ - i)}
  \]
  \hfill(2.0 points)
  \[
    \theta_0 = \arcsin\left(\frac{\sin\phi_o}{\sin i}\right) = 69.9^\circ
  \]
  \[
    \theta_M = \arcsin\left(\frac{\sin\phi_M}{\sin i}\right) = -12.0^\circ
  \]
  The angle the satellite traveled will depend on how many complete orbits it
  took and the angle from the final whole orbit to Macei\'o:
  \[
    \Delta\theta = n\cdot360^\circ - (\theta_0 - \theta_M)
  \]
  \hfill(3.0 points)

  by the graph, $n = 3$, which results in $\Delta\theta = 998.1^\circ$.
  \hfill(1.0 points)

  Knowing its orbit is circular,
  \[
    \Delta t = T\cdot\frac{\Delta\theta}{360^\circ} = 8.32\ \mathrm{h}
  \]
  \hfill(1.0 points)

  The local time at Chorz\'ow over which the satellite passes by will be,
  considering $\lambda_M = 35.7^\circ W = -35.7^\circ$,
  \[
    T_C = T_M - \Delta t + \omega_\oplus^{-1}(\lambda_C - \lambda_M)
  \]
  \hfill(3.0 points)
  \[
    \boxed{T_C \approx 7\mathrm{h}20\mathrm{min}}
  \]
  \hfill(1.0 points)
\end{description}

\subsubsection*{Part II: Tundra orbits}

\begin{description}
\item[(f)] For the inclination, we need to take the maximum latitude reached
  by the satellite, which gives us $\boxed{i = 63^\circ}$.
  \hfill(1.0 points)

  For the period, we have to realize that in one orbit, the Earth's rotation
  does not cause a shift in the orbit, so it must be geosynchronous, that is
  $T = T_\oplus = \boxed{1436\ \mathrm{min}}$. \hfill(1.0 points)

  For the semi-major axis, we apply Kepler's Third Law:
  \[
    \frac{a^3}{T^2} = \frac{GM}{4\pi^2}
    \rightarrow \boxed{a = 42.164\ \mathrm{km}}
  \]
  \hfill(2.0 points)

\item[(g)] We will first find the expression for the area divided by the
  semi-latus rectum as a function of eccentricity so then we can apply the
  Kepler's Second Law to find the required time. \hfill(2.0 points)

  To do this, we will use the projection of the area of a circle when it is
  inclined by an angle $\arccos\left(\frac{b}{a}\right)$, as shown in the
  figure below:

\ioaafig{2024_TH_T11-s2.pdf}

  Thus, we will first calculate the hatched area of the circle figure
  $A_{circle}$
  \[
    A_{circle} = A_{sector} - A_{triangle}
      = \frac{\theta}{2\pi}\cdot\pi a^2 - \frac{ae\cdot b}{2}
  \]
  \hfill(2.0 points)

  We can also use that $\theta = \frac{\pi}{2} - \alpha$ and
  $\alpha = \sin^{-1}(e)$. Therefore:
  \[
    A_{circle} = \left(\frac{\pi\cdot a^2}{4}
      - \frac{\sin^{-1}(e)\cdot a^2}{2} - \frac{ab\cdot e}{2}\right)
  \]
  Now, we project this area,
  \[
    A_{ellipse} = A_{circle}\cdot\frac{b}{a}
  \]
  \[
    A_{ellipse} = \left(\frac{\pi ab}{4}
      - \frac{\sin^{-1}(e)\cdot ab}{2} - \frac{b^2 e}{2}\right)
  \]
  \hfill(4.0 points)

  Thus, the area covered by the satellite's position vector from one
  semi-latus rectum to the other will be given by twice the area found above,
  as we can see from the image, so that the area on the perigee side is given
  by
  \[
    A_{per} = \frac{\pi ab}{2} - b(a\sin^{-1}(e) + eb)
  \]
  For the apogee side, it is enough to know that
  $A_{ap} = \pi ab - A_{per}$, where $\pi ab$ is the area of the whole
  ellipse. Therefore, we have:
  \[
    A_{ap} = \frac{\pi ab}{2} + b\left(a\sin^{-1}(e) + eb\right)
  \]
  \hfill(2.0 points)

  Using Kepler's Second Law, we can write:
  \[
    \frac{A_{ap}}{A_{ellipse}} = \frac{T'}{T}
  \]
  \hfill(2.0 points)

  Using that $A_{ellipse} = \pi ab$ and $b = a\sqrt{1 - e^2}$, we have:
  \[
    \boxed{\frac{T'}{T} = \frac{1}{2} + \frac{\sin^{-1}(e)}{\pi}
      + \frac{e}{\pi}\cdot\sqrt{1 - e^2}}
  \]

  \textbf{Alternative Solution}

\ioaafig{2024_TH_T11-s3.pdf}

  The relation between the true anomaly and the mean anomaly is given by
  \[
    M = E - e\cdot\sin(E)
  \]
  From the figure (improve figure later), we observe that the satellite has % sic: "(improve figure later)" left in the official solutions
  the following properties:
  \begin{itemize}
    \item At perigee:
      \begin{itemize}
        \item $E = 0$
        \item $\theta = 0$
      \end{itemize}
    \item At the semi-latus rectum:
      \begin{itemize}
        \item $E = \cos^{-1}(e)$
        \item $\theta = 90^\circ$
      \end{itemize}
  \end{itemize}
  \hfill(2.0 points)

  Therefore, at the moment of passage through the semi-latus rectum, knowing
  that $\sin(E) = \sqrt{1 - e^2}$, we have
  \[
    M = \cos^{-1}(e) - e\cdot\sqrt{1 - e^2}
  \]
  \hfill(4.0 points)

  If we also use the perigee as the initial reference for the mean anomaly, we
  can associate the time of passage from the perigee to a certain point in the
  orbit by
  \[
    t = \frac{M}{\omega}
  \]
  where $\omega$ is the mean angular velocity of the orbit, given by
  $\omega = \frac{2\pi}{T}$, with $T$ being the orbital period. Thus, the time
  elapsed from the perigee to the semi-latus rectum is given by
  \[
    t = T\frac{\cos^{-1}(e) - e\cdot\sqrt{1 - e^2}}{2\pi}
  \]
  \hfill(2.0 points)

  Finally, the time $T'$ spent in the northern hemisphere can be found as
  follows:
  \[
    T' + 2t = T \rightarrow T' = T - 2t
  \]
  Thus, the desired ratio $T'/T$ will be:
  \[
    \frac{T'}{T} = 1 - \frac{2t}{T}
  \]
  \hfill(2.0 points)

  Substituting the expression for $t$, we get:
  \[
    \frac{T'}{T} = 1 - \frac{\cos^{-1}(e) + e\cdot\sqrt{1 - e^2}}{\pi}
  \]
  If we use the relation $\cos^{-1}(e) + \sin^{-1}(e) = \frac{\pi}{2}$, we
  finally obtain:
  \[
    \boxed{\frac{T'}{T} = \frac{1}{2} + \frac{\sin^{-1}(e)}{\pi}
      + \frac{e}{\pi}\cdot\sqrt{1 - e^2}}
  \]

\item[(h)] Between the two consecutive passages through the semi-latus rectum,
  we can use the following equation:
  \[
    \Delta\lambda = \pi - \omega_\oplus\cdot T'
  \]
  \hfill(2.0 points)

  Where $\Delta\lambda$ is the variation in longitude of the satellite between
  its passages through the semi-latus rectum, as shown in the image below.

\ioaafig{2024_TH_T11-s4.pdf}

  Therefore, we can write using the expression of the last item
  \[
    \Delta\lambda + \omega_\oplus\cdot T\cdot
      \left(\frac{1}{2} + \frac{\sin^{-1}(e)}{\pi}
      + \frac{e}{\pi}\cdot\sqrt{1 - e^2}\right) = \pi
  \]
  As $\omega_\oplus = \frac{2\pi}{T}$ (the orbit is geosynchronous), we will
  get:
  \begin{align*}
    \Delta\lambda + 2\cdot\sin^{-1}(e) + 2e\cdot\sqrt{1 - e^2} &= 0 \\
    \sin^{-1}(e) + e\cdot\sqrt{1 - e^2} &= -\frac{\Delta\lambda}{2}
  \end{align*}
  \hfill(2.0 points)

  Considering that $\Delta\lambda \approx -67.5^\circ$ from the graph (it is
  important to note that this value should be negative by the figure), we find
  from the above formula, by iteration, $\boxed{e \approx 0.3}$.
  \hfill(2.0 points)

  If we use the approximations $\sin(e) \approx e$ and $e^2 \ll 1$ given in
  the statement, we find:
  \[
    2e = -\frac{\Delta\lambda}{2}
  \]
  \hfill(3.0 points)

  which gives us an approximate answer $\boxed{e \approx 0.295 \approx 0.3}$
  \hfill(1.0 points)

\item[(i)] At the start and end points of retrograde motion, we have
  $\omega_{RA} = \omega_\oplus$, where $\omega_{RA}$ is the right ascension
  angular velocity of the satellite. \hfill(2.0 points)

  To find this velocity, we use the following spherical triangle:

\ioaafig{2024_TH_T11-s5.pdf}

  where $\theta_2 = \theta - 90^\circ$, and it is known that:
  \[
    \cos(\theta_2) = \cos(\delta)\cos(\alpha) \quad\text{(Cosine Law)}
  \]
  \hfill(1.0 points)
  \[
    \frac{\sin(\theta_2)}{\sin(\alpha)} = \frac{\cos(\delta)}{\cos(i)}
    \quad\text{(Sine Law)}
  \]
  \hfill(1.0 points)

  We substitute $\cos(\delta)$ from the first expression into the second
  expression, obtaining:
  \begin{align*}
    \frac{\sin(\theta_2)}{\sin(\alpha)}
      &= \frac{\cos(\theta_2)}{\cos(\alpha)\cdot\cos(i)} \\
    \tan(\theta_2)\cdot\cos(i) &= \tan(\alpha)
  \end{align*}
  We then derive the above expression with respect to time --- it is possible
  to skip the derivative step by just decomposing the satellite's angular
  velocity vector.
  \[
    \frac{\cos(i)}{\cos^2(\theta_2)}\,\dot\theta
      = \frac{\dot\alpha}{\cos^2(\alpha)}
  \]
  \hfill(2.0 points)

  Using the Cosine Law expression to replace $\cos(\alpha)$, we finally have:
  \[
    \omega_{RA} = \dot\alpha = \frac{\cos(i)}{\cos^2(\delta)}\,\dot\theta
  \]
  \hfill(2.0 points)

  We can then calculate the value of $\dot\theta$ as a function of the
  distance of the satellite from the center of the Earth, using the
  conservation of angular momentum:
  \[
    \dot\theta = \frac{2\pi}{T}\,\frac{a^2\sqrt{1 - e^2}}{r^2}
  \]
  \hfill(2.0 points)

  We still have the polar expression for the distance $r$ as a function of the
  polar radius, given by: % sic: "as a function of the polar radius"
  \[
    r = \frac{a(1 - e^2)}{1 + e\cdot\cos(\theta)}
  \]
  \hfill(2.0 points)

  Substituting everything into the expression that gives the condition for the
  start or end of retrograde motion, we find:
  \begin{align*}
    \frac{2\pi}{T}\,
      \frac{(1 + e\cos\theta)^2}{(1 - e^2)^{\frac{3}{2}}}\,
      \frac{\cos(i)}{\cos^2(\delta)}
      &= \omega_{earth} = \frac{2\pi}{T} \\
    \frac{(1 + e\cos\theta)^2}{(1 - e^2)^{\frac{3}{2}}}\,
      \frac{\cos(i)}{\cos^2(\delta)} &= 1
  \end{align*}
  \hfill(2.0 points)

  From the graph, we find the latitude --- or declination --- of the start or
  end of retrograde motion, given by $\delta \approx -32^\circ$.
  \hfill(2.0 points)

  Thus, substituting all the variables into the above equation and solving for
  $\theta$, we find: $\boxed{\theta = 54^\circ}$ and
  $\boxed{\theta = 306^\circ}$ as solutions. \hfill(2.0 points)

\item[(j)] For the figure-8 shape not to form, we need the inversion of motion
  in the northern hemisphere to occur exactly at the apogee position.
  Therefore, we will use exactly the same expression as the previous item, but
  with $\theta = 180^\circ$ and eccentricity as the unknown variable.
  \hfill(2.0 + 1.0 points)

  For the value of $\delta$, as the satellite is at the position of maximum
  declination, $\delta = i$. We have:
  \[
    \frac{(1 - e)^2}{(1 - e^2)^{\frac{3}{2}}}\,\frac{1}{\cos(i)} = 1
  \]
  \hfill(2.0 points)

  Solving the above equation, we obtain: $\boxed{e \approx 0.4}$
  \hfill(1.0 points)
\end{description}

\solhead{2.43}{Asteroid}
% source id: 2024_TH_T3
\begin{description}
\item[(a)] In this scenario, the asteroid would fall directly towards the
  star. However, in order to simplify the calculations, it is possible to
  consider that the asteroid would be in a degenerate elliptical orbit. In
  that case, the semi-minor axis would be infinitesimally small and the
  asteroid would still practically be moving on a straight line. It is also
  important to notice that the focii would virtually be at the periapsis and % sic: "focii", "apapsis"
  apapsis points. The degenerate elliptical orbit is shown in the following
  figure: \hfill(2.0 points)

\ioaafig{2024_TH_T3-s1.pdf}

  The semi-major axis of this elliptical orbit can be obtained through the
  following expression, where $m$ is the mass of the asteroid:
  \begin{align*}
    E = \frac{mv^2}{2} - \frac{GMm}{d} &= -\frac{GMm}{2a} \\
    \frac{GMm}{2d} - \frac{GMm}{d} &= -\frac{GMm}{2a} \\
    \frac{1}{2d} &= \frac{1}{2a} \\
    \therefore\ d &= a
  \end{align*}
  \hfill(2.0 points)

  Now, it is possible to use Kepler's third law to find an expression for the
  period of the orbit:
  \[
    \frac{T^2}{d^3} = \frac{4\pi^2}{GM}
      \longrightarrow T = 2\pi d\sqrt{\frac{d}{GM}}
  \]
  \hfill(1.0 points)

  If the asteroid is initially moving towards the star, it sweeps area I to
  reach the star. Using Kepler's second law, it is possible to calculate how
  long it takes for that to happen:
  \[
    \frac{\Delta t}{T} = \frac{A_I}{A_{Total}}
      = \frac{\frac{\pi ab}{4} - \frac{ab}{2}}{\pi ab}
  \]
  \hfill(2.0 points)
  \[
    \Delta t = \left(\frac{1}{4} - \frac{1}{2\pi}\right)T
  \]
  \[
    \boxed{\Delta t = \left(\frac{\pi}{2} - 1\right)d\sqrt{\frac{d}{GM}}}
  \]
  \hfill(1.0 points)

  Note that the radius of the star is negligible in these calculations since
  it is significantly smaller than $d$.

\item[(b)] If the asteroid is initially moving away from the star, it sweeps
  areas II and III before reaching the star. It is possible to again use
  Kepler's second law to find the time interval:
  \[
    \frac{\Delta t}{T} = \frac{A_{II} + A_{III}}{A_{Total}}
      = \frac{\frac{ab}{2} + \frac{\pi ab}{4} + \frac{\pi ab}{2}}{\pi ab}
  \]
  \hfill(1.5 points)
  \[
    \Delta t = \left(\frac{3}{4} + \frac{1}{2\pi}\right)T
  \]
  \[
    \boxed{\Delta t = \left(\frac{3\pi}{2} + 1\right)d\sqrt{\frac{d}{GM}}}
  \]
  \hfill(0.5 points)
\end{description}

\solhead{2.44}{Minor Planet}
% source id: 2020_DA_2
\begin{description}
\item[(a)] We consider the $\triangle ESP$, the vertices of which are Earth
  (E), Sun (S) and the minor planet (P).

\ioaafig{2020_DA_2-s1.pdf}

  where $\Delta\lambda = |\lambda - \lambda_\odot|$. \hfill(2.0 points)
  \begin{align*}
    r\ \text{(in m)} &= \frac{R_\oplus}{p\ \text{(in rad)}} \\
    \therefore\ r\ \text{(in au)} &= \frac{206265 R_\oplus}{p\cdot a_\oplus}
  \end{align*}
  \hfill(1.0 points)

  By using the cosine rule and sine rule,
  \begin{align*}
    r_h^2 &= a_\oplus^2 + r^2 - 2a_\oplus r\cos(\Delta\lambda) \\
    r_h &= \sqrt{a_\oplus^2 + r^2 - 2a_\oplus r\cos(\Delta\lambda)}
  \end{align*}
  \hfill(1.0 points)
  \begin{align*}
    \frac{\sin\theta_h}{r} &= \frac{\sin(\Delta\lambda)}{r_h} \\
    \therefore\ \theta_h &=
      \sin^{-1}\left(\frac{r\sin(\Delta\lambda)}{r_h}\right)
  \end{align*}
  \hfill(1.0 points)

  An angle $\psi_P$ is defined as angle swept by the vector
  $\overrightarrow{SP}$ from an arbitrary x axis (see figure). Let us fix the
  direction of the x-axis as the direction $\overrightarrow{SP}$ at the
  initial moment $t = 2012.3$. In other words, $\psi_P(2012.3) = 0$. Using the
  same x-axis, let us define corresponding angle for the Earth,
  $\psi_E$. Thus, we have
  \[
    \psi_P = \psi_E - \theta_h
  \]
  If we take any two consecutive measurements,
  \[
    \Delta\psi_E = \Delta\lambda_\odot
  \]
  \hfill(2.0 points)

  This is preferred over the conversion of time difference to angle as this
  data has better accuracy.

  Using these formulae, one can find $r$, $r_h$, $\psi_E$ and $\psi_P$ for
  every measurement instance in the data. $\psi_P$ and $r_h$ can be used to
  plot locations of the minor planet.

  Alternatively, one can also convert these to rectangular coordinates, i.e.
  \begin{align*}
    x &= r_h\cos\psi_P, \\
    y &= r_h\sin\psi_P;
  \end{align*}

  \begin{center}
  \begin{tabular}{rrrrrrrrr}
    \multicolumn{1}{c}{$t$} & \multicolumn{1}{c}{$r$} &
    \multicolumn{1}{c}{$r_h$} & \multicolumn{1}{c}{$\sin\theta_h$} &
    \multicolumn{1}{c}{$\theta_h$} & \multicolumn{1}{c}{$\psi_E$} &
    \multicolumn{1}{c}{$\psi_P$} & \multicolumn{1}{c}{$x$} &
    \multicolumn{1}{c}{$y$} \\
    \multicolumn{1}{c}{[year]} & \multicolumn{1}{c}{[au]} &
    \multicolumn{1}{c}{[au]} & &
    \multicolumn{1}{c}{[$^\circ$]} & \multicolumn{1}{c}{[$^\circ$]} &
    \multicolumn{1}{c}{[$^\circ$]} & \multicolumn{1}{c}{[au]} &
    \multicolumn{1}{c}{[au]} \\
    \hline
    2012.3 & 2.302 & 2.073 & $-1.0000$ & 270.00 & 270.00 & 0 & 2.073 & 0 \\
    2012.6 & 1.215 & 2.020 & $-0.4510$ & $-26.81$ & 3.88 & 30.69 & 1.737 & 1.031 \\
    2012.9 & 1.240 & 2.229 & 0.1097 & 6.30 & 111.13 & 104.83 & $-0.571$ & 2.155 \\
    2013.4 & 3.664 & 2.839 & 0.6391 & 140.27 & 293.89 & 153.62 & $-2.543$ & 1.262 \\
    2013.6 & 4.071 & 3.106 & $-0.2991$ & 197.40 & 3.64 & 166.24 & $-3.017$ & 0.739 \\
    2013.9 & 3.198 & 3.309 & $-0.9655$ & $-74.92$ & 110.87 & 185.79 & $-3.292$ & $-0.334$ \\
    2014.2 & 2.783 & 2.007 & 0.7357 & 132.63 & 222.34 & 89.71 & 0.010 & 2.007 \\
    2014.5 & 4.419 & 3.845 & 0.8735 & 119.13 & 328.09 & 208.96 & $-3.364$ & $-1.862$ \\
    2014.8 & 4.805 & 3.962 & $-0.5914$ & 216.26 & 74.50 & 218.24 & $-3.112$ & $-2.453$ \\
    2015.0 & 4.332 & 3.994 & $-0.9738$ & $-76.85$ & 149.24 & 226.09 & $-2.770$ & $-2.877$ \\
    2015.3 & 3.035 & 3.985 & 0.2736 & 15.88 & 257.6 & 241.72 & $-1.888$ & $-3.509$ \\
    \hline
  \end{tabular}
  \end{center}
  % In the original the 2014.2 row is printed in red (the outlier).

  there are 11 rows. Two points for each row. \hfill(22.0 points)

  \textbf{Note.} The parallax value for $p(2014.2) = 3.16''$ is obviously an
  outlier. Due to this, the corresponding position of the minor planet is also
  an outlier. Therefore, it should not be included in the plot or in the next
  analysis. \hfill(1.0 points)

\ioaafig{2020_DA_2-s2.pdf}

  These 4 points are for merely drawing the approximate polar or rectangular
  plot. Points for the line of apsides are separate (see below). Marking the
  earth positions or drawing the visually fit ellipse (blue line) is NOT
  expected.

  One may try to identify a minimum and a maximum and maximum in the % sic: "a maximum and maximum"
  calculated values of $r_h$ and take these as the aphelion and perihelion. In
  the figure these two points are shown and the line between drawn, i.e.\ the
  magenta line. Clearly these are not the aphelion and perihelion (as
  evidenced by the fact that the line connecting them does not pass through
  the focus), but they lie close. \hfill(1.0 + 1.0 points)

  We see that the points on either side of supposed aphelion are reasonably
  close (and the planet will be moving slower close to aphelion). However,
  there is a big gap between the second and third point, we take the aphelion
  position to be the correct one and draw the line of apsides to pass through
  aphelion and (0,0). (see below for the actual calculation of $a_p$ that does
  not rely on the graphical identification of the line of the apsides.)
  \hfill(2.0 points)

\item[(b)]
  \begin{description}
  \item[(i)] Assuming that the points above are close to the perihelion and
    aphelion, we should realise that the distance between them is roughly
    equal to the length of the major axis. Using the cosine rule,
    \begin{align*}
      \psi_1 &= 30.69^\circ \\
      r_1 &= 2.020\,\mathrm{au} \\
      \psi_2 &= 226.09^\circ \\
      r_2 &= 3.994\,\mathrm{au} \\
      \therefore\ \Delta\psi_P &= 226.09^\circ - 30.69^\circ \\
        &= 195.40^\circ \\
      2a_P &\approx \sqrt{r_1^2 + r_2^2 - 2r_1 r_2\cos\Delta\psi_P} \\
        &= \sqrt{2.020^2 + 3.994^2
           - 2\times2.02\times3.994\times\cos195.40^\circ} \\
        &= 5.965\,\mathrm{au} \\
      \therefore\ a_P &\approx 2.98\,\mathrm{au}
    \end{align*}
    \hfill(3.0 points)

  \item[(ii)]
    \[
      e = \frac{r_a - r_p}{r_a + r_p}
        \approx \frac{3.994 - 2.02}{3.994 + 2.02}
        \approx \frac{1}{3} = 0.33
    \]
    \hfill(2.0 points)

    Some students may visually fit an ellipse to the data and draw this line
    up to this visually fit ellipse. For the fit ellipse in the figure, the
    length of the major axis (green line) is $2a = 2.98$\,au and
    $e = \frac{1}{3}$, rotated by $52^\circ$. % sic: 2a = 2.98 au (cf. 2a_P = 5.965 au above)

  \item[(iii)] From Kepler's Laws,
    \begin{align*}
      P &= \sqrt{a_P^3} = \sqrt{2.98^3} \\
        &= 5.15\,\mathrm{yr}
    \end{align*}
    \hfill(1.0 points)

    The semi-period will be 2.58\,yr. Note that the time difference between
    our chosen points is only 2.4\,yr. This is fine, as we can see from the
    best fit ellipse that these points are not exactly along the line of
    apsides.

    At the same time, this is precisely the reason why one should not take
    $\Delta t$ between these two points as the semi-period.
  \end{description}

\item[(c)] The error in the semimajor axis is determined by the error in $r_1$
  and $r_2$ which, in turn, depend on $p$ and $\Delta\psi$. The latter is
  dominated by the uncertainty in the orientation of the ellipse (i.e., the
  assumption that $r_1$ and $r_2$ are close to the line of the apsides).

  In order to take this into account we take as $\Delta\psi_P$ the difference
  with the value that corresponds to the difference in angles of the actual
  perihelion and aphelion ($180^\circ$); i.e.,
  $\delta\Delta\psi_P \approx 15/180$, corresponding to
  $\Delta(\Delta\psi) \approx 0.08\pi \approx 0.26$. The geocentric distance
  is determined via parallax, the other quantities are constants, so it may be
  written $r_g \propto p^{-1}$. The immediate consequence is
  $\delta r_g = \delta p$ where $\delta$ is the designation for the relative
  error. Thus,
  \[
    \Delta r = r\delta r = r\delta p
  \]
  This gives $\Delta r_1 = 1.01\times10^{-2}$\,au and
  $\Delta r_2 = 1.997\times10^{-2}$\,au.

  The error in $a_p$ is then
  \begin{align*}
    \Delta a_P &= \frac{\sqrt{(r_1 - r_2\cos\Delta\Psi_p)^2\Delta r_1^2
      + (r_2 - r_1\cos\Delta\Psi_p)^2\Delta r_2^2
      + (r_1 r_2\sin\Delta\Psi_p)^2\Delta(\Delta\Psi_p)^2}}{2\sqrt{2}a_P} \\
      &= 0.04\,\mathrm{au} \\
    \therefore\ a_P &= (2.98\pm0.04)\,\mathrm{yr} % sic: units given as yr in source, should be au
  \end{align*}
  \hfill(3.0 points)

  The uncertainty in the period can then be obtained from Kepler's law and is
  \begin{align*}
    \delta P &= 3/2\,\delta a_p = 2\times10^{-2} \\
    \therefore\ P &= (5.2\pm0.1)\,\mathrm{yr}
  \end{align*}
  \hfill(1.0 points)

  The eccentricity can be calculated from $r_p = a_p(1 - e)$. Error
  propagation leads to
  \begin{align*}
    \Delta e &= \sqrt{\left((1 - e)\delta r\right)^2
      + \left(\frac{1 - e}{a_p}\Delta a_p\right)^2} \\
      &= 9\times10^{-3} \approx 0.01 \\
    \therefore\ e &= 0.33\pm0.01
  \end{align*}
  \hfill(2.0 points)
\end{description}
